% Rozdzial 16. Fragment do skryptu Nowoczesny C++ (bez preambuly).
% C++20. Wymagane pakiety obecne w poprzednich rozdzialach:
% listings, amsmath, amssymb, booktabs, tabularx, array, hyperref, enumitem.
% Wymagane makra: \code, \term, \nota (jak w pliku glownym skryptu).
% Dolaczenie: % Rozdzial 16. Fragment do skryptu Nowoczesny C++ (bez preambuly).
% C++20. Wymagane pakiety obecne w poprzednich rozdzialach:
% listings, amsmath, amssymb, booktabs, tabularx, array, hyperref, enumitem.
% Wymagane makra: \code, \term, \nota (jak w pliku glownym skryptu).
% Dolaczenie: % Rozdzial 16. Fragment do skryptu Nowoczesny C++ (bez preambuly).
% C++20. Wymagane pakiety obecne w poprzednich rozdzialach:
% listings, amsmath, amssymb, booktabs, tabularx, array, hyperref, enumitem.
% Wymagane makra: \code, \term, \nota (jak w pliku glownym skryptu).
% Dolaczenie: % Rozdzial 16. Fragment do skryptu Nowoczesny C++ (bez preambuly).
% C++20. Wymagane pakiety obecne w poprzednich rozdzialach:
% listings, amsmath, amssymb, booktabs, tabularx, array, hyperref, enumitem.
% Wymagane makra: \code, \term, \nota (jak w pliku glownym skryptu).
% Dolaczenie: \input{rozdzial_16_zaawansowana_wspolbieznosc_model_pamieci_cpp.tex}
% Kazdy listing opisany jako kompletny program jest samodzielny.

\chapter{Zaawansowana współbieżność, model pamięci i algorytmy bez blokad}
\label{chap:cpp16-advanced-concurrency}

\section{Cel rozdziału i powiązanie z podstawami}

W rozdziale~\ref{chap:concurrency} poznaliśmy wątki, mutexy, zmienne warunkowe, operacje
atomowe, \code{std::async}, pule wątków i~OpenMP. Sama znajomość tych
interfejsów nie wystarcza jednak do samodzielnego projektowania
algorytmów współbieżnych. Trzeba odpowiedzieć na trudniejsze pytania:
\begin{itemize}
\item Kiedy zapis wykonany przez jeden wątek \emph{musi} być widoczny dla drugiego?
\item Które wyniki są dozwolone przez model pamięci, nawet gdy wydają się
      sprzeczne z intuicją opartą na sekwencyjnym wykonywaniu?
\item Czym różni się atomowość pojedynczej operacji od niepodzielności
      całego algorytmu?
\item Jak zaprojektować \emph{postęp} algorytmu i~bezpiecznie zwalniać pamięć,
      gdy inne wątki mogą jeszcze posiadać wskaźniki do obiektów?
\item Kiedy rozwiązanie bez mutexa rzeczywiście poprawia wydajność?
\end{itemize}

Skupimy się na \emph{uzasadnianiu poprawności}, a~nie tylko na składni.
W~szczególności przedstawimy proste testy modelu pamięci (\emph{litmus tests}),
algorytm kolejki SPSC, publikację niezmiennych struktur danych, problem ABA,
techniki odzyskiwania pamięci oraz praktyczne zasady pomiaru i~diagnostyki.
Nie zakładamy, że każdy \code{std::atomic<T>} jest realizowany bez blokad.

\nota{Wszystkie gwarancje sformułowane w~rozdziale dotyczą abstrakcyjnego
modelu języka C++. Rzeczywisty procesor może wykonywać instrukcje
w~innej kolejności i~stosować bufory zapisów, ale program musi być
poprawny na gruncie standardu, a~nie tylko na jednym komputerze.}

\section{Formalny porządek zdarzeń w~wielowątkowym programie}
\label{sec:cpp16-formal}

\subsection{Zdarzenia i~relacje \emph{sequenced-before}, \emph{synchronizes-with}}

Rozpatrzmy pewne operacje na pamięci jako \emph{zdarzenia}. Wśród relacji
używanych do opisu ich uporządkowania najważniejsze są:
\begin{description}[style=nextline]
\item[\emph{Sequenced-before} (\(\mathrm{sb}\))]
  Porządek określony regułami języka pomiędzy wybranymi operacjami
  w~obrębie jednego wątku. Nie należy go mylić z~dowolną kolejnością
  pojawiania się instrukcji w~kodzie źródłowym: język zawiera także
  wyrażenia z~nieokreśloną kolejnością ewaluacji.
\item[\emph{Synchronizes-with} (\(\mathrm{sw}\))]
  Synchronizacja pomiędzy odpowiednimi operacjami różnych wątków:
  na przykład zapis \code{release} i~odczyt \code{acquire}, który czyta
  wartość z~odpowiedniego zapisu lub jego sekwencji \emph{release}.
  Podobną synchronizację zapewniają właściwe operacje na mutexach.
\item[\emph{Happens-before} (\(\mathrm{hb}\))]
  Relacja porządkująca zdarzenia, budowana w~typowych przykładach
  przez \(\mathrm{sb}\), \(\mathrm{sw}\) i~ich przechodnie złożenia.
  Jeśli \(A\;\mathrm{hb}\;B\), skutki \(A\) są uporządkowane przed \(B\)
  w~sensie formalnego modelu; nie jest to stwierdzenie o~czasie zegarowym.
\end{description}

Dla przejrzystości będziemy używać następującego schematu:
\[
\underbrace{A\;\xrightarrow{\mathrm{sb}}\;R}_{\text{producent}}
\quad\xrightarrow{\mathrm{sw}}\quad
\underbrace{Q\;\xrightarrow{\mathrm{sb}}\;B}_{\text{konsument}}
\quad\Longrightarrow\quad A\;\xrightarrow{\mathrm{hb}}\;B.
\]
Operacja \(R\) jest publikacją, \(Q\) odpowiadającym jej odbiorem;
\(A\) i~\(B\) mogą dotyczyć zwykłej, nieatomowej zmiennej.

\subsection{Atomowość nie zastępuje synchronizacji}

W~przeciwieństwie do zwykłego odczytu lub zapisu operacje na
\code{std::atomic<T>} nie powodują wyścigu danych pomiędzy sobą.
Nie oznacza to jednak, że każdy odczyt atomowy widzi ostatni zapis
według zegara ściennego. Taki globalny \emph{ostatni zapis} zwykle
w~ogóle nie istnieje jako zdefiniowane pojęcie modelu.

Z~kolei jeśli dwa wątki odczytują i~modyfikują to samo miejsce pamięci
nieatomowo, a~konfliktowe operacje nie są uporządkowane relacją
\emph{happens-before}, program ma \emph{data race} i~niezdefiniowane zachowanie.
Samo oznaczenie obiektu przez \code{volatile} nie rozwiązuje tego problemu.

\subsection{Porządek modyfikacji (\emph{modification order}) i~spójność}

Dla \emph{każdego pojedynczego} obiektu atomowego istnieje całkowity
porządek jego modyfikacji (\emph{modification order}). Nie wynika z~tego
jeden wspólny porządek modyfikacji \emph{wszystkich} atomików.
Ważne konsekwencje to między innymi:
\begin{itemize}
\item zapis do \code{x} uporządkowany przed innym zapisem do \code{x}
      musi być odpowiednio umieszczony przed nim w~porządku modyfikacji;
\item jeśli odczyt \code{x} jest uporządkowany przed kolejnym odczytem
      \code{x} w~tym samym wątku, drugi odczyt nie może cofnąć się
      do wcześniejszej modyfikacji w~porządku tego atomika;
\item ograniczenia spójności dla jednego atomika nie zapewniają
      automatycznie spójnego obrazu \emph{dwóch różnych} atomików.
\end{itemize}

\section{Testy modelu pamięci: wyniki, które zaskakują}
\label{sec:cpp16-litmus}

\subsection{\emph{Store buffering}: dwa odczyty zera}

Załóżmy, że \code{x} i~\code{y} początkowo są równe zero, a~wątki
wykonują następujące operacje:
\[
\begin{array}{c|c}
\text{Wątek 1} & \text{Wątek 2}\\ \hline
x\gets 1 & y\gets 1\\
r_1\gets y & r_2\gets x
\end{array}
\]

\begin{lstlisting}[caption={Store buffering: kompletna, wolna od wyścigu danych próba (C++20).}]
#include <atomic>
#include <iostream>
#include <thread>

int main() {
    std::atomic<int> x{0}, y{0};
    int r1 = -1, r2 = -1;

    std::jthread a([&] {
        x.store(1, std::memory_order_relaxed);
        r1 = y.load(std::memory_order_relaxed);
    });
    std::jthread b([&] {
        y.store(1, std::memory_order_relaxed);
        r2 = x.load(std::memory_order_relaxed);
    });
    a.join();
    b.join();
    std::cout << r1 << ' ' << r2 << '\n';
}
\end{lstlisting}

W~tym programie \(r_1=r_2=0\) jest \textbf{dozwolonym} wynikiem
modelu C++ dla \code{relaxed}; nie ma gwarancji, że zostanie
zaobserwowany w~praktycznym uruchomieniu. Jeśli \emph{wszystkie cztery}
operacje atomowe zamienimy na \code{memory_order_seq_cst}, wynik
\((0,0)\) stanie się niemożliwy: wspólny porządek operacji
sekwencyjnie spójnych musiałby przeczyć kolejności w~każdym z~wątków.

\nota{Wynik pojedynczego eksperymentu, nawet powtórzony milion razy,
nie dowodzi niedopuszczalności innego zachowania. Test tego rodzaju
to materiał do rozumowania nad modelem, a~nie pełny dowód poprawności.}

\subsection{\emph{Message passing}: publikacja zwykłych danych}

\begin{lstlisting}[caption={Publikacja danych przy użyciu release/acquire (C++11).}]
#include <atomic>
#include <cassert>
#include <thread>

int main() {
    int payload = 0;
    std::atomic<bool> published{false};

    std::thread producer([&] {
        payload = 42;
        published.store(true, std::memory_order_release);
    });
    std::thread consumer([&] {
        while (!published.load(std::memory_order_acquire)) {
            std::this_thread::yield();
        }
        assert(payload == 42);
    });

    producer.join();
    consumer.join();
}
\end{lstlisting}

Odczyt flagi, który widzi \code{true}, synchronizuje się z~zapisem
\code{release}. Zapis \code{payload = 42} następuje \emph{happens-before}
odczytu \code{payload} przez konsumenta. Flaga \emph{publikuje}
wcześniejszy stan; nie należy po publikacji modyfikować
\code{payload} bez dalszej synchronizacji.

\section{Porządki atomowe i~ich dokładne zastosowania}
\label{sec:cpp16-orders}

\begin{center}
\small
\begin{tabularx}{\textwidth}{@{}lX@{}}
\toprule
\textbf{Porządek} & \textbf{Typowe zastosowanie i~ograniczenie} \\
\midrule
\code{relaxed} & Liczniki, statystyki, pola kontrolne, gdy
sama atomowość wystarcza, a~relacja synchronizacji wynika skądinąd. \\
\code{release} & Zapis lub operacja RMW publikująca wcześniejsze efekty
wątku. \\
\code{acquire} & Odczyt lub operacja RMW odbierająca stan od publikującego wątku. \\
\code{acq_rel} & Operacja odczyt--modyfikacja--zapis wymagająca obu kierunków. \\
\code{seq_cst} & Operacje z~dodatkową gwarancją pojedynczego wspólnego
porządku dla operacji sekwencyjnie spójnych. \\
\bottomrule
\end{tabularx}
\end{center}

\code{memory_order_consume} ma długą i~skomplikowaną historię
standaryzacji. W~praktycznym kodzie C++20 należy preferować
\code{acquire} zamiast opierania przenośnej poprawności na
\code{consume}. Uwaga: porządek \code{release} jest właściwy dla
zapisu, \code{acquire} dla odczytu, a~\code{acq_rel} dla operacji
odczyt--modyfikacja--zapis; nie wszystkie porządki są poprawne dla
każdego przeciążenia funkcji atomowej.

\subsection{Operacje \emph{read--modify--write}}

Metody \code{fetch_add}, \code{fetch_sub}, \code{exchange} i~udany
\code{compare_exchange_*} wykonują operację RMW. Są one niepodzielne
w~odniesieniu do danego atomika, ale nie tworzą automatycznie
transakcji obejmującej inne pola.

\begin{lstlisting}[caption={Przydzielanie kolejnych numerów: wystarcza relaxed (C++11).}]
#include <atomic>
#include <cassert>
#include <thread>

int main() {
    std::atomic<int> next_id{0};
    int first = -1, second = -1;
    std::thread a([&] {
        first = next_id.fetch_add(1, std::memory_order_relaxed);
    });
    std::thread b([&] {
        second = next_id.fetch_add(1, std::memory_order_relaxed);
    });
    a.join();
    b.join();
    assert(first != second);
    assert(first + second == 1);
}
\end{lstlisting}

Uzyskujemy unikatowe wartości, bo RMW mają jednoznaczny porządek
modyfikacji \code{next_id}. Nie uzyskujemy przez to żadnej dodatkowej
synchronizacji niezwiązanych danych.

\subsection{\emph{Release sequence}: łańcuch operacji RMW}

W~C++20 zapis z~porządkiem \code{release} może rozpoczynać sekwencję
\emph{release} składającą się z~następujących po nim operacji RMW
na \emph{tym samym} atomiku, wykonywanych także przez inne wątki.
Odczyt \code{acquire} wartości pochodzącej z~tej sekwencji może
synchronizować się z~początkowym zapisem \code{release}.

\begin{lstlisting}[caption={Przekazanie publikacji przez relaxed RMW (C++20).}]
#include <atomic>
#include <cassert>
#include <thread>

int main() {
    int payload = 7;
    std::atomic<int> state{0};

    std::thread producer([&] {
        payload = 42;
        state.store(1, std::memory_order_release);
    });
    std::thread intermediary([&] {
        int expected = 0;
        while ((expected = state.load(std::memory_order_acquire))
               == 0) {
            std::this_thread::yield();
        }
        state.fetch_add(1, std::memory_order_relaxed);
    });
    std::thread consumer([&] {
        while (state.load(std::memory_order_acquire) != 2) {
            std::this_thread::yield();
        }
        assert(payload == 42);
    });

    producer.join();
    intermediary.join();
    consumer.join();
}
\end{lstlisting}

Wątek pośredni czeka, aż stan osiągnie 1, i~wykonuje RMW, który
zapisuje 2. Odczyt wartości 2 z~\code{acquire} synchronizuje się
z~pierwotnym \code{release} przez sekwencję \emph{release}.
Nie można w~ogólności zastąpić \code{fetch_add} zwykłym
\code{store(2, relaxed)} i~utrzymać tej samej argumentacji.

\section{CAS: porównaj i~zamień, pętle ponawiania i~porządek niepowodzenia}

\begin{center}\small
\begin{tabularx}{\textwidth}{@{}p{0.38\textwidth}X@{}}
\toprule
Operacja atomowa & Działanie \\
\midrule
\code{compare_exchange_strong} & Jednorazowo porównuje i próbuje podmienić wartość; nie zgłasza pozornego niepowodzenia. \\
\code{compare_exchange_weak} & Może zgłosić pozorne niepowodzenie; nadaje się do pętli ponawiania. \\
\code{expected} przy CAS & Przy niepowodzeniu otrzymuje rzeczywistą wartość atomika, którą trzeba uwzględnić przy kolejnej próbie. \\
\code{memory_order_acq_rel} & Przy udanej operacji odczytująco-zapisującej łączy własności acquire i release. \\
\code{memory_order_relaxed} & Zapewnia atomowość bez dodatkowej synchronizacji z innymi danymi. \\
\bottomrule
\end{tabularx}
\end{center}
\label{sec:cpp16-cas}

\code{compare_exchange_weak(expected, desired, success, failure)}
atomowo porównuje bieżącą wartość z~\code{expected}. Gdy porównanie
się powiedzie, zapisuje \code{desired}; w~przeciwnym razie aktualizuje
\code{expected} do wartości zastanej w~atomiku. Wersja \code{weak}
może ponadto nie powieść się pozornie (\emph{spurious failure}).
W~pętli jest zwykle dobrym wyborem.

\begin{lstlisting}[caption={Atomowe maksimum z~pętlą CAS (C++11).}]
#include <atomic>
#include <cassert>

void atomic_max(std::atomic<int>& target, int candidate) {
    int old = target.load(std::memory_order_relaxed);
    while (old < candidate &&
           !target.compare_exchange_weak(
               old, candidate,
               std::memory_order_relaxed,
               std::memory_order_relaxed)) {
        // Przy niepowodzeniu old jest aktualizowane.
    }
}

int main() {
    std::atomic<int> maximum{10};
    atomic_max(maximum, 8);
    atomic_max(maximum, 17);
    assert(maximum.load() == 17);
}
\end{lstlisting}

Parametr \code{failure} opisuje porządek operacji \emph{odczytu}
w~przypadku niepowodzenia. Nie może to być \code{release} ani
\code{acq_rel}; musi też spełniać ograniczenia względem porządku
sukcesu określone standardem. Typowy wzorzec publikacji nowej
wartości to sukces z~\code{acq_rel} i~niepowodzenie z~\code{acquire}
(albo \code{relaxed}, jeśli tak wynika z~dowodu).

\subsection{\code{weak} kontra \code{strong}}

\code{compare_exchange_strong} nie zgłasza pozornego niepowodzenia
przy poprawnie porównanych wartościach. Może być użyteczne dla
pojedynczej próby; \code{weak} bywa wygodniejsze w~pętlach CAS.
Nie ma uniwersalnej gwarancji, że jedna z~tych wersji zawsze działa
szybciej na danej platformie.

\section{Bariery pamięci i~\texttt{std::atomic\_thread\_fence}}
\label{sec:cpp16-fences}

Bariera \code{std::atomic_thread_fence(order)} jest operacją
porządkowania pamięci, która sama nie odczytuje ani nie zapisuje
konkretnego atomika. Synchronizacja z~innym wątkiem wciąż wymaga
odpowiednich operacji atomowych łączących obie strony. Samo
wstawienie \code{fence} obok zwykłej zmiennej nie usuwa wyścigu danych.

\begin{lstlisting}[caption={Publikacja ze zwolnieniową i~przejmującą barierą pamięci (C++11).}]
#include <atomic>
#include <cassert>
#include <thread>

int main() {
    int payload = 0;
    std::atomic<bool> flag{false};

    std::thread producer([&] {
        payload = 42;
        std::atomic_thread_fence(std::memory_order_release);
        flag.store(true, std::memory_order_relaxed);
    });
    std::thread consumer([&] {
        while (!flag.load(std::memory_order_relaxed)) {
            std::this_thread::yield();
        }
        std::atomic_thread_fence(std::memory_order_acquire);
        assert(payload == 42);
    });

    producer.join();
    consumer.join();
}
\end{lstlisting}

To przykład synchronizacji \emph{fence--fence}: odczyt atomika
widzi zapis następujący po barierze \code{release}, a~bariera
\code{acquire} jest wykonywana po tym odczycie. W~nowym kodzie
zwykle prostsze i~czytelniejsze jest zastosowanie pary operacji
\code{release}/\code{acquire} bez jawnych barier. Bariery nie są
substytutem przemyślanego protokołu komunikacji.

\section{Gwarancje postępu: \emph{blocking}, \emph{lock-free} i~\emph{wait-free}}
\label{sec:cpp16-progress}

Wielowątkowy algorytm można analizować pod względem \emph{bezpieczeństwa}
(\emph{safety}: nic niedozwolonego się nie dzieje) i~\emph{żywotności}
(\emph{liveness}: coś pożądanego w~końcu się wydarzy). Są to różne kwestie.

\begin{description}[style=nextline]
\item[\emph{Blocking}]
  Wątek może oczekiwać na odblokowanie mutexa przez konkretny inny wątek.
  Jeśli właściciel przestanie robić postępy, pozostali również mogą utknąć.
\item[\emph{Lock-free}]
  W~standardowym ujęciu postępu algorytmowego, mimo rywalizacji wątków,
  system jako całość może kontynuować wykonywanie operacji:
  niepowodzenia jednej próby CAS zwykle oznaczają, że inny uczestnik
  wykonał postęp. Nie gwarantuje to postępu \emph{każdego} wątku.
\item[\emph{Wait-free}]
  Każda operacja danego wątku ma ograniczoną liczbę własnych kroków
  niezależnie od działań innych wątków (przy przyjętym modelu).
  To silniejszy warunek niż \emph{lock-free}.
\item[\emph{Obstruction-free}]
  Operacja postępuje, jeśli wątek może działać dostatecznie długo
  bez konfliktów ze strony innych wątków. Jest to słabsza gwarancja.
\end{description}

\emph{Lock-free} nie znaczy ,,zawsze szybciej''. Nie oznacza również
braku aktywnego oczekiwania ani braku zagłodzenia konkretnego wątku.
Odróżniajmy też własność implementacji pojedynczego atomika
(\code{is_lock_free()}) od postępu kompletnego algorytmu.

\subsection{Prosta blokada wirująca (\emph{spinlock})}

\begin{lstlisting}[caption={Poprawna blokada wirująca z flagą atomową; ilustracja (C++20).}]
#include <atomic>
#include <cassert>
#include <mutex>
#include <thread>

class SpinLock {
    std::atomic_flag flag = ATOMIC_FLAG_INIT;
public:
    void lock() noexcept {
        while (flag.test_and_set(std::memory_order_acquire)) {
            // Zamiast stale rywalizowac o cache line, czekamy.
            flag.wait(true, std::memory_order_relaxed);
        }
    }
    void unlock() noexcept {
        flag.clear(std::memory_order_release);
        flag.notify_one();
    }
};

int main() {
    SpinLock lock;
    int counter = 0;
    auto work = [&] {
        for (int i = 0; i < 10000; ++i) {
            std::lock_guard<SpinLock> guard(lock);
            ++counter;
        }
    };
    std::jthread a(work), b(work);
    a.join();
    b.join();
    assert(counter == 20000);
}
\end{lstlisting}

\code{test_and_set} ma semantykę RMW; odblokowanie jest operacją
\code{release}. Metoda \code{wait(true)} może blokować, a~zatem
powyższy \code{SpinLock} \emph{nie} jest algorytmem \emph{lock-free}.
Ponadto nie gwarantuje sprawiedliwości. W~realnym oprogramowaniu
\code{std::mutex} zwykle pozostaje lepszym wyborem, zwłaszcza gdy
sekcja krytyczna może być długa lub obciążenie jest wysokie.

\section{Problem ABA: poprawne CAS nie wystarcza}
\label{sec:cpp16-aba}

Rozważmy wierzchołek stosu reprezentowany wskaźnikiem atomowym.
Wątek \(A\) czyta wskaźnik \(p\) i~zamierza wykonać
\code{compare_exchange} z~wartości \(p\) na \code{p->next}.
W~tym czasie wątek \(B\) może:
\begin{enumerate}
\item zdjąć element \(p\) ze stosu;
\item zdjąć jeszcze jeden element;
\item ponownie umieścić \(p\) na wierzchołku stosu.
\end{enumerate}

Wartość wskaźnika uległa zmianie \(A\to B\to A\) (stąd \emph{ABA}),
ale CAS wątku \(A\) może zobaczyć ponownie tę samą wartość \(p\).
W~międzyczasie struktura stosu mogła się całkowicie zmienić.

\begin{center}
\(\boxed{\text{head}=p}\;\longrightarrow\;
  \boxed{\text{head}=q}\;\longrightarrow\;
  \boxed{\text{head}=p}\)
\end{center}

Dwa problemy trzeba odróżnić:
\begin{enumerate}
\item \emph{ABA logiczne}: porównanie tej samej wartości nie wykrywa
      zmian stanu struktury.
\item \emph{Bezpieczeństwo pamięci}: wątek \(A\) może próbować odczytać
      \code{p->next}, gdy inny wątek zniszczył już obiekt \code{*p}.
      Taki odczyt nie staje się poprawny tylko dlatego, że \code{head}
      jest atomowy.
\end{enumerate}

\subsection{Wskaźnik z~licznikiem generacji}

Jedna z~technik polega na porównywaniu pary
\((\text{wskaźnik},\text{licznik wersji})\) zamiast samego wskaźnika.
Wtedy ponowne pojawienie się tego samego adresu nie oznacza równości
całego stanu. Licznik również ma ograniczony zakres i~może się
przepełnić; trzeba uzasadnić brak niebezpiecznego \emph{wrap-around}.

\nota{Nie zakładaj, że \code{std::atomic<SomeStruct>} będzie bez blokad:
implementacja może wykorzystać wewnętrzny mutex lub bibliotekę
atomowych operacji. Możliwość zapisania stanu złożonego w~\code{atomic}
nie rozwiązuje też samodzielnie problemu czasu życia obiektów.}

\subsection{Hazard pointers, epoch-based reclamation i~RCU}

\term{Hazard pointers} pozwalają wątkom publikować informację:
,,zamierzam korzystać z~tego obiektu''. Wątek usuwający obiekt
odkłada go na listę \emph{retired} i~odracza zniszczenie do chwili,
gdy żaden aktywny hazard pointer go nie wskazuje. Poprawna procedura
zwykle musi ponownie zweryfikować odczyt wskaźnika po publikacji
hazardu; samo opublikowanie dowolnego, potencjalnie już nieważnego
adresu nie wystarcza.

\term{Epoch-based reclamation} grupuje operacje w~epoki. Usunięte
obiekty mogą zostać odzyskane dopiero wtedy, gdy wszystkie aktywne
wątki zakończyły okresy korzystania z~danych obejmujące ich usunięcie.
Dłużej zatrzymany wątek może opóźniać zwolnienie pamięci.

\term{Read-copy-update (RCU)} korzysta z~odczytów starych wersji
danych i~publikacji nowych wersji; usuwanie poprzednich odbywa się
po odpowiednim okresie ochronnym (\emph{grace period}). Szczegóły
różnią się pomiędzy implementacjami.

\begin{center}
\small
\begin{tabularx}{\textwidth}{@{}lX@{}}
\toprule
\textbf{Metoda} & \textbf{Kiedy rozważyć} \\
\midrule
Hazard pointers & Wiele odczytów wskaźników, częste usuwanie
węzłów; możliwa kosztowna administracja hazardami. \\
Epoki & Duży ruch, kontrolowane sekcje dostępu; zatrzymany czytelnik
może opóźnić odzyskanie pamięci. \\
RCU & Wzorce ,,wielu czytelników, niewielu pisarzy'', publikacja
niezmiennych wersji danych. \\
\code{shared_ptr} & Gdy łatwość dowodu poprawności jest ważniejsza
od kosztu liczenia referencji; nie obiecuje \emph{lock-free}. \\
\bottomrule
\end{tabularx}
\end{center}

W~C++20 standardowa biblioteka nie udostępnia jeszcze kompletnej,
przenośnej abstrakcji hazard pointerów. Bezpieczna implementacja
wymaga precyzyjnego protokołu pamięciowego, dlatego zamiast podawać
uproszczony, ale błędny stos ze zwalnianiem węzłów, przedstawimy
algorytm bez dynamicznego odzyskiwania pamięci: kolejkę SPSC.

\section{Kolejka SPSC: praktyczny algorytm bez mutexa}
\label{sec:cpp16-spsc}

\emph{Single Producer, Single Consumer} (SPSC) oznacza dokładnie
jednego producenta i~jednego konsumenta. Bufor cykliczny może wtedy
używać dwóch liczników:
\begin{itemize}
\item \code{head} --- numer kolejnego miejsca zapisu, modyfikowany
      wyłącznie przez producenta;
\item \code{tail} --- numer kolejnego miejsca odczytu, modyfikowany
      wyłącznie przez konsumenta.
\end{itemize}

Producent zapisuje komórkę, a~następnie publikuje nowy \code{head}
operacją \code{release}. Konsument odczytuje \code{head} przez
\code{acquire}, a~więc może bezpiecznie zobaczyć zawartość wcześniej
zapisanej komórki. Analogicznie konsument zwalnia komórkę, publikując
nowy \code{tail}, a~producent odbiera tę informację operacją
\code{acquire}. Warunek pełności bufora o~pojemności \(C\) to
\(\mathrm{head}-\mathrm{tail}=C\), a~pustości:
\(\mathrm{head}=\mathrm{tail}\).

\begin{lstlisting}[caption={Kompletna kolejka SPSC dla int, ze statyczną pamięcią (C++20).}]
#include <array>
#include <atomic>
#include <cassert>
#include <cstddef>
#include <thread>

// Tylko jeden producent i dokladnie jeden konsument.
template<std::size_t Capacity>
class SpscQueue {
    static_assert(Capacity > 0);
    static_assert((Capacity & (Capacity - 1)) == 0);
    std::array<int, Capacity> buffer_{};
    std::atomic<std::size_t> head_{0};
    std::atomic<std::size_t> tail_{0};
public:
    bool try_push(int value) {
        const auto head = head_.load(std::memory_order_relaxed);
        const auto tail = tail_.load(std::memory_order_acquire);
        if (head - tail == Capacity) return false;
        buffer_[head & (Capacity - 1)] = value;
        head_.store(head + 1, std::memory_order_release);
        return true;
    }

    bool try_pop(int& result) {
        const auto tail = tail_.load(std::memory_order_relaxed);
        const auto head = head_.load(std::memory_order_acquire);
        if (tail == head) return false;
        result = buffer_[tail & (Capacity - 1)];
        tail_.store(tail + 1, std::memory_order_release);
        return true;
    }
};

int main() {
    constexpr int Count = 100000;
    SpscQueue<1024> queue;
    long long result = 0;

    std::jthread producer([&] {
        for (int x = 1; x <= Count; ++x) {
            while (!queue.try_push(x)) std::this_thread::yield();
        }
    });
    std::jthread consumer([&] {
        int item = 0;
        for (int n = 0; n < Count; ++n) {
            while (!queue.try_pop(item)) std::this_thread::yield();
            result += item;
        }
    });
    producer.join();
    consumer.join();
    assert(result == 1LL * Count * (Count + 1) / 2);
}
\end{lstlisting}

\subsection{Dlaczego kod jest poprawny?}

Dla danych w~\code{buffer_} nie stosujemy operacji atomowych.
Nie ma jednak konfliktowego dostępu bez synchronizacji, ponieważ:
\begin{enumerate}
\item przed odczytem komórki konsument otrzymuje informację o~zapisie
      przez \code{head_.load(acquire)};
\item przed ponownym nadpisaniem tej komórki producent otrzymuje
      informację o~jej zwolnieniu przez \code{tail_.load(acquire)};
\item każdy z~liczników ma jednego piszącego, więc nie potrzeba CAS
      do aktualizowania \code{head_} ani \code{tail_}.
\end{enumerate}

Przykład zakłada, że operacji \code{try_push} nie wywołuje kilku
producentów, a~\code{try_pop} kilku konsumentów, nawet jeśli
same metody są dostępne publicznie. Złamane założenie SPSC
powoduje utratę poprawności. Ograniczenia bufora, przepełnienie
liczników i~kwestie typów niebanalnych wymagają osobnej analizy.
Użyliśmy potęgi dwójki, żeby indeksowanie przez maskę pozostało
poprawne także przy modularnym zawijaniu licznika bez znaku.

\nota{Ten przykład jest edukacyjną, \emph{nieblokującą} konstrukcją
w~sensie, że nie przejmuje mutexa, ale aktywne pętle producenta
i~konsumenta mogą czekać dowolnie długo. Metody \code{try_push}
i~\code{try_pop} są krótkie i~nie wykonują pętli CAS. Bez uzasadnienia
własności użytych atomików na danej platformie nie deklarujemy
silniejszej gwarancji wydajnościowej.}

\section{Atomowa publikacja niezmiennych wersji danych}
\label{sec:cpp16-snapshot}

Wiele rzeczywistych aplikacji (konfiguracja serwera, słowniki,
parametry modeli) posiada licznych czytelników i~rzadkich pisarzy.
Dobrą architekturą bywa utworzenie nowej, niezmiennej wersji
struktury i~jednorazowe opublikowanie wskaźnika do niej.
W~C++20 można do tego wykorzystać specjalizację
\code{std::atomic<std::shared_ptr<T>>}.

\begin{lstlisting}[caption={Bezpieczny snapshot konfiguracji (C++20).}]
#include <atomic>
#include <cassert>
#include <memory>
#include <thread>

struct Config {
    int threshold;
    double scale;
};

int main() {
    using Ptr = std::shared_ptr<const Config>;
    std::atomic<Ptr> current{
        std::make_shared<const Config>(Config{10, 1.0})
    };

    std::jthread writer([&] {
        Ptr next = std::make_shared<const Config>(Config{20, 2.0});
        current.store(std::move(next), std::memory_order_release);
    });
    std::jthread reader([&] {
        Ptr snap = current.load(std::memory_order_acquire);
        assert((snap->threshold == 10 && snap->scale == 1.0) ||
               (snap->threshold == 20 && snap->scale == 2.0));
    });
    writer.join();
    reader.join();
}
\end{lstlisting}

Czytelnik otrzymuje \emph{spójny} snapshot obu pól, a~stara wersja
jest automatycznie utrzymywana przy życiu, dopóki istnieją do niej
uchwyty \code{shared_ptr}. Nie oznacza to, że przechowywanie
\code{shared_ptr} w~atomiku jest \emph{lock-free}. Referencyjne
liczenie i~alokacje bywają istotnym kosztem przy bardzo częstych
odczytach. Dla obciążenia z~dużą przewagą czytelników należy
porównać to podejście z~\code{shared_mutex}, kopiowaniem małych
struktur i~technikami RCU.

\subsection{\code{std::atomic_ref}: synchronizacja istniejącego obiektu}

\code{std::atomic_ref<T>} (C++20) pozwala wykonywać operacje atomowe
na istniejącym, odpowiednio wyrównanym obiekcie typu \code{T}.
Przydatne bywa przy interoperacyjności ze strukturami danych,
które nie przechowują \code{std::atomic<T>} jako pola.

\begin{lstlisting}[caption={Operacje atomowe na istniejącej zmiennej (C++20).}]
#include <atomic>
#include <cassert>
#include <thread>

int main() {
    alignas(std::atomic_ref<int>::required_alignment) int counter = 0;
    {
        std::atomic_ref<int> shared_counter(counter);
        auto add = [&] {
            for (int i = 0; i < 10000; ++i)
                shared_counter.fetch_add(1, std::memory_order_relaxed);
        };
        std::jthread a(add), b(add);
        a.join();
        b.join();
    } // koniec zycia atomic_ref
    assert(counter == 20000);
}
\end{lstlisting}

Dopóki istnieje \code{atomic_ref} odnoszący się do obiektu, dostęp
do tego obiektu należy realizować przez instancje \code{atomic_ref}
--- nie tylko w~chwilach równoczesnego dostępu. Dlatego zwykłe
odczytanie \code{counter} wykonujemy \emph{po} zakończeniu czasu
życia referencji atomowej. Sam obiekt musi żyć dostatecznie długo,
a~jego wyrównanie musi spełniać wymagania \code{required_alignment}.

\section{Planowanie zadań: pule wątków, \emph{work stealing} i~\emph{backpressure}}
\label{sec:cpp16-scheduling}

Poza poprawną synchronizacją potrzebujemy rozsądnej strategii
przypisywania prac do wątków. Wątek systemowy ma koszt utworzenia,
przełączenia kontekstu i~utrzymania stosu. Z~tego powodu tysiące
krótkich zadań nie powinny automatycznie oznaczać tysiąca wątków.

\begin{description}[style=nextline]
\item[\emph{Thread pool}]
  Ustalona liczba wątków pobiera zadania z~kolejki. Ogranicza koszt
  częstego tworzenia wątków i~może ograniczać równoległe obciążenie.
\item[\emph{Work stealing}]
  Każdy wątek ma własną kolejkę, a~bezczynny pracownik próbuje
  przejąć część pracy od innych. Dobrze sprawdza się przy nierównych
  czasach trwania podzadań. Wymaga ostrożnego protokołu synchronizacji
  kolejek, szczególnie gdy korzystają z~nich różni producenci.
\item[\emph{Backpressure}]
  Mechanizm spowalniania producenta, gdy konsumenci nie nadążają.
  Bez niego kolejka bez ograniczenia pojemności może zużyć całą pamięć.
\item[\emph{Structured concurrency}]
  Sposób organizacji zadań, w~którym ich czas życia, anulowanie i~błędy
  pozostają związane z~zakresem nadrzędnej operacji. Sam \code{std::async}
  nie implementuje kompletnego frameworka takiego stylu.
\end{description}

\subsection{Pula wątków: bezpieczne zamknięcie i~propagacja wyjątków}

Następny przykład rozwija pulę z~rozdziału~\ref{chap:concurrency}. Używa mutexa
oraz zmiennej warunkowej. Łatwy dowód poprawności i~przewidywalne
zachowanie przy zamykaniu są często ważniejsze niż
niepotwierdzony \emph{lock-free}.

\begin{lstlisting}[caption={Prosta pula wątków odprowadzająca kolejkę przed zamknięciem (C++20).}]
#include <condition_variable>
#include <functional>
#include <future>
#include <iostream>
#include <memory>
#include <mutex>
#include <queue>
#include <stdexcept>
#include <thread>
#include <type_traits>
#include <utility>
#include <vector>

class ThreadPool {
    std::mutex mutex_;
    std::condition_variable cv_;
    std::queue<std::function<void()>> jobs_;
    bool closed_ = false;
    std::vector<std::jthread> workers_;

    void worker_loop() {
        for (;;) {
            std::function<void()> job;
            {
                std::unique_lock lock(mutex_);
                cv_.wait(lock, [&] {
                    return closed_ || !jobs_.empty();
                });
                if (closed_ && jobs_.empty()) return;
                job = std::move(jobs_.front());
                jobs_.pop();
            }
            job();  // packaged_task przechwyci wyjatek
        }
    }

public:
    explicit ThreadPool(unsigned count) {
        if (count == 0) throw std::invalid_argument("count=0");
        try {
            for (unsigned i = 0; i < count; ++i)
                workers_.emplace_back([this] { worker_loop(); });
        } catch (...) {
            close();  // obudz pracownikow; ich destruktory wykonaja join
            throw;
        }
    }
    ThreadPool(const ThreadPool&) = delete;
    ThreadPool& operator=(const ThreadPool&) = delete;
    ~ThreadPool() { close(); }

    void close() {
        {
            std::lock_guard lock(mutex_);
            closed_ = true;
        }
        cv_.notify_all();
    }

    template<class F>
    auto submit(F&& callable)
        -> std::future<std::invoke_result_t<std::decay_t<F>&>> {
        using R = std::invoke_result_t<std::decay_t<F>&>;
        auto task = std::make_shared<std::packaged_task<R()>>(
            std::forward<F>(callable)
        );
        auto result = task->get_future();
        {
            std::lock_guard lock(mutex_);
            if (closed_) throw std::runtime_error("pool closed");
            jobs_.emplace([task] { (*task)(); });
        }
        cv_.notify_one();
        return result;
    }
};

int main() {
    ThreadPool pool(2);
    auto a = pool.submit([] { return 40; });
    auto b = pool.submit([] { return 2; });
    std::cout << a.get() + b.get() << '\n'; // 42
    pool.close(); // od tej chwili bez nowych zadan
}
\end{lstlisting}

Ważne własności pokazanej implementacji:
\begin{itemize}
\item \code{close()} jest bezpieczne wobec równoległego \code{submit()}:
      oba działania są serializowane przez ten sam mutex;
\item po \code{close()} pracownicy \emph{kończą kolejkę}, a~dopiero
      następnie opuszczają pętlę; destruktory \code{jthread} dołączają
      pracowników, zanim zniszczone zostaną mutex i~kolejka;
\item \code{packaged_task} zapisuje wartość lub wyjątek w~stanie
      asynchronicznym, więc zwykłe wyjątki z~zadań nie wychodzą
      bezpośrednio przez \code{worker_loop};
\item ta minimalna pula nie ma ograniczenia długości kolejki,
      priorytetów, terminów ani anulowania uruchomionych prac.
\end{itemize}

\nota{Nie wolno niszczyć takiej puli z~zadania wykonywanego w~tej
samej puli: dołączenie własnego wątku prowadzi do błędu.
Podobnie zadanie czekające na inne zadanie z~tej samej puli może
zakleszczyć program, jeśli wszyscy pracownicy pozostają zajęci
oczekiwaniem. Produkcyjny executor wymaga rozwiązań dla takich sytuacji.}

\subsection{Kiedy potrzebny jest \emph{work stealing}?}

Załóżmy, że przetwarzamy drzewo lub dzielimy dużą macierz na bloki.
Część bloków kończy się po mikrosekundach, inne zajmują milisekundy.
Statyczny przydział zadań może zostawić jeden wątek przeciążony,
gdy inne nie mają już pracy. Przechwytywanie pracy pomiędzy lokalnymi
kolejkami poprawia \emph{równoważenie obciążenia}.

Typowy projekt stosuje:
\begin{enumerate}
\item właścicielski dostęp pracownika do jednego końca jego kolejki;
\item kradzież elementów przez inne wątki z~drugiego końca;
\item atomowe indeksy i~protokół obsługujący konflikt właściciela
      ze złodziejem o~ostatni element;
\item kontrolę czasu życia tablicy zadań oraz jej ewentualnej realokacji.
\end{enumerate}

Takie struktury są znacznie bardziej subtelne niż \code{std::deque}
z~mutexem. W~praktyce warto użyć sprawdzonego executora/biblioteki,
a~własną implementację traktować jako osobny, formalnie weryfikowany projekt.

\section{Anulowanie, sygnały stopu i~punkty oczekiwania}
\label{sec:cpp16-cancel}

Kooperacyjne zatrzymywanie w~C++20 jest oparte na
\code{std::stop_source}, \code{std::stop_token} i~\code{std::stop_callback}.
Żądanie zatrzymania nie zabija wątku; umożliwia mu dobrowolne
zakończenie po zweryfikowaniu stanu i~odpowiednim uporządkowaniu zasobów.

\begin{lstlisting}[caption={Wybudzanie i anulowanie czekania (C++20).}]
#include <condition_variable>
#include <iostream>
#include <mutex>
#include <thread>

int main() {
    std::mutex mutex;
    std::condition_variable_any cv;
    bool ready = false;

    std::jthread worker([&](std::stop_token stop) {
        std::unique_lock lock(mutex);
        bool completed = cv.wait(lock, stop, [&] { return ready; });
        std::cout << (completed ? "ready\n" : "cancelled\n");
    });

    worker.request_stop();
    worker.join();
}
\end{lstlisting}

W~tym przeciążeniu \code{wait} zwraca wartość predykatu po zakończeniu
czekania. Jeśli żądanie zatrzymania nastąpi przed spełnieniem
predykatu, przykład wypisze \code{cancelled}. Jeśli stan zmieniłby
się w~międzyczasie, wynik mógłby być inny. Zwykła
\code{std::condition_variable} nie posiada w~C++20 odpowiednika
tego przeciążenia z~\code{stop_token}.

\subsection{Anulowanie i~bezpieczeństwo wyjątkowe}

W~dobrym API warto określić, co dzieje się z~wynikiem po anulowaniu:
\begin{itemize}
\item czy zwracany jest częściowy wynik;
\item czy do \code{future} trafia własny wyjątek oznaczający anulowanie;
\item czy potrzebne jest odwrócenie skutków wcześniej wykonanych operacji;
\item jak zapewnić, że wszystkie zależne zadania zostaną zakończone
      przed zniszczeniem współdzielonych danych.
\end{itemize}

Żadne z~tych pytań nie ma jednej odpowiedzi dla wszystkich bibliotek.
Sama metoda \code{future.wait_for()} nie anuluje zadań, a~korutyna
z~rozdziału~\ref{chap:cpp15-coroutines} wymaga mechanizmu propagacji żądania anulowania
przez awaitery i~executor.

\section{Fałszywe współdzielenie (\emph{false sharing}) i~cache coherence}
\label{sec:cpp16-false-sharing}

Dwa niezależne atomiki mogą być zapisane blisko siebie w~pamięci,
tak że mieszczą się w~tej samej jednostce przenoszenia danych
pomiędzy pamięciami podręcznymi. Wtedy zapis przez różne rdzenie
może wymagać kosztownych transferów linii cache, choć z~punktu
widzenia C++ nie ma wyścigu danych.

\begin{lstlisting}[caption={Rozdzielenie intensywnie zapisywanych liczników (C++20; stała wyrównania od C++17).}]
#include <atomic>
#include <cstdint>
#include <iostream>
#include <new>
#include <thread>

struct alignas(std::hardware_destructive_interference_size) Counter {
    std::atomic<std::uint64_t> value{0};
};

int main() {
    Counter a, b;
    auto increment = [](Counter& c) {
        for (int i = 0; i < 100000; ++i)
            c.value.fetch_add(1, std::memory_order_relaxed);
    };
    std::jthread first([&] { increment(a); });
    std::jthread second([&] { increment(b); });
    first.join();
    second.join();
    std::cout << a.value << ' ' << b.value << '\n';
}
\end{lstlisting}

\code{std::hardware_destructive_interference_size} jest wskazówką
implementacji dotyczącą rozdzielania obiektów często modyfikowanych
przez różne wątki; w~praktyce jej obsługa zależy od kompilatora
i~biblioteki standardowej. \code{alignas} może zwiększać zużycie
pamięci, zatem nie należy bez pomiarów wyrównywać w~ten sposób
każdego pola. Rozdział~\ref{chap:cpp-memory-performance} omawia szerzej cache locality i~organizację
danych AoS/SoA.

\subsection{Co mierzyć przy ocenie skalowalności?}

Pojedynczy czas wykonania jest mało informacyjny. Przy porównywaniu
wersji jednowątkowej i~wielowątkowej warto mierzyć:
\begin{itemize}
\item przepustowość (operacje/sekundę), a~nie tylko opóźnienie
      pojedynczego zadania;
\item \emph{speedup} \(S(p)=T(1)/T(p)\), dla \(p\) wątków;
\item rozkład opóźnień, np. medianę i~wysokie percentyle;
\item liczbę alokacji, przełączeń kontekstu, braków w~cache i~konfliktów CAS;
\item oddzielnie okres rozgrzewki oraz właściwy okres pomiarowy.
\end{itemize}

Według prawa Amdahla, jeśli część \(s\) programu musi pozostać
sekwencyjna, to w~idealizacji
\[
S(p)\leq\frac{1}{s+\frac{1-s}{p}}.
\]
Wzór ten nie obejmuje kosztów synchronizacji, a~więc rzeczywiste
przyspieszenie może być jeszcze niższe. Jednocześnie program
z~dominującym kosztem dostępu do pamięci może nie skalować się
nawet przy bardzo małej części sekwencyjnej.

\section{Linearizowalność i~niezmienniki współdzielonych struktur}
\label{sec:cpp16-linearizability}

Dla współbieżnych kolekcji często potrzebujemy mocniejszej własności
niż brak wyścigów danych: \emph{linearizowalności}.
Operacja jest linearizowalna, jeśli można przypisać jej jeden
umowny moment pomiędzy wywołaniem i~zakończeniem, w~którym
wywołanie ma taki sam skutek jak operacja wykonana sekwencyjnie.

Dla prostej atomowej inkrementacji punktem linearyzacji jest
sam RMW. Dla kolejki algorytm jest bardziej subtelny:
\begin{itemize}
\item producent najpierw zapisuje komórkę, a~następnie publikuje
      nowy \code{head};
\item konsument najpierw odczytuje i~przetwarza komórkę, a~następnie
      publikuje nowy \code{tail};
\item sama komórka nie musi być atomowa, o~ile zachowana jest
      właściwa synchronizacja i~protokół własności miejsca w~buforze.
\end{itemize}

W~algorytmach złożonych, np. z~dwoma atomikami \code{size} i~\code{head},
każde z~pól może być wolne od wyścigu danych, ale ich jednoczesny
odczyt nie musi stanowić spójnego snapshotu. Osobny \emph{niezmiennik}
struktury wymaga własnego uzasadnienia.

\subsection{Dlaczego prosty \emph{seqlock} nie wystarczy dla zwykłych danych?}

Popularny pomysł polega na tym, że pisarz zwiększa licznik wersji
przed i~po aktualizacji obiektu, a~czytelnik ponawia odczyt, jeżeli
wersja zmieniła się podczas odczytu. Jeżeli pola obiektu są zwykłymi
zmiennymi C++, czytelnik nadal może czytać je jednocześnie z~zapisem,
co oznacza \emph{data race}. Ponowienie odczytu \emph{po fakcie}
nie naprawia niezdefiniowanego zachowania, które już nastąpiło.

Poprawny wariant dla przenośnego C++ musi zastosować odpowiednią
synchronizację lub atomowe pola danych, a~nie tylko atomowy licznik
wersji. Alternatywą jest pokazana wcześniej publikacja niezmiennych
snapshotów.

\section{Współbieżność a~parallel STL, Ranges i~OpenMP}
\label{sec:cpp16-interop}

Nie każda równoległa praca wymaga własnej kolejki czy atomowych
liczników. Gdy zadanie polega na niezależnym przetwarzaniu elementów,
lepszą abstrakcją może być algorytm równoległy lub dyrektywa OpenMP.

\subsection{Polityki wykonania standardowych algorytmów}

Nagłówek \code{<execution>} udostępnia od C++17 między innymi:
\begin{description}[style=nextline]
\item[\code{std::execution::seq}] Wykonanie sekwencyjne.
\item[\code{std::execution::par}] Implementacja może wykonywać pracę
      równolegle; nie daje gwarancji, że rzeczywiście użyje wielu wątków.
\item[\code{std::execution::par_unseq}] Zezwala dodatkowo na
      niesekwencyjne wykonywanie, w~tym wektoryzację; nakłada ważne
      ograniczenia na wykonywane funkcje i~synchronizację.
\end{description}

\begin{lstlisting}[caption={Równoległa transformacja bez współdzielonych zapisów (C++17).}]
#include <algorithm>
#include <execution>
#include <numeric>
#include <vector>

int main() {
    std::vector<double> a(10000), b(a.size());
    std::iota(a.begin(), a.end(), 0.0);
    std::transform(std::execution::par,
                   a.begin(), a.end(), b.begin(),
                   [](double x) { return x * x; });
}
\end{lstlisting}

Kod nie zakłada kolejności przetwarzania elementów. Wyjściowy zakres
\code{b} nie nakłada się z~\code{a}, a~każde wywołanie funkcji
zapisuje inny element. Rzeczywiste zrównoleglenie zależy od
implementacji biblioteki standardowej i~konfiguracji środowiska;
niektóre kompilacje wymagają dodatkowej biblioteki wykonawczej.

\subsection{C++ i~OpenMP: dwa poziomy abstrakcji}

OpenMP, omawiany szczegółowo w~rozdziale~\ref{chap:concurrency}, umożliwia między innymi
równoległe pętle, redukcje i~zależności pomiędzy zadaniami.
W~obliczeniach numerycznych typowe rozwiązanie wygląda tak:

\begin{lstlisting}[caption={Wektorowa redukcja z OpenMP; wymagane -fopenmp.}]
#include <iostream>
#include <vector>

int main() {
    const std::vector<double> x(100000, 1.0);
    double sum = 0.0;

    #pragma omp parallel for reduction(+:sum) schedule(static)
    for (int i = 0; i < static_cast<int>(x.size()); ++i) {
        sum += x[static_cast<std::size_t>(i)];
    }
    std::cout << sum << '\n';
}
\end{lstlisting}

Kodowanie redukcji przez \code{reduction} pozwala uniknąć współdzielenia
jednego licznika zapisywanego przez wszystkie wątki. Kolejność
sumowania liczb zmiennoprzecinkowych może być odmienna niż przy
obliczeniu sekwencyjnym, a~zatem wynik może różnić się przez
zaokrąglenia. Liczbę wątków można ustawić m.in. przez
\code{OMP_NUM_THREADS}; \code{OMP_SCHEDULE} wpływa na harmonogram
tylko przy \code{schedule(runtime)}. Powyższy przykład jawnie
wybiera \code{schedule(static)}.

\nota{Mieszanie pul wątków, OpenMP, równoległych algorytmów STL
oraz zadań asynchronicznych może prowadzić do \emph{oversubscription}:
liczba aktywnych wątków znacząco przewyższa dostępne zasoby CPU.
Przy nakładaniu kilku warstw równoległości trzeba świadomie
kontrolować budżet wątków i~mierzyć efekt końcowy.}

\section{Diagnostyka: kompilatory, sanitizatory, testy przeciążeniowe}
\label{sec:cpp16-testing}

\subsection{Kompilacja i~uruchamianie przykładów}

Dla przykładów wykorzystujących \code{std::jthread}, \code{atomic::wait}
i~\code{atomic<shared_ptr<T>>} wybieramy tryb C++20:

\begin{lstlisting}[language=bash,caption={Kompilacja wielowątkowego programu C++20.}]
g++ -std=c++20 -O2 -Wall -Wextra -Wpedantic -pthread file.cpp -o app
./app
\end{lstlisting}

Dla przykładów OpenMP w~GCC lub odpowiednio skonfigurowanym Clangu:
\begin{lstlisting}[language=bash,caption={Włączenie OpenMP.}]
g++ -std=c++20 -O2 -Wall -Wextra -fopenmp file.cpp -o app
./app
\end{lstlisting}

\subsection{ThreadSanitizer i~inne narzędzia}

\emph{ThreadSanitizer} (TSan) potrafi wykrywać wiele wyścigów danych
i~wybrane błędy synchronizacji w~uruchamianym programie:

\begin{lstlisting}[language=bash,caption={Budowanie i uruchamianie z ThreadSanitizer.}]
g++ -std=c++20 -O1 -g -pthread -fsanitize=thread file.cpp -o app_tsan
./app_tsan
\end{lstlisting}

Nie wszystkie środowiska poprawnie wspierają TSan; szczególnie
przy nietypowych bibliotekach wykonawczych, OpenMP i~własnych
mechanizmach synchronizacji należy weryfikować zakres instrumentacji.
AddressSanitizer i~UndefinedBehaviorSanitizer mogą dodatkowo wykrywać
część błędów czasu życia oraz naruszeń zasad języka, lecz nie są
zamiennikami TSan ani formalnego dowodu synchronizacji.

\subsection{Jak testować algorytm współbieżny?}

Dla poważniejszej struktury danych zaleca się następującą procedurę:
\begin{enumerate}
\item \textbf{Zdefiniuj kontrakt}: ilu producentów i~konsumentów,
      dozwolone operacje oraz zasady zamknięcia i~niszczenia struktury.
\item \textbf{Zidentyfikuj własność danych}: kto może zapisywać dane
      w~każdej chwili; kiedy przechodzi własność konkretnej komórki.
\item \textbf{Narysuj relacje}: wskaż atomik, który publikuje zapis
      i~odczyt, który tę publikację odbiera.
\item \textbf{Wskaż punkt linearyzacji}: jeśli API ma być linearizowalne,
      określ, kiedy operacja staje się widoczna logicznie.
\item \textbf{Przeanalizuj zakończenie}: co dzieje się przy anulowaniu,
      wyjątku, opóźnieniu konsumenta i~niszczeniu obiektu.
\item \textbf{Testuj wielokrotnie}: zmieniaj liczbę wątków, rozmiary
      danych i~częstotliwość przełączeń; uruchamiaj TSan tam,
      gdzie jest dostępny.
\item \textbf{Benchmarkuj osobno}: włącz optymalizacje i~mierz
      przy realistycznym obciążeniu, ale nie myl szybkiego testu
      z~dowodem poprawności.
\end{enumerate}

\section{Typowe błędy i~zalecenia projektowe}
\label{sec:cpp16-guidelines}

\begin{center}
\small
\begin{tabularx}{\textwidth}{@{}lX@{}}
\toprule
\textbf{Błąd} & \textbf{Dlaczego jest groźny / co wybrać} \\
\midrule
\code{volatile} jako synchronizacja & Nie eliminuje \emph{data race};
użyj mutexa lub atomików z~właściwym porządkiem. \\
\code{relaxed} do publikacji zwykłych danych & Nie wytwarza
relacji synchronizacji; zazwyczaj potrzebne \code{release/acquire}. \\
Bezrefleksyjne \code{seq_cst} & Zwykle poprawniejsze niż źle
dobrane słabsze porządki, ale nie zastąpi logicznego kontraktu obiektu. \\
Kasowanie węzła po udanym CAS & Inny wątek może nadal używać
wskaźnika; zaprojektuj odzyskiwanie pamięci. \\
Dwa atomiki jako ,,jedna transakcja'' & Brak automatycznego
atomowego snapshotu pary; użyj blokady lub niezmiennej wersji. \\
\emph{Spinlock} dla długiej pracy & Prowadzi do rywalizacji
o~CPU; zwykle użyj \code{std::mutex}. \\
Nieograniczona kolejka zadań & Może zużyć pamięć; dodaj limit
i~mechanizm \emph{backpressure}. \\
Nadmierna liczba wątków & Konkurencja o~rdzenie i~cache;
kontroluj równoległość i~profiluj. \\
Mieszanie OpenMP i~własnej puli bez planu & Może prowadzić
do \emph{oversubscription} i~spadku wydajności. \\
\bottomrule
\end{tabularx}
\end{center}

\subsection{Praktyczne reguły wyboru}

\begin{itemize}
\item \textbf{Mutex} wybieraj jako punkt wyjścia, jeśli współdzielisz
      kilka powiązanych pól i~potrzebujesz łatwo sprawdzalnych niezmienników.
\item \textbf{Atomik \code{relaxed}} jest właściwy dla prostego licznika
      niezależnego od publikacji innych danych.
\item \textbf{\code{release/acquire}} stosuj, gdy istnieje wyraźny
      punkt przekazania danych lub własności pomiędzy wątkami.
\item \textbf{Lock-free} rozważaj dopiero wtedy, gdy pomiary pokażą
      realny koszt blokad, a~projekt ma przejrzysty dowód poprawności
      oraz plan odzyskiwania pamięci.
\item \textbf{OpenMP / parallel STL} wybieraj do operacji na niezależnych
      elementach, redukcji i~równoległości danych.
\item \textbf{Thread pool} wybieraj do dużej liczby krótkich zadań,
      gdy potrzebujesz ograniczać liczbę wątków, priorytety lub
      przepustowość; osobno zdefiniuj zasady zamykania.
\end{itemize}

\section*{Podsumowanie rozdziału}
\addcontentsline{toc}{section}{Podsumowanie rozdziału}

Zaawansowana współbieżność nie polega na mechanicznym zastępowaniu
mutexów przez \code{std::atomic}. Najpierw trzeba określić
\emph{własność danych i~kontrakt algorytmu}, następnie uzasadnić
relacje \emph{happens-before} i~wyznaczyć punkty widoczności
poszczególnych operacji. Dopiero wtedy warto wybierać porządki
pamięciowe, typy kolejek, planistę i~sposób odzyskiwania pamięci.

W~szczególności brak wyścigu danych nie oznacza jeszcze poprawności
logicznej, a~nieblokujący dostęp do atomika nie gwarantuje
nieblokującego ani szybkiego zachowania całego programu. Bezpieczna
współbieżność to połączenie \emph{modelu formalnego}, \emph{kontroli
czasu życia} i~\emph{rzetelnego pomiaru}.

\section*{Literatura i dokumentacja}
\addcontentsline{toc}{section}{Literatura i dokumentacja}
\begin{itemize}
\item Anthony Williams, \emph{C++ Concurrency in Action}, wyd.~2,
      Manning, 2019.
\item ISO/IEC 14882:2020, rozdziały o~modelu pamięci, wątkach,
      atomikach i~synchronizacji.
\item \url{https://en.cppreference.com/w/cpp/atomic/memory_order}
      --- porządki pamięci i~relacje synchronizacji.
\item \url{https://en.cppreference.com/w/cpp/atomic/atomic}
      --- typy atomowe i~operacje RMW.
\item \url{https://en.cppreference.com/w/cpp/atomic/atomic_ref}
      --- atomowe odwołania do istniejących obiektów.
\item \url{https://en.cppreference.com/w/cpp/thread/jthread}
      --- \code{jthread}, czas życia i~kooperacyjne zatrzymywanie.
\item OpenMP Architecture Review Board,
      \emph{OpenMP Application Programming Interface};
      \url{https://www.openmp.org/specifications/}.
\end{itemize}

% Kazdy listing opisany jako kompletny program jest samodzielny.

\chapter{Zaawansowana współbieżność, model pamięci i algorytmy bez blokad}
\label{chap:cpp16-advanced-concurrency}

\section{Cel rozdziału i powiązanie z podstawami}

W rozdziale~\ref{chap:concurrency} poznaliśmy wątki, mutexy, zmienne warunkowe, operacje
atomowe, \code{std::async}, pule wątków i~OpenMP. Sama znajomość tych
interfejsów nie wystarcza jednak do samodzielnego projektowania
algorytmów współbieżnych. Trzeba odpowiedzieć na trudniejsze pytania:
\begin{itemize}
\item Kiedy zapis wykonany przez jeden wątek \emph{musi} być widoczny dla drugiego?
\item Które wyniki są dozwolone przez model pamięci, nawet gdy wydają się
      sprzeczne z intuicją opartą na sekwencyjnym wykonywaniu?
\item Czym różni się atomowość pojedynczej operacji od niepodzielności
      całego algorytmu?
\item Jak zaprojektować \emph{postęp} algorytmu i~bezpiecznie zwalniać pamięć,
      gdy inne wątki mogą jeszcze posiadać wskaźniki do obiektów?
\item Kiedy rozwiązanie bez mutexa rzeczywiście poprawia wydajność?
\end{itemize}

Skupimy się na \emph{uzasadnianiu poprawności}, a~nie tylko na składni.
W~szczególności przedstawimy proste testy modelu pamięci (\emph{litmus tests}),
algorytm kolejki SPSC, publikację niezmiennych struktur danych, problem ABA,
techniki odzyskiwania pamięci oraz praktyczne zasady pomiaru i~diagnostyki.
Nie zakładamy, że każdy \code{std::atomic<T>} jest realizowany bez blokad.

\nota{Wszystkie gwarancje sformułowane w~rozdziale dotyczą abstrakcyjnego
modelu języka C++. Rzeczywisty procesor może wykonywać instrukcje
w~innej kolejności i~stosować bufory zapisów, ale program musi być
poprawny na gruncie standardu, a~nie tylko na jednym komputerze.}

\section{Formalny porządek zdarzeń w~wielowątkowym programie}
\label{sec:cpp16-formal}

\subsection{Zdarzenia i~relacje \emph{sequenced-before}, \emph{synchronizes-with}}

Rozpatrzmy pewne operacje na pamięci jako \emph{zdarzenia}. Wśród relacji
używanych do opisu ich uporządkowania najważniejsze są:
\begin{description}[style=nextline]
\item[\emph{Sequenced-before} (\(\mathrm{sb}\))]
  Porządek określony regułami języka pomiędzy wybranymi operacjami
  w~obrębie jednego wątku. Nie należy go mylić z~dowolną kolejnością
  pojawiania się instrukcji w~kodzie źródłowym: język zawiera także
  wyrażenia z~nieokreśloną kolejnością ewaluacji.
\item[\emph{Synchronizes-with} (\(\mathrm{sw}\))]
  Synchronizacja pomiędzy odpowiednimi operacjami różnych wątków:
  na przykład zapis \code{release} i~odczyt \code{acquire}, który czyta
  wartość z~odpowiedniego zapisu lub jego sekwencji \emph{release}.
  Podobną synchronizację zapewniają właściwe operacje na mutexach.
\item[\emph{Happens-before} (\(\mathrm{hb}\))]
  Relacja porządkująca zdarzenia, budowana w~typowych przykładach
  przez \(\mathrm{sb}\), \(\mathrm{sw}\) i~ich przechodnie złożenia.
  Jeśli \(A\;\mathrm{hb}\;B\), skutki \(A\) są uporządkowane przed \(B\)
  w~sensie formalnego modelu; nie jest to stwierdzenie o~czasie zegarowym.
\end{description}

Dla przejrzystości będziemy używać następującego schematu:
\[
\underbrace{A\;\xrightarrow{\mathrm{sb}}\;R}_{\text{producent}}
\quad\xrightarrow{\mathrm{sw}}\quad
\underbrace{Q\;\xrightarrow{\mathrm{sb}}\;B}_{\text{konsument}}
\quad\Longrightarrow\quad A\;\xrightarrow{\mathrm{hb}}\;B.
\]
Operacja \(R\) jest publikacją, \(Q\) odpowiadającym jej odbiorem;
\(A\) i~\(B\) mogą dotyczyć zwykłej, nieatomowej zmiennej.

\subsection{Atomowość nie zastępuje synchronizacji}

W~przeciwieństwie do zwykłego odczytu lub zapisu operacje na
\code{std::atomic<T>} nie powodują wyścigu danych pomiędzy sobą.
Nie oznacza to jednak, że każdy odczyt atomowy widzi ostatni zapis
według zegara ściennego. Taki globalny \emph{ostatni zapis} zwykle
w~ogóle nie istnieje jako zdefiniowane pojęcie modelu.

Z~kolei jeśli dwa wątki odczytują i~modyfikują to samo miejsce pamięci
nieatomowo, a~konfliktowe operacje nie są uporządkowane relacją
\emph{happens-before}, program ma \emph{data race} i~niezdefiniowane zachowanie.
Samo oznaczenie obiektu przez \code{volatile} nie rozwiązuje tego problemu.

\subsection{Porządek modyfikacji (\emph{modification order}) i~spójność}

Dla \emph{każdego pojedynczego} obiektu atomowego istnieje całkowity
porządek jego modyfikacji (\emph{modification order}). Nie wynika z~tego
jeden wspólny porządek modyfikacji \emph{wszystkich} atomików.
Ważne konsekwencje to między innymi:
\begin{itemize}
\item zapis do \code{x} uporządkowany przed innym zapisem do \code{x}
      musi być odpowiednio umieszczony przed nim w~porządku modyfikacji;
\item jeśli odczyt \code{x} jest uporządkowany przed kolejnym odczytem
      \code{x} w~tym samym wątku, drugi odczyt nie może cofnąć się
      do wcześniejszej modyfikacji w~porządku tego atomika;
\item ograniczenia spójności dla jednego atomika nie zapewniają
      automatycznie spójnego obrazu \emph{dwóch różnych} atomików.
\end{itemize}

\section{Testy modelu pamięci: wyniki, które zaskakują}
\label{sec:cpp16-litmus}

\subsection{\emph{Store buffering}: dwa odczyty zera}

Załóżmy, że \code{x} i~\code{y} początkowo są równe zero, a~wątki
wykonują następujące operacje:
\[
\begin{array}{c|c}
\text{Wątek 1} & \text{Wątek 2}\\ \hline
x\gets 1 & y\gets 1\\
r_1\gets y & r_2\gets x
\end{array}
\]

\begin{lstlisting}[caption={Store buffering: kompletna, wolna od wyścigu danych próba (C++20).}]
#include <atomic>
#include <iostream>
#include <thread>

int main() {
    std::atomic<int> x{0}, y{0};
    int r1 = -1, r2 = -1;

    std::jthread a([&] {
        x.store(1, std::memory_order_relaxed);
        r1 = y.load(std::memory_order_relaxed);
    });
    std::jthread b([&] {
        y.store(1, std::memory_order_relaxed);
        r2 = x.load(std::memory_order_relaxed);
    });
    a.join();
    b.join();
    std::cout << r1 << ' ' << r2 << '\n';
}
\end{lstlisting}

W~tym programie \(r_1=r_2=0\) jest \textbf{dozwolonym} wynikiem
modelu C++ dla \code{relaxed}; nie ma gwarancji, że zostanie
zaobserwowany w~praktycznym uruchomieniu. Jeśli \emph{wszystkie cztery}
operacje atomowe zamienimy na \code{memory_order_seq_cst}, wynik
\((0,0)\) stanie się niemożliwy: wspólny porządek operacji
sekwencyjnie spójnych musiałby przeczyć kolejności w~każdym z~wątków.

\nota{Wynik pojedynczego eksperymentu, nawet powtórzony milion razy,
nie dowodzi niedopuszczalności innego zachowania. Test tego rodzaju
to materiał do rozumowania nad modelem, a~nie pełny dowód poprawności.}

\subsection{\emph{Message passing}: publikacja zwykłych danych}

\begin{lstlisting}[caption={Publikacja danych przy użyciu release/acquire (C++11).}]
#include <atomic>
#include <cassert>
#include <thread>

int main() {
    int payload = 0;
    std::atomic<bool> published{false};

    std::thread producer([&] {
        payload = 42;
        published.store(true, std::memory_order_release);
    });
    std::thread consumer([&] {
        while (!published.load(std::memory_order_acquire)) {
            std::this_thread::yield();
        }
        assert(payload == 42);
    });

    producer.join();
    consumer.join();
}
\end{lstlisting}

Odczyt flagi, który widzi \code{true}, synchronizuje się z~zapisem
\code{release}. Zapis \code{payload = 42} następuje \emph{happens-before}
odczytu \code{payload} przez konsumenta. Flaga \emph{publikuje}
wcześniejszy stan; nie należy po publikacji modyfikować
\code{payload} bez dalszej synchronizacji.

\section{Porządki atomowe i~ich dokładne zastosowania}
\label{sec:cpp16-orders}

\begin{center}
\small
\begin{tabularx}{\textwidth}{@{}lX@{}}
\toprule
\textbf{Porządek} & \textbf{Typowe zastosowanie i~ograniczenie} \\
\midrule
\code{relaxed} & Liczniki, statystyki, pola kontrolne, gdy
sama atomowość wystarcza, a~relacja synchronizacji wynika skądinąd. \\
\code{release} & Zapis lub operacja RMW publikująca wcześniejsze efekty
wątku. \\
\code{acquire} & Odczyt lub operacja RMW odbierająca stan od publikującego wątku. \\
\code{acq_rel} & Operacja odczyt--modyfikacja--zapis wymagająca obu kierunków. \\
\code{seq_cst} & Operacje z~dodatkową gwarancją pojedynczego wspólnego
porządku dla operacji sekwencyjnie spójnych. \\
\bottomrule
\end{tabularx}
\end{center}

\code{memory_order_consume} ma długą i~skomplikowaną historię
standaryzacji. W~praktycznym kodzie C++20 należy preferować
\code{acquire} zamiast opierania przenośnej poprawności na
\code{consume}. Uwaga: porządek \code{release} jest właściwy dla
zapisu, \code{acquire} dla odczytu, a~\code{acq_rel} dla operacji
odczyt--modyfikacja--zapis; nie wszystkie porządki są poprawne dla
każdego przeciążenia funkcji atomowej.

\subsection{Operacje \emph{read--modify--write}}

Metody \code{fetch_add}, \code{fetch_sub}, \code{exchange} i~udany
\code{compare_exchange_*} wykonują operację RMW. Są one niepodzielne
w~odniesieniu do danego atomika, ale nie tworzą automatycznie
transakcji obejmującej inne pola.

\begin{lstlisting}[caption={Przydzielanie kolejnych numerów: wystarcza relaxed (C++11).}]
#include <atomic>
#include <cassert>
#include <thread>

int main() {
    std::atomic<int> next_id{0};
    int first = -1, second = -1;
    std::thread a([&] {
        first = next_id.fetch_add(1, std::memory_order_relaxed);
    });
    std::thread b([&] {
        second = next_id.fetch_add(1, std::memory_order_relaxed);
    });
    a.join();
    b.join();
    assert(first != second);
    assert(first + second == 1);
}
\end{lstlisting}

Uzyskujemy unikatowe wartości, bo RMW mają jednoznaczny porządek
modyfikacji \code{next_id}. Nie uzyskujemy przez to żadnej dodatkowej
synchronizacji niezwiązanych danych.

\subsection{\emph{Release sequence}: łańcuch operacji RMW}

W~C++20 zapis z~porządkiem \code{release} może rozpoczynać sekwencję
\emph{release} składającą się z~następujących po nim operacji RMW
na \emph{tym samym} atomiku, wykonywanych także przez inne wątki.
Odczyt \code{acquire} wartości pochodzącej z~tej sekwencji może
synchronizować się z~początkowym zapisem \code{release}.

\begin{lstlisting}[caption={Przekazanie publikacji przez relaxed RMW (C++20).}]
#include <atomic>
#include <cassert>
#include <thread>

int main() {
    int payload = 7;
    std::atomic<int> state{0};

    std::thread producer([&] {
        payload = 42;
        state.store(1, std::memory_order_release);
    });
    std::thread intermediary([&] {
        int expected = 0;
        while ((expected = state.load(std::memory_order_acquire))
               == 0) {
            std::this_thread::yield();
        }
        state.fetch_add(1, std::memory_order_relaxed);
    });
    std::thread consumer([&] {
        while (state.load(std::memory_order_acquire) != 2) {
            std::this_thread::yield();
        }
        assert(payload == 42);
    });

    producer.join();
    intermediary.join();
    consumer.join();
}
\end{lstlisting}

Wątek pośredni czeka, aż stan osiągnie 1, i~wykonuje RMW, który
zapisuje 2. Odczyt wartości 2 z~\code{acquire} synchronizuje się
z~pierwotnym \code{release} przez sekwencję \emph{release}.
Nie można w~ogólności zastąpić \code{fetch_add} zwykłym
\code{store(2, relaxed)} i~utrzymać tej samej argumentacji.

\section{CAS: porównaj i~zamień, pętle ponawiania i~porządek niepowodzenia}

\begin{center}\small
\begin{tabularx}{\textwidth}{@{}p{0.38\textwidth}X@{}}
\toprule
Operacja atomowa & Działanie \\
\midrule
\code{compare_exchange_strong} & Jednorazowo porównuje i próbuje podmienić wartość; nie zgłasza pozornego niepowodzenia. \\
\code{compare_exchange_weak} & Może zgłosić pozorne niepowodzenie; nadaje się do pętli ponawiania. \\
\code{expected} przy CAS & Przy niepowodzeniu otrzymuje rzeczywistą wartość atomika, którą trzeba uwzględnić przy kolejnej próbie. \\
\code{memory_order_acq_rel} & Przy udanej operacji odczytująco-zapisującej łączy własności acquire i release. \\
\code{memory_order_relaxed} & Zapewnia atomowość bez dodatkowej synchronizacji z innymi danymi. \\
\bottomrule
\end{tabularx}
\end{center}
\label{sec:cpp16-cas}

\code{compare_exchange_weak(expected, desired, success, failure)}
atomowo porównuje bieżącą wartość z~\code{expected}. Gdy porównanie
się powiedzie, zapisuje \code{desired}; w~przeciwnym razie aktualizuje
\code{expected} do wartości zastanej w~atomiku. Wersja \code{weak}
może ponadto nie powieść się pozornie (\emph{spurious failure}).
W~pętli jest zwykle dobrym wyborem.

\begin{lstlisting}[caption={Atomowe maksimum z~pętlą CAS (C++11).}]
#include <atomic>
#include <cassert>

void atomic_max(std::atomic<int>& target, int candidate) {
    int old = target.load(std::memory_order_relaxed);
    while (old < candidate &&
           !target.compare_exchange_weak(
               old, candidate,
               std::memory_order_relaxed,
               std::memory_order_relaxed)) {
        // Przy niepowodzeniu old jest aktualizowane.
    }
}

int main() {
    std::atomic<int> maximum{10};
    atomic_max(maximum, 8);
    atomic_max(maximum, 17);
    assert(maximum.load() == 17);
}
\end{lstlisting}

Parametr \code{failure} opisuje porządek operacji \emph{odczytu}
w~przypadku niepowodzenia. Nie może to być \code{release} ani
\code{acq_rel}; musi też spełniać ograniczenia względem porządku
sukcesu określone standardem. Typowy wzorzec publikacji nowej
wartości to sukces z~\code{acq_rel} i~niepowodzenie z~\code{acquire}
(albo \code{relaxed}, jeśli tak wynika z~dowodu).

\subsection{\code{weak} kontra \code{strong}}

\code{compare_exchange_strong} nie zgłasza pozornego niepowodzenia
przy poprawnie porównanych wartościach. Może być użyteczne dla
pojedynczej próby; \code{weak} bywa wygodniejsze w~pętlach CAS.
Nie ma uniwersalnej gwarancji, że jedna z~tych wersji zawsze działa
szybciej na danej platformie.

\section{Bariery pamięci i~\texttt{std::atomic\_thread\_fence}}
\label{sec:cpp16-fences}

Bariera \code{std::atomic_thread_fence(order)} jest operacją
porządkowania pamięci, która sama nie odczytuje ani nie zapisuje
konkretnego atomika. Synchronizacja z~innym wątkiem wciąż wymaga
odpowiednich operacji atomowych łączących obie strony. Samo
wstawienie \code{fence} obok zwykłej zmiennej nie usuwa wyścigu danych.

\begin{lstlisting}[caption={Publikacja ze zwolnieniową i~przejmującą barierą pamięci (C++11).}]
#include <atomic>
#include <cassert>
#include <thread>

int main() {
    int payload = 0;
    std::atomic<bool> flag{false};

    std::thread producer([&] {
        payload = 42;
        std::atomic_thread_fence(std::memory_order_release);
        flag.store(true, std::memory_order_relaxed);
    });
    std::thread consumer([&] {
        while (!flag.load(std::memory_order_relaxed)) {
            std::this_thread::yield();
        }
        std::atomic_thread_fence(std::memory_order_acquire);
        assert(payload == 42);
    });

    producer.join();
    consumer.join();
}
\end{lstlisting}

To przykład synchronizacji \emph{fence--fence}: odczyt atomika
widzi zapis następujący po barierze \code{release}, a~bariera
\code{acquire} jest wykonywana po tym odczycie. W~nowym kodzie
zwykle prostsze i~czytelniejsze jest zastosowanie pary operacji
\code{release}/\code{acquire} bez jawnych barier. Bariery nie są
substytutem przemyślanego protokołu komunikacji.

\section{Gwarancje postępu: \emph{blocking}, \emph{lock-free} i~\emph{wait-free}}
\label{sec:cpp16-progress}

Wielowątkowy algorytm można analizować pod względem \emph{bezpieczeństwa}
(\emph{safety}: nic niedozwolonego się nie dzieje) i~\emph{żywotności}
(\emph{liveness}: coś pożądanego w~końcu się wydarzy). Są to różne kwestie.

\begin{description}[style=nextline]
\item[\emph{Blocking}]
  Wątek może oczekiwać na odblokowanie mutexa przez konkretny inny wątek.
  Jeśli właściciel przestanie robić postępy, pozostali również mogą utknąć.
\item[\emph{Lock-free}]
  W~standardowym ujęciu postępu algorytmowego, mimo rywalizacji wątków,
  system jako całość może kontynuować wykonywanie operacji:
  niepowodzenia jednej próby CAS zwykle oznaczają, że inny uczestnik
  wykonał postęp. Nie gwarantuje to postępu \emph{każdego} wątku.
\item[\emph{Wait-free}]
  Każda operacja danego wątku ma ograniczoną liczbę własnych kroków
  niezależnie od działań innych wątków (przy przyjętym modelu).
  To silniejszy warunek niż \emph{lock-free}.
\item[\emph{Obstruction-free}]
  Operacja postępuje, jeśli wątek może działać dostatecznie długo
  bez konfliktów ze strony innych wątków. Jest to słabsza gwarancja.
\end{description}

\emph{Lock-free} nie znaczy ,,zawsze szybciej''. Nie oznacza również
braku aktywnego oczekiwania ani braku zagłodzenia konkretnego wątku.
Odróżniajmy też własność implementacji pojedynczego atomika
(\code{is_lock_free()}) od postępu kompletnego algorytmu.

\subsection{Prosta blokada wirująca (\emph{spinlock})}

\begin{lstlisting}[caption={Poprawna blokada wirująca z flagą atomową; ilustracja (C++20).}]
#include <atomic>
#include <cassert>
#include <mutex>
#include <thread>

class SpinLock {
    std::atomic_flag flag = ATOMIC_FLAG_INIT;
public:
    void lock() noexcept {
        while (flag.test_and_set(std::memory_order_acquire)) {
            // Zamiast stale rywalizowac o cache line, czekamy.
            flag.wait(true, std::memory_order_relaxed);
        }
    }
    void unlock() noexcept {
        flag.clear(std::memory_order_release);
        flag.notify_one();
    }
};

int main() {
    SpinLock lock;
    int counter = 0;
    auto work = [&] {
        for (int i = 0; i < 10000; ++i) {
            std::lock_guard<SpinLock> guard(lock);
            ++counter;
        }
    };
    std::jthread a(work), b(work);
    a.join();
    b.join();
    assert(counter == 20000);
}
\end{lstlisting}

\code{test_and_set} ma semantykę RMW; odblokowanie jest operacją
\code{release}. Metoda \code{wait(true)} może blokować, a~zatem
powyższy \code{SpinLock} \emph{nie} jest algorytmem \emph{lock-free}.
Ponadto nie gwarantuje sprawiedliwości. W~realnym oprogramowaniu
\code{std::mutex} zwykle pozostaje lepszym wyborem, zwłaszcza gdy
sekcja krytyczna może być długa lub obciążenie jest wysokie.

\section{Problem ABA: poprawne CAS nie wystarcza}
\label{sec:cpp16-aba}

Rozważmy wierzchołek stosu reprezentowany wskaźnikiem atomowym.
Wątek \(A\) czyta wskaźnik \(p\) i~zamierza wykonać
\code{compare_exchange} z~wartości \(p\) na \code{p->next}.
W~tym czasie wątek \(B\) może:
\begin{enumerate}
\item zdjąć element \(p\) ze stosu;
\item zdjąć jeszcze jeden element;
\item ponownie umieścić \(p\) na wierzchołku stosu.
\end{enumerate}

Wartość wskaźnika uległa zmianie \(A\to B\to A\) (stąd \emph{ABA}),
ale CAS wątku \(A\) może zobaczyć ponownie tę samą wartość \(p\).
W~międzyczasie struktura stosu mogła się całkowicie zmienić.

\begin{center}
\(\boxed{\text{head}=p}\;\longrightarrow\;
  \boxed{\text{head}=q}\;\longrightarrow\;
  \boxed{\text{head}=p}\)
\end{center}

Dwa problemy trzeba odróżnić:
\begin{enumerate}
\item \emph{ABA logiczne}: porównanie tej samej wartości nie wykrywa
      zmian stanu struktury.
\item \emph{Bezpieczeństwo pamięci}: wątek \(A\) może próbować odczytać
      \code{p->next}, gdy inny wątek zniszczył już obiekt \code{*p}.
      Taki odczyt nie staje się poprawny tylko dlatego, że \code{head}
      jest atomowy.
\end{enumerate}

\subsection{Wskaźnik z~licznikiem generacji}

Jedna z~technik polega na porównywaniu pary
\((\text{wskaźnik},\text{licznik wersji})\) zamiast samego wskaźnika.
Wtedy ponowne pojawienie się tego samego adresu nie oznacza równości
całego stanu. Licznik również ma ograniczony zakres i~może się
przepełnić; trzeba uzasadnić brak niebezpiecznego \emph{wrap-around}.

\nota{Nie zakładaj, że \code{std::atomic<SomeStruct>} będzie bez blokad:
implementacja może wykorzystać wewnętrzny mutex lub bibliotekę
atomowych operacji. Możliwość zapisania stanu złożonego w~\code{atomic}
nie rozwiązuje też samodzielnie problemu czasu życia obiektów.}

\subsection{Hazard pointers, epoch-based reclamation i~RCU}

\term{Hazard pointers} pozwalają wątkom publikować informację:
,,zamierzam korzystać z~tego obiektu''. Wątek usuwający obiekt
odkłada go na listę \emph{retired} i~odracza zniszczenie do chwili,
gdy żaden aktywny hazard pointer go nie wskazuje. Poprawna procedura
zwykle musi ponownie zweryfikować odczyt wskaźnika po publikacji
hazardu; samo opublikowanie dowolnego, potencjalnie już nieważnego
adresu nie wystarcza.

\term{Epoch-based reclamation} grupuje operacje w~epoki. Usunięte
obiekty mogą zostać odzyskane dopiero wtedy, gdy wszystkie aktywne
wątki zakończyły okresy korzystania z~danych obejmujące ich usunięcie.
Dłużej zatrzymany wątek może opóźniać zwolnienie pamięci.

\term{Read-copy-update (RCU)} korzysta z~odczytów starych wersji
danych i~publikacji nowych wersji; usuwanie poprzednich odbywa się
po odpowiednim okresie ochronnym (\emph{grace period}). Szczegóły
różnią się pomiędzy implementacjami.

\begin{center}
\small
\begin{tabularx}{\textwidth}{@{}lX@{}}
\toprule
\textbf{Metoda} & \textbf{Kiedy rozważyć} \\
\midrule
Hazard pointers & Wiele odczytów wskaźników, częste usuwanie
węzłów; możliwa kosztowna administracja hazardami. \\
Epoki & Duży ruch, kontrolowane sekcje dostępu; zatrzymany czytelnik
może opóźnić odzyskanie pamięci. \\
RCU & Wzorce ,,wielu czytelników, niewielu pisarzy'', publikacja
niezmiennych wersji danych. \\
\code{shared_ptr} & Gdy łatwość dowodu poprawności jest ważniejsza
od kosztu liczenia referencji; nie obiecuje \emph{lock-free}. \\
\bottomrule
\end{tabularx}
\end{center}

W~C++20 standardowa biblioteka nie udostępnia jeszcze kompletnej,
przenośnej abstrakcji hazard pointerów. Bezpieczna implementacja
wymaga precyzyjnego protokołu pamięciowego, dlatego zamiast podawać
uproszczony, ale błędny stos ze zwalnianiem węzłów, przedstawimy
algorytm bez dynamicznego odzyskiwania pamięci: kolejkę SPSC.

\section{Kolejka SPSC: praktyczny algorytm bez mutexa}
\label{sec:cpp16-spsc}

\emph{Single Producer, Single Consumer} (SPSC) oznacza dokładnie
jednego producenta i~jednego konsumenta. Bufor cykliczny może wtedy
używać dwóch liczników:
\begin{itemize}
\item \code{head} --- numer kolejnego miejsca zapisu, modyfikowany
      wyłącznie przez producenta;
\item \code{tail} --- numer kolejnego miejsca odczytu, modyfikowany
      wyłącznie przez konsumenta.
\end{itemize}

Producent zapisuje komórkę, a~następnie publikuje nowy \code{head}
operacją \code{release}. Konsument odczytuje \code{head} przez
\code{acquire}, a~więc może bezpiecznie zobaczyć zawartość wcześniej
zapisanej komórki. Analogicznie konsument zwalnia komórkę, publikując
nowy \code{tail}, a~producent odbiera tę informację operacją
\code{acquire}. Warunek pełności bufora o~pojemności \(C\) to
\(\mathrm{head}-\mathrm{tail}=C\), a~pustości:
\(\mathrm{head}=\mathrm{tail}\).

\begin{lstlisting}[caption={Kompletna kolejka SPSC dla int, ze statyczną pamięcią (C++20).}]
#include <array>
#include <atomic>
#include <cassert>
#include <cstddef>
#include <thread>

// Tylko jeden producent i dokladnie jeden konsument.
template<std::size_t Capacity>
class SpscQueue {
    static_assert(Capacity > 0);
    static_assert((Capacity & (Capacity - 1)) == 0);
    std::array<int, Capacity> buffer_{};
    std::atomic<std::size_t> head_{0};
    std::atomic<std::size_t> tail_{0};
public:
    bool try_push(int value) {
        const auto head = head_.load(std::memory_order_relaxed);
        const auto tail = tail_.load(std::memory_order_acquire);
        if (head - tail == Capacity) return false;
        buffer_[head & (Capacity - 1)] = value;
        head_.store(head + 1, std::memory_order_release);
        return true;
    }

    bool try_pop(int& result) {
        const auto tail = tail_.load(std::memory_order_relaxed);
        const auto head = head_.load(std::memory_order_acquire);
        if (tail == head) return false;
        result = buffer_[tail & (Capacity - 1)];
        tail_.store(tail + 1, std::memory_order_release);
        return true;
    }
};

int main() {
    constexpr int Count = 100000;
    SpscQueue<1024> queue;
    long long result = 0;

    std::jthread producer([&] {
        for (int x = 1; x <= Count; ++x) {
            while (!queue.try_push(x)) std::this_thread::yield();
        }
    });
    std::jthread consumer([&] {
        int item = 0;
        for (int n = 0; n < Count; ++n) {
            while (!queue.try_pop(item)) std::this_thread::yield();
            result += item;
        }
    });
    producer.join();
    consumer.join();
    assert(result == 1LL * Count * (Count + 1) / 2);
}
\end{lstlisting}

\subsection{Dlaczego kod jest poprawny?}

Dla danych w~\code{buffer_} nie stosujemy operacji atomowych.
Nie ma jednak konfliktowego dostępu bez synchronizacji, ponieważ:
\begin{enumerate}
\item przed odczytem komórki konsument otrzymuje informację o~zapisie
      przez \code{head_.load(acquire)};
\item przed ponownym nadpisaniem tej komórki producent otrzymuje
      informację o~jej zwolnieniu przez \code{tail_.load(acquire)};
\item każdy z~liczników ma jednego piszącego, więc nie potrzeba CAS
      do aktualizowania \code{head_} ani \code{tail_}.
\end{enumerate}

Przykład zakłada, że operacji \code{try_push} nie wywołuje kilku
producentów, a~\code{try_pop} kilku konsumentów, nawet jeśli
same metody są dostępne publicznie. Złamane założenie SPSC
powoduje utratę poprawności. Ograniczenia bufora, przepełnienie
liczników i~kwestie typów niebanalnych wymagają osobnej analizy.
Użyliśmy potęgi dwójki, żeby indeksowanie przez maskę pozostało
poprawne także przy modularnym zawijaniu licznika bez znaku.

\nota{Ten przykład jest edukacyjną, \emph{nieblokującą} konstrukcją
w~sensie, że nie przejmuje mutexa, ale aktywne pętle producenta
i~konsumenta mogą czekać dowolnie długo. Metody \code{try_push}
i~\code{try_pop} są krótkie i~nie wykonują pętli CAS. Bez uzasadnienia
własności użytych atomików na danej platformie nie deklarujemy
silniejszej gwarancji wydajnościowej.}

\section{Atomowa publikacja niezmiennych wersji danych}
\label{sec:cpp16-snapshot}

Wiele rzeczywistych aplikacji (konfiguracja serwera, słowniki,
parametry modeli) posiada licznych czytelników i~rzadkich pisarzy.
Dobrą architekturą bywa utworzenie nowej, niezmiennej wersji
struktury i~jednorazowe opublikowanie wskaźnika do niej.
W~C++20 można do tego wykorzystać specjalizację
\code{std::atomic<std::shared_ptr<T>>}.

\begin{lstlisting}[caption={Bezpieczny snapshot konfiguracji (C++20).}]
#include <atomic>
#include <cassert>
#include <memory>
#include <thread>

struct Config {
    int threshold;
    double scale;
};

int main() {
    using Ptr = std::shared_ptr<const Config>;
    std::atomic<Ptr> current{
        std::make_shared<const Config>(Config{10, 1.0})
    };

    std::jthread writer([&] {
        Ptr next = std::make_shared<const Config>(Config{20, 2.0});
        current.store(std::move(next), std::memory_order_release);
    });
    std::jthread reader([&] {
        Ptr snap = current.load(std::memory_order_acquire);
        assert((snap->threshold == 10 && snap->scale == 1.0) ||
               (snap->threshold == 20 && snap->scale == 2.0));
    });
    writer.join();
    reader.join();
}
\end{lstlisting}

Czytelnik otrzymuje \emph{spójny} snapshot obu pól, a~stara wersja
jest automatycznie utrzymywana przy życiu, dopóki istnieją do niej
uchwyty \code{shared_ptr}. Nie oznacza to, że przechowywanie
\code{shared_ptr} w~atomiku jest \emph{lock-free}. Referencyjne
liczenie i~alokacje bywają istotnym kosztem przy bardzo częstych
odczytach. Dla obciążenia z~dużą przewagą czytelników należy
porównać to podejście z~\code{shared_mutex}, kopiowaniem małych
struktur i~technikami RCU.

\subsection{\code{std::atomic_ref}: synchronizacja istniejącego obiektu}

\code{std::atomic_ref<T>} (C++20) pozwala wykonywać operacje atomowe
na istniejącym, odpowiednio wyrównanym obiekcie typu \code{T}.
Przydatne bywa przy interoperacyjności ze strukturami danych,
które nie przechowują \code{std::atomic<T>} jako pola.

\begin{lstlisting}[caption={Operacje atomowe na istniejącej zmiennej (C++20).}]
#include <atomic>
#include <cassert>
#include <thread>

int main() {
    alignas(std::atomic_ref<int>::required_alignment) int counter = 0;
    {
        std::atomic_ref<int> shared_counter(counter);
        auto add = [&] {
            for (int i = 0; i < 10000; ++i)
                shared_counter.fetch_add(1, std::memory_order_relaxed);
        };
        std::jthread a(add), b(add);
        a.join();
        b.join();
    } // koniec zycia atomic_ref
    assert(counter == 20000);
}
\end{lstlisting}

Dopóki istnieje \code{atomic_ref} odnoszący się do obiektu, dostęp
do tego obiektu należy realizować przez instancje \code{atomic_ref}
--- nie tylko w~chwilach równoczesnego dostępu. Dlatego zwykłe
odczytanie \code{counter} wykonujemy \emph{po} zakończeniu czasu
życia referencji atomowej. Sam obiekt musi żyć dostatecznie długo,
a~jego wyrównanie musi spełniać wymagania \code{required_alignment}.

\section{Planowanie zadań: pule wątków, \emph{work stealing} i~\emph{backpressure}}
\label{sec:cpp16-scheduling}

Poza poprawną synchronizacją potrzebujemy rozsądnej strategii
przypisywania prac do wątków. Wątek systemowy ma koszt utworzenia,
przełączenia kontekstu i~utrzymania stosu. Z~tego powodu tysiące
krótkich zadań nie powinny automatycznie oznaczać tysiąca wątków.

\begin{description}[style=nextline]
\item[\emph{Thread pool}]
  Ustalona liczba wątków pobiera zadania z~kolejki. Ogranicza koszt
  częstego tworzenia wątków i~może ograniczać równoległe obciążenie.
\item[\emph{Work stealing}]
  Każdy wątek ma własną kolejkę, a~bezczynny pracownik próbuje
  przejąć część pracy od innych. Dobrze sprawdza się przy nierównych
  czasach trwania podzadań. Wymaga ostrożnego protokołu synchronizacji
  kolejek, szczególnie gdy korzystają z~nich różni producenci.
\item[\emph{Backpressure}]
  Mechanizm spowalniania producenta, gdy konsumenci nie nadążają.
  Bez niego kolejka bez ograniczenia pojemności może zużyć całą pamięć.
\item[\emph{Structured concurrency}]
  Sposób organizacji zadań, w~którym ich czas życia, anulowanie i~błędy
  pozostają związane z~zakresem nadrzędnej operacji. Sam \code{std::async}
  nie implementuje kompletnego frameworka takiego stylu.
\end{description}

\subsection{Pula wątków: bezpieczne zamknięcie i~propagacja wyjątków}

Następny przykład rozwija pulę z~rozdziału~\ref{chap:concurrency}. Używa mutexa
oraz zmiennej warunkowej. Łatwy dowód poprawności i~przewidywalne
zachowanie przy zamykaniu są często ważniejsze niż
niepotwierdzony \emph{lock-free}.

\begin{lstlisting}[caption={Prosta pula wątków odprowadzająca kolejkę przed zamknięciem (C++20).}]
#include <condition_variable>
#include <functional>
#include <future>
#include <iostream>
#include <memory>
#include <mutex>
#include <queue>
#include <stdexcept>
#include <thread>
#include <type_traits>
#include <utility>
#include <vector>

class ThreadPool {
    std::mutex mutex_;
    std::condition_variable cv_;
    std::queue<std::function<void()>> jobs_;
    bool closed_ = false;
    std::vector<std::jthread> workers_;

    void worker_loop() {
        for (;;) {
            std::function<void()> job;
            {
                std::unique_lock lock(mutex_);
                cv_.wait(lock, [&] {
                    return closed_ || !jobs_.empty();
                });
                if (closed_ && jobs_.empty()) return;
                job = std::move(jobs_.front());
                jobs_.pop();
            }
            job();  // packaged_task przechwyci wyjatek
        }
    }

public:
    explicit ThreadPool(unsigned count) {
        if (count == 0) throw std::invalid_argument("count=0");
        try {
            for (unsigned i = 0; i < count; ++i)
                workers_.emplace_back([this] { worker_loop(); });
        } catch (...) {
            close();  // obudz pracownikow; ich destruktory wykonaja join
            throw;
        }
    }
    ThreadPool(const ThreadPool&) = delete;
    ThreadPool& operator=(const ThreadPool&) = delete;
    ~ThreadPool() { close(); }

    void close() {
        {
            std::lock_guard lock(mutex_);
            closed_ = true;
        }
        cv_.notify_all();
    }

    template<class F>
    auto submit(F&& callable)
        -> std::future<std::invoke_result_t<std::decay_t<F>&>> {
        using R = std::invoke_result_t<std::decay_t<F>&>;
        auto task = std::make_shared<std::packaged_task<R()>>(
            std::forward<F>(callable)
        );
        auto result = task->get_future();
        {
            std::lock_guard lock(mutex_);
            if (closed_) throw std::runtime_error("pool closed");
            jobs_.emplace([task] { (*task)(); });
        }
        cv_.notify_one();
        return result;
    }
};

int main() {
    ThreadPool pool(2);
    auto a = pool.submit([] { return 40; });
    auto b = pool.submit([] { return 2; });
    std::cout << a.get() + b.get() << '\n'; // 42
    pool.close(); // od tej chwili bez nowych zadan
}
\end{lstlisting}

Ważne własności pokazanej implementacji:
\begin{itemize}
\item \code{close()} jest bezpieczne wobec równoległego \code{submit()}:
      oba działania są serializowane przez ten sam mutex;
\item po \code{close()} pracownicy \emph{kończą kolejkę}, a~dopiero
      następnie opuszczają pętlę; destruktory \code{jthread} dołączają
      pracowników, zanim zniszczone zostaną mutex i~kolejka;
\item \code{packaged_task} zapisuje wartość lub wyjątek w~stanie
      asynchronicznym, więc zwykłe wyjątki z~zadań nie wychodzą
      bezpośrednio przez \code{worker_loop};
\item ta minimalna pula nie ma ograniczenia długości kolejki,
      priorytetów, terminów ani anulowania uruchomionych prac.
\end{itemize}

\nota{Nie wolno niszczyć takiej puli z~zadania wykonywanego w~tej
samej puli: dołączenie własnego wątku prowadzi do błędu.
Podobnie zadanie czekające na inne zadanie z~tej samej puli może
zakleszczyć program, jeśli wszyscy pracownicy pozostają zajęci
oczekiwaniem. Produkcyjny executor wymaga rozwiązań dla takich sytuacji.}

\subsection{Kiedy potrzebny jest \emph{work stealing}?}

Załóżmy, że przetwarzamy drzewo lub dzielimy dużą macierz na bloki.
Część bloków kończy się po mikrosekundach, inne zajmują milisekundy.
Statyczny przydział zadań może zostawić jeden wątek przeciążony,
gdy inne nie mają już pracy. Przechwytywanie pracy pomiędzy lokalnymi
kolejkami poprawia \emph{równoważenie obciążenia}.

Typowy projekt stosuje:
\begin{enumerate}
\item właścicielski dostęp pracownika do jednego końca jego kolejki;
\item kradzież elementów przez inne wątki z~drugiego końca;
\item atomowe indeksy i~protokół obsługujący konflikt właściciela
      ze złodziejem o~ostatni element;
\item kontrolę czasu życia tablicy zadań oraz jej ewentualnej realokacji.
\end{enumerate}

Takie struktury są znacznie bardziej subtelne niż \code{std::deque}
z~mutexem. W~praktyce warto użyć sprawdzonego executora/biblioteki,
a~własną implementację traktować jako osobny, formalnie weryfikowany projekt.

\section{Anulowanie, sygnały stopu i~punkty oczekiwania}
\label{sec:cpp16-cancel}

Kooperacyjne zatrzymywanie w~C++20 jest oparte na
\code{std::stop_source}, \code{std::stop_token} i~\code{std::stop_callback}.
Żądanie zatrzymania nie zabija wątku; umożliwia mu dobrowolne
zakończenie po zweryfikowaniu stanu i~odpowiednim uporządkowaniu zasobów.

\begin{lstlisting}[caption={Wybudzanie i anulowanie czekania (C++20).}]
#include <condition_variable>
#include <iostream>
#include <mutex>
#include <thread>

int main() {
    std::mutex mutex;
    std::condition_variable_any cv;
    bool ready = false;

    std::jthread worker([&](std::stop_token stop) {
        std::unique_lock lock(mutex);
        bool completed = cv.wait(lock, stop, [&] { return ready; });
        std::cout << (completed ? "ready\n" : "cancelled\n");
    });

    worker.request_stop();
    worker.join();
}
\end{lstlisting}

W~tym przeciążeniu \code{wait} zwraca wartość predykatu po zakończeniu
czekania. Jeśli żądanie zatrzymania nastąpi przed spełnieniem
predykatu, przykład wypisze \code{cancelled}. Jeśli stan zmieniłby
się w~międzyczasie, wynik mógłby być inny. Zwykła
\code{std::condition_variable} nie posiada w~C++20 odpowiednika
tego przeciążenia z~\code{stop_token}.

\subsection{Anulowanie i~bezpieczeństwo wyjątkowe}

W~dobrym API warto określić, co dzieje się z~wynikiem po anulowaniu:
\begin{itemize}
\item czy zwracany jest częściowy wynik;
\item czy do \code{future} trafia własny wyjątek oznaczający anulowanie;
\item czy potrzebne jest odwrócenie skutków wcześniej wykonanych operacji;
\item jak zapewnić, że wszystkie zależne zadania zostaną zakończone
      przed zniszczeniem współdzielonych danych.
\end{itemize}

Żadne z~tych pytań nie ma jednej odpowiedzi dla wszystkich bibliotek.
Sama metoda \code{future.wait_for()} nie anuluje zadań, a~korutyna
z~rozdziału~\ref{chap:cpp15-coroutines} wymaga mechanizmu propagacji żądania anulowania
przez awaitery i~executor.

\section{Fałszywe współdzielenie (\emph{false sharing}) i~cache coherence}
\label{sec:cpp16-false-sharing}

Dwa niezależne atomiki mogą być zapisane blisko siebie w~pamięci,
tak że mieszczą się w~tej samej jednostce przenoszenia danych
pomiędzy pamięciami podręcznymi. Wtedy zapis przez różne rdzenie
może wymagać kosztownych transferów linii cache, choć z~punktu
widzenia C++ nie ma wyścigu danych.

\begin{lstlisting}[caption={Rozdzielenie intensywnie zapisywanych liczników (C++20; stała wyrównania od C++17).}]
#include <atomic>
#include <cstdint>
#include <iostream>
#include <new>
#include <thread>

struct alignas(std::hardware_destructive_interference_size) Counter {
    std::atomic<std::uint64_t> value{0};
};

int main() {
    Counter a, b;
    auto increment = [](Counter& c) {
        for (int i = 0; i < 100000; ++i)
            c.value.fetch_add(1, std::memory_order_relaxed);
    };
    std::jthread first([&] { increment(a); });
    std::jthread second([&] { increment(b); });
    first.join();
    second.join();
    std::cout << a.value << ' ' << b.value << '\n';
}
\end{lstlisting}

\code{std::hardware_destructive_interference_size} jest wskazówką
implementacji dotyczącą rozdzielania obiektów często modyfikowanych
przez różne wątki; w~praktyce jej obsługa zależy od kompilatora
i~biblioteki standardowej. \code{alignas} może zwiększać zużycie
pamięci, zatem nie należy bez pomiarów wyrównywać w~ten sposób
każdego pola. Rozdział~\ref{chap:cpp-memory-performance} omawia szerzej cache locality i~organizację
danych AoS/SoA.

\subsection{Co mierzyć przy ocenie skalowalności?}

Pojedynczy czas wykonania jest mało informacyjny. Przy porównywaniu
wersji jednowątkowej i~wielowątkowej warto mierzyć:
\begin{itemize}
\item przepustowość (operacje/sekundę), a~nie tylko opóźnienie
      pojedynczego zadania;
\item \emph{speedup} \(S(p)=T(1)/T(p)\), dla \(p\) wątków;
\item rozkład opóźnień, np. medianę i~wysokie percentyle;
\item liczbę alokacji, przełączeń kontekstu, braków w~cache i~konfliktów CAS;
\item oddzielnie okres rozgrzewki oraz właściwy okres pomiarowy.
\end{itemize}

Według prawa Amdahla, jeśli część \(s\) programu musi pozostać
sekwencyjna, to w~idealizacji
\[
S(p)\leq\frac{1}{s+\frac{1-s}{p}}.
\]
Wzór ten nie obejmuje kosztów synchronizacji, a~więc rzeczywiste
przyspieszenie może być jeszcze niższe. Jednocześnie program
z~dominującym kosztem dostępu do pamięci może nie skalować się
nawet przy bardzo małej części sekwencyjnej.

\section{Linearizowalność i~niezmienniki współdzielonych struktur}
\label{sec:cpp16-linearizability}

Dla współbieżnych kolekcji często potrzebujemy mocniejszej własności
niż brak wyścigów danych: \emph{linearizowalności}.
Operacja jest linearizowalna, jeśli można przypisać jej jeden
umowny moment pomiędzy wywołaniem i~zakończeniem, w~którym
wywołanie ma taki sam skutek jak operacja wykonana sekwencyjnie.

Dla prostej atomowej inkrementacji punktem linearyzacji jest
sam RMW. Dla kolejki algorytm jest bardziej subtelny:
\begin{itemize}
\item producent najpierw zapisuje komórkę, a~następnie publikuje
      nowy \code{head};
\item konsument najpierw odczytuje i~przetwarza komórkę, a~następnie
      publikuje nowy \code{tail};
\item sama komórka nie musi być atomowa, o~ile zachowana jest
      właściwa synchronizacja i~protokół własności miejsca w~buforze.
\end{itemize}

W~algorytmach złożonych, np. z~dwoma atomikami \code{size} i~\code{head},
każde z~pól może być wolne od wyścigu danych, ale ich jednoczesny
odczyt nie musi stanowić spójnego snapshotu. Osobny \emph{niezmiennik}
struktury wymaga własnego uzasadnienia.

\subsection{Dlaczego prosty \emph{seqlock} nie wystarczy dla zwykłych danych?}

Popularny pomysł polega na tym, że pisarz zwiększa licznik wersji
przed i~po aktualizacji obiektu, a~czytelnik ponawia odczyt, jeżeli
wersja zmieniła się podczas odczytu. Jeżeli pola obiektu są zwykłymi
zmiennymi C++, czytelnik nadal może czytać je jednocześnie z~zapisem,
co oznacza \emph{data race}. Ponowienie odczytu \emph{po fakcie}
nie naprawia niezdefiniowanego zachowania, które już nastąpiło.

Poprawny wariant dla przenośnego C++ musi zastosować odpowiednią
synchronizację lub atomowe pola danych, a~nie tylko atomowy licznik
wersji. Alternatywą jest pokazana wcześniej publikacja niezmiennych
snapshotów.

\section{Współbieżność a~parallel STL, Ranges i~OpenMP}
\label{sec:cpp16-interop}

Nie każda równoległa praca wymaga własnej kolejki czy atomowych
liczników. Gdy zadanie polega na niezależnym przetwarzaniu elementów,
lepszą abstrakcją może być algorytm równoległy lub dyrektywa OpenMP.

\subsection{Polityki wykonania standardowych algorytmów}

Nagłówek \code{<execution>} udostępnia od C++17 między innymi:
\begin{description}[style=nextline]
\item[\code{std::execution::seq}] Wykonanie sekwencyjne.
\item[\code{std::execution::par}] Implementacja może wykonywać pracę
      równolegle; nie daje gwarancji, że rzeczywiście użyje wielu wątków.
\item[\code{std::execution::par_unseq}] Zezwala dodatkowo na
      niesekwencyjne wykonywanie, w~tym wektoryzację; nakłada ważne
      ograniczenia na wykonywane funkcje i~synchronizację.
\end{description}

\begin{lstlisting}[caption={Równoległa transformacja bez współdzielonych zapisów (C++17).}]
#include <algorithm>
#include <execution>
#include <numeric>
#include <vector>

int main() {
    std::vector<double> a(10000), b(a.size());
    std::iota(a.begin(), a.end(), 0.0);
    std::transform(std::execution::par,
                   a.begin(), a.end(), b.begin(),
                   [](double x) { return x * x; });
}
\end{lstlisting}

Kod nie zakłada kolejności przetwarzania elementów. Wyjściowy zakres
\code{b} nie nakłada się z~\code{a}, a~każde wywołanie funkcji
zapisuje inny element. Rzeczywiste zrównoleglenie zależy od
implementacji biblioteki standardowej i~konfiguracji środowiska;
niektóre kompilacje wymagają dodatkowej biblioteki wykonawczej.

\subsection{C++ i~OpenMP: dwa poziomy abstrakcji}

OpenMP, omawiany szczegółowo w~rozdziale~\ref{chap:concurrency}, umożliwia między innymi
równoległe pętle, redukcje i~zależności pomiędzy zadaniami.
W~obliczeniach numerycznych typowe rozwiązanie wygląda tak:

\begin{lstlisting}[caption={Wektorowa redukcja z OpenMP; wymagane -fopenmp.}]
#include <iostream>
#include <vector>

int main() {
    const std::vector<double> x(100000, 1.0);
    double sum = 0.0;

    #pragma omp parallel for reduction(+:sum) schedule(static)
    for (int i = 0; i < static_cast<int>(x.size()); ++i) {
        sum += x[static_cast<std::size_t>(i)];
    }
    std::cout << sum << '\n';
}
\end{lstlisting}

Kodowanie redukcji przez \code{reduction} pozwala uniknąć współdzielenia
jednego licznika zapisywanego przez wszystkie wątki. Kolejność
sumowania liczb zmiennoprzecinkowych może być odmienna niż przy
obliczeniu sekwencyjnym, a~zatem wynik może różnić się przez
zaokrąglenia. Liczbę wątków można ustawić m.in. przez
\code{OMP_NUM_THREADS}; \code{OMP_SCHEDULE} wpływa na harmonogram
tylko przy \code{schedule(runtime)}. Powyższy przykład jawnie
wybiera \code{schedule(static)}.

\nota{Mieszanie pul wątków, OpenMP, równoległych algorytmów STL
oraz zadań asynchronicznych może prowadzić do \emph{oversubscription}:
liczba aktywnych wątków znacząco przewyższa dostępne zasoby CPU.
Przy nakładaniu kilku warstw równoległości trzeba świadomie
kontrolować budżet wątków i~mierzyć efekt końcowy.}

\section{Diagnostyka: kompilatory, sanitizatory, testy przeciążeniowe}
\label{sec:cpp16-testing}

\subsection{Kompilacja i~uruchamianie przykładów}

Dla przykładów wykorzystujących \code{std::jthread}, \code{atomic::wait}
i~\code{atomic<shared_ptr<T>>} wybieramy tryb C++20:

\begin{lstlisting}[language=bash,caption={Kompilacja wielowątkowego programu C++20.}]
g++ -std=c++20 -O2 -Wall -Wextra -Wpedantic -pthread file.cpp -o app
./app
\end{lstlisting}

Dla przykładów OpenMP w~GCC lub odpowiednio skonfigurowanym Clangu:
\begin{lstlisting}[language=bash,caption={Włączenie OpenMP.}]
g++ -std=c++20 -O2 -Wall -Wextra -fopenmp file.cpp -o app
./app
\end{lstlisting}

\subsection{ThreadSanitizer i~inne narzędzia}

\emph{ThreadSanitizer} (TSan) potrafi wykrywać wiele wyścigów danych
i~wybrane błędy synchronizacji w~uruchamianym programie:

\begin{lstlisting}[language=bash,caption={Budowanie i uruchamianie z ThreadSanitizer.}]
g++ -std=c++20 -O1 -g -pthread -fsanitize=thread file.cpp -o app_tsan
./app_tsan
\end{lstlisting}

Nie wszystkie środowiska poprawnie wspierają TSan; szczególnie
przy nietypowych bibliotekach wykonawczych, OpenMP i~własnych
mechanizmach synchronizacji należy weryfikować zakres instrumentacji.
AddressSanitizer i~UndefinedBehaviorSanitizer mogą dodatkowo wykrywać
część błędów czasu życia oraz naruszeń zasad języka, lecz nie są
zamiennikami TSan ani formalnego dowodu synchronizacji.

\subsection{Jak testować algorytm współbieżny?}

Dla poważniejszej struktury danych zaleca się następującą procedurę:
\begin{enumerate}
\item \textbf{Zdefiniuj kontrakt}: ilu producentów i~konsumentów,
      dozwolone operacje oraz zasady zamknięcia i~niszczenia struktury.
\item \textbf{Zidentyfikuj własność danych}: kto może zapisywać dane
      w~każdej chwili; kiedy przechodzi własność konkretnej komórki.
\item \textbf{Narysuj relacje}: wskaż atomik, który publikuje zapis
      i~odczyt, który tę publikację odbiera.
\item \textbf{Wskaż punkt linearyzacji}: jeśli API ma być linearizowalne,
      określ, kiedy operacja staje się widoczna logicznie.
\item \textbf{Przeanalizuj zakończenie}: co dzieje się przy anulowaniu,
      wyjątku, opóźnieniu konsumenta i~niszczeniu obiektu.
\item \textbf{Testuj wielokrotnie}: zmieniaj liczbę wątków, rozmiary
      danych i~częstotliwość przełączeń; uruchamiaj TSan tam,
      gdzie jest dostępny.
\item \textbf{Benchmarkuj osobno}: włącz optymalizacje i~mierz
      przy realistycznym obciążeniu, ale nie myl szybkiego testu
      z~dowodem poprawności.
\end{enumerate}

\section{Typowe błędy i~zalecenia projektowe}
\label{sec:cpp16-guidelines}

\begin{center}
\small
\begin{tabularx}{\textwidth}{@{}lX@{}}
\toprule
\textbf{Błąd} & \textbf{Dlaczego jest groźny / co wybrać} \\
\midrule
\code{volatile} jako synchronizacja & Nie eliminuje \emph{data race};
użyj mutexa lub atomików z~właściwym porządkiem. \\
\code{relaxed} do publikacji zwykłych danych & Nie wytwarza
relacji synchronizacji; zazwyczaj potrzebne \code{release/acquire}. \\
Bezrefleksyjne \code{seq_cst} & Zwykle poprawniejsze niż źle
dobrane słabsze porządki, ale nie zastąpi logicznego kontraktu obiektu. \\
Kasowanie węzła po udanym CAS & Inny wątek może nadal używać
wskaźnika; zaprojektuj odzyskiwanie pamięci. \\
Dwa atomiki jako ,,jedna transakcja'' & Brak automatycznego
atomowego snapshotu pary; użyj blokady lub niezmiennej wersji. \\
\emph{Spinlock} dla długiej pracy & Prowadzi do rywalizacji
o~CPU; zwykle użyj \code{std::mutex}. \\
Nieograniczona kolejka zadań & Może zużyć pamięć; dodaj limit
i~mechanizm \emph{backpressure}. \\
Nadmierna liczba wątków & Konkurencja o~rdzenie i~cache;
kontroluj równoległość i~profiluj. \\
Mieszanie OpenMP i~własnej puli bez planu & Może prowadzić
do \emph{oversubscription} i~spadku wydajności. \\
\bottomrule
\end{tabularx}
\end{center}

\subsection{Praktyczne reguły wyboru}

\begin{itemize}
\item \textbf{Mutex} wybieraj jako punkt wyjścia, jeśli współdzielisz
      kilka powiązanych pól i~potrzebujesz łatwo sprawdzalnych niezmienników.
\item \textbf{Atomik \code{relaxed}} jest właściwy dla prostego licznika
      niezależnego od publikacji innych danych.
\item \textbf{\code{release/acquire}} stosuj, gdy istnieje wyraźny
      punkt przekazania danych lub własności pomiędzy wątkami.
\item \textbf{Lock-free} rozważaj dopiero wtedy, gdy pomiary pokażą
      realny koszt blokad, a~projekt ma przejrzysty dowód poprawności
      oraz plan odzyskiwania pamięci.
\item \textbf{OpenMP / parallel STL} wybieraj do operacji na niezależnych
      elementach, redukcji i~równoległości danych.
\item \textbf{Thread pool} wybieraj do dużej liczby krótkich zadań,
      gdy potrzebujesz ograniczać liczbę wątków, priorytety lub
      przepustowość; osobno zdefiniuj zasady zamykania.
\end{itemize}

\section*{Podsumowanie rozdziału}
\addcontentsline{toc}{section}{Podsumowanie rozdziału}

Zaawansowana współbieżność nie polega na mechanicznym zastępowaniu
mutexów przez \code{std::atomic}. Najpierw trzeba określić
\emph{własność danych i~kontrakt algorytmu}, następnie uzasadnić
relacje \emph{happens-before} i~wyznaczyć punkty widoczności
poszczególnych operacji. Dopiero wtedy warto wybierać porządki
pamięciowe, typy kolejek, planistę i~sposób odzyskiwania pamięci.

W~szczególności brak wyścigu danych nie oznacza jeszcze poprawności
logicznej, a~nieblokujący dostęp do atomika nie gwarantuje
nieblokującego ani szybkiego zachowania całego programu. Bezpieczna
współbieżność to połączenie \emph{modelu formalnego}, \emph{kontroli
czasu życia} i~\emph{rzetelnego pomiaru}.

\section*{Literatura i dokumentacja}
\addcontentsline{toc}{section}{Literatura i dokumentacja}
\begin{itemize}
\item Anthony Williams, \emph{C++ Concurrency in Action}, wyd.~2,
      Manning, 2019.
\item ISO/IEC 14882:2020, rozdziały o~modelu pamięci, wątkach,
      atomikach i~synchronizacji.
\item \url{https://en.cppreference.com/w/cpp/atomic/memory_order}
      --- porządki pamięci i~relacje synchronizacji.
\item \url{https://en.cppreference.com/w/cpp/atomic/atomic}
      --- typy atomowe i~operacje RMW.
\item \url{https://en.cppreference.com/w/cpp/atomic/atomic_ref}
      --- atomowe odwołania do istniejących obiektów.
\item \url{https://en.cppreference.com/w/cpp/thread/jthread}
      --- \code{jthread}, czas życia i~kooperacyjne zatrzymywanie.
\item OpenMP Architecture Review Board,
      \emph{OpenMP Application Programming Interface};
      \url{https://www.openmp.org/specifications/}.
\end{itemize}

% Kazdy listing opisany jako kompletny program jest samodzielny.

\chapter{Zaawansowana współbieżność, model pamięci i algorytmy bez blokad}
\label{chap:cpp16-advanced-concurrency}

\section{Cel rozdziału i powiązanie z podstawami}

W rozdziale~\ref{chap:concurrency} poznaliśmy wątki, mutexy, zmienne warunkowe, operacje
atomowe, \code{std::async}, pule wątków i~OpenMP. Sama znajomość tych
interfejsów nie wystarcza jednak do samodzielnego projektowania
algorytmów współbieżnych. Trzeba odpowiedzieć na trudniejsze pytania:
\begin{itemize}
\item Kiedy zapis wykonany przez jeden wątek \emph{musi} być widoczny dla drugiego?
\item Które wyniki są dozwolone przez model pamięci, nawet gdy wydają się
      sprzeczne z intuicją opartą na sekwencyjnym wykonywaniu?
\item Czym różni się atomowość pojedynczej operacji od niepodzielności
      całego algorytmu?
\item Jak zaprojektować \emph{postęp} algorytmu i~bezpiecznie zwalniać pamięć,
      gdy inne wątki mogą jeszcze posiadać wskaźniki do obiektów?
\item Kiedy rozwiązanie bez mutexa rzeczywiście poprawia wydajność?
\end{itemize}

Skupimy się na \emph{uzasadnianiu poprawności}, a~nie tylko na składni.
W~szczególności przedstawimy proste testy modelu pamięci (\emph{litmus tests}),
algorytm kolejki SPSC, publikację niezmiennych struktur danych, problem ABA,
techniki odzyskiwania pamięci oraz praktyczne zasady pomiaru i~diagnostyki.
Nie zakładamy, że każdy \code{std::atomic<T>} jest realizowany bez blokad.

\nota{Wszystkie gwarancje sformułowane w~rozdziale dotyczą abstrakcyjnego
modelu języka C++. Rzeczywisty procesor może wykonywać instrukcje
w~innej kolejności i~stosować bufory zapisów, ale program musi być
poprawny na gruncie standardu, a~nie tylko na jednym komputerze.}

\section{Formalny porządek zdarzeń w~wielowątkowym programie}
\label{sec:cpp16-formal}

\subsection{Zdarzenia i~relacje \emph{sequenced-before}, \emph{synchronizes-with}}

Rozpatrzmy pewne operacje na pamięci jako \emph{zdarzenia}. Wśród relacji
używanych do opisu ich uporządkowania najważniejsze są:
\begin{description}[style=nextline]
\item[\emph{Sequenced-before} (\(\mathrm{sb}\))]
  Porządek określony regułami języka pomiędzy wybranymi operacjami
  w~obrębie jednego wątku. Nie należy go mylić z~dowolną kolejnością
  pojawiania się instrukcji w~kodzie źródłowym: język zawiera także
  wyrażenia z~nieokreśloną kolejnością ewaluacji.
\item[\emph{Synchronizes-with} (\(\mathrm{sw}\))]
  Synchronizacja pomiędzy odpowiednimi operacjami różnych wątków:
  na przykład zapis \code{release} i~odczyt \code{acquire}, który czyta
  wartość z~odpowiedniego zapisu lub jego sekwencji \emph{release}.
  Podobną synchronizację zapewniają właściwe operacje na mutexach.
\item[\emph{Happens-before} (\(\mathrm{hb}\))]
  Relacja porządkująca zdarzenia, budowana w~typowych przykładach
  przez \(\mathrm{sb}\), \(\mathrm{sw}\) i~ich przechodnie złożenia.
  Jeśli \(A\;\mathrm{hb}\;B\), skutki \(A\) są uporządkowane przed \(B\)
  w~sensie formalnego modelu; nie jest to stwierdzenie o~czasie zegarowym.
\end{description}

Dla przejrzystości będziemy używać następującego schematu:
\[
\underbrace{A\;\xrightarrow{\mathrm{sb}}\;R}_{\text{producent}}
\quad\xrightarrow{\mathrm{sw}}\quad
\underbrace{Q\;\xrightarrow{\mathrm{sb}}\;B}_{\text{konsument}}
\quad\Longrightarrow\quad A\;\xrightarrow{\mathrm{hb}}\;B.
\]
Operacja \(R\) jest publikacją, \(Q\) odpowiadającym jej odbiorem;
\(A\) i~\(B\) mogą dotyczyć zwykłej, nieatomowej zmiennej.

\subsection{Atomowość nie zastępuje synchronizacji}

W~przeciwieństwie do zwykłego odczytu lub zapisu operacje na
\code{std::atomic<T>} nie powodują wyścigu danych pomiędzy sobą.
Nie oznacza to jednak, że każdy odczyt atomowy widzi ostatni zapis
według zegara ściennego. Taki globalny \emph{ostatni zapis} zwykle
w~ogóle nie istnieje jako zdefiniowane pojęcie modelu.

Z~kolei jeśli dwa wątki odczytują i~modyfikują to samo miejsce pamięci
nieatomowo, a~konfliktowe operacje nie są uporządkowane relacją
\emph{happens-before}, program ma \emph{data race} i~niezdefiniowane zachowanie.
Samo oznaczenie obiektu przez \code{volatile} nie rozwiązuje tego problemu.

\subsection{Porządek modyfikacji (\emph{modification order}) i~spójność}

Dla \emph{każdego pojedynczego} obiektu atomowego istnieje całkowity
porządek jego modyfikacji (\emph{modification order}). Nie wynika z~tego
jeden wspólny porządek modyfikacji \emph{wszystkich} atomików.
Ważne konsekwencje to między innymi:
\begin{itemize}
\item zapis do \code{x} uporządkowany przed innym zapisem do \code{x}
      musi być odpowiednio umieszczony przed nim w~porządku modyfikacji;
\item jeśli odczyt \code{x} jest uporządkowany przed kolejnym odczytem
      \code{x} w~tym samym wątku, drugi odczyt nie może cofnąć się
      do wcześniejszej modyfikacji w~porządku tego atomika;
\item ograniczenia spójności dla jednego atomika nie zapewniają
      automatycznie spójnego obrazu \emph{dwóch różnych} atomików.
\end{itemize}

\section{Testy modelu pamięci: wyniki, które zaskakują}
\label{sec:cpp16-litmus}

\subsection{\emph{Store buffering}: dwa odczyty zera}

Załóżmy, że \code{x} i~\code{y} początkowo są równe zero, a~wątki
wykonują następujące operacje:
\[
\begin{array}{c|c}
\text{Wątek 1} & \text{Wątek 2}\\ \hline
x\gets 1 & y\gets 1\\
r_1\gets y & r_2\gets x
\end{array}
\]

\begin{lstlisting}[caption={Store buffering: kompletna, wolna od wyścigu danych próba (C++20).}]
#include <atomic>
#include <iostream>
#include <thread>

int main() {
    std::atomic<int> x{0}, y{0};
    int r1 = -1, r2 = -1;

    std::jthread a([&] {
        x.store(1, std::memory_order_relaxed);
        r1 = y.load(std::memory_order_relaxed);
    });
    std::jthread b([&] {
        y.store(1, std::memory_order_relaxed);
        r2 = x.load(std::memory_order_relaxed);
    });
    a.join();
    b.join();
    std::cout << r1 << ' ' << r2 << '\n';
}
\end{lstlisting}

W~tym programie \(r_1=r_2=0\) jest \textbf{dozwolonym} wynikiem
modelu C++ dla \code{relaxed}; nie ma gwarancji, że zostanie
zaobserwowany w~praktycznym uruchomieniu. Jeśli \emph{wszystkie cztery}
operacje atomowe zamienimy na \code{memory_order_seq_cst}, wynik
\((0,0)\) stanie się niemożliwy: wspólny porządek operacji
sekwencyjnie spójnych musiałby przeczyć kolejności w~każdym z~wątków.

\nota{Wynik pojedynczego eksperymentu, nawet powtórzony milion razy,
nie dowodzi niedopuszczalności innego zachowania. Test tego rodzaju
to materiał do rozumowania nad modelem, a~nie pełny dowód poprawności.}

\subsection{\emph{Message passing}: publikacja zwykłych danych}

\begin{lstlisting}[caption={Publikacja danych przy użyciu release/acquire (C++11).}]
#include <atomic>
#include <cassert>
#include <thread>

int main() {
    int payload = 0;
    std::atomic<bool> published{false};

    std::thread producer([&] {
        payload = 42;
        published.store(true, std::memory_order_release);
    });
    std::thread consumer([&] {
        while (!published.load(std::memory_order_acquire)) {
            std::this_thread::yield();
        }
        assert(payload == 42);
    });

    producer.join();
    consumer.join();
}
\end{lstlisting}

Odczyt flagi, który widzi \code{true}, synchronizuje się z~zapisem
\code{release}. Zapis \code{payload = 42} następuje \emph{happens-before}
odczytu \code{payload} przez konsumenta. Flaga \emph{publikuje}
wcześniejszy stan; nie należy po publikacji modyfikować
\code{payload} bez dalszej synchronizacji.

\section{Porządki atomowe i~ich dokładne zastosowania}
\label{sec:cpp16-orders}

\begin{center}
\small
\begin{tabularx}{\textwidth}{@{}lX@{}}
\toprule
\textbf{Porządek} & \textbf{Typowe zastosowanie i~ograniczenie} \\
\midrule
\code{relaxed} & Liczniki, statystyki, pola kontrolne, gdy
sama atomowość wystarcza, a~relacja synchronizacji wynika skądinąd. \\
\code{release} & Zapis lub operacja RMW publikująca wcześniejsze efekty
wątku. \\
\code{acquire} & Odczyt lub operacja RMW odbierająca stan od publikującego wątku. \\
\code{acq_rel} & Operacja odczyt--modyfikacja--zapis wymagająca obu kierunków. \\
\code{seq_cst} & Operacje z~dodatkową gwarancją pojedynczego wspólnego
porządku dla operacji sekwencyjnie spójnych. \\
\bottomrule
\end{tabularx}
\end{center}

\code{memory_order_consume} ma długą i~skomplikowaną historię
standaryzacji. W~praktycznym kodzie C++20 należy preferować
\code{acquire} zamiast opierania przenośnej poprawności na
\code{consume}. Uwaga: porządek \code{release} jest właściwy dla
zapisu, \code{acquire} dla odczytu, a~\code{acq_rel} dla operacji
odczyt--modyfikacja--zapis; nie wszystkie porządki są poprawne dla
każdego przeciążenia funkcji atomowej.

\subsection{Operacje \emph{read--modify--write}}

Metody \code{fetch_add}, \code{fetch_sub}, \code{exchange} i~udany
\code{compare_exchange_*} wykonują operację RMW. Są one niepodzielne
w~odniesieniu do danego atomika, ale nie tworzą automatycznie
transakcji obejmującej inne pola.

\begin{lstlisting}[caption={Przydzielanie kolejnych numerów: wystarcza relaxed (C++11).}]
#include <atomic>
#include <cassert>
#include <thread>

int main() {
    std::atomic<int> next_id{0};
    int first = -1, second = -1;
    std::thread a([&] {
        first = next_id.fetch_add(1, std::memory_order_relaxed);
    });
    std::thread b([&] {
        second = next_id.fetch_add(1, std::memory_order_relaxed);
    });
    a.join();
    b.join();
    assert(first != second);
    assert(first + second == 1);
}
\end{lstlisting}

Uzyskujemy unikatowe wartości, bo RMW mają jednoznaczny porządek
modyfikacji \code{next_id}. Nie uzyskujemy przez to żadnej dodatkowej
synchronizacji niezwiązanych danych.

\subsection{\emph{Release sequence}: łańcuch operacji RMW}

W~C++20 zapis z~porządkiem \code{release} może rozpoczynać sekwencję
\emph{release} składającą się z~następujących po nim operacji RMW
na \emph{tym samym} atomiku, wykonywanych także przez inne wątki.
Odczyt \code{acquire} wartości pochodzącej z~tej sekwencji może
synchronizować się z~początkowym zapisem \code{release}.

\begin{lstlisting}[caption={Przekazanie publikacji przez relaxed RMW (C++20).}]
#include <atomic>
#include <cassert>
#include <thread>

int main() {
    int payload = 7;
    std::atomic<int> state{0};

    std::thread producer([&] {
        payload = 42;
        state.store(1, std::memory_order_release);
    });
    std::thread intermediary([&] {
        int expected = 0;
        while ((expected = state.load(std::memory_order_acquire))
               == 0) {
            std::this_thread::yield();
        }
        state.fetch_add(1, std::memory_order_relaxed);
    });
    std::thread consumer([&] {
        while (state.load(std::memory_order_acquire) != 2) {
            std::this_thread::yield();
        }
        assert(payload == 42);
    });

    producer.join();
    intermediary.join();
    consumer.join();
}
\end{lstlisting}

Wątek pośredni czeka, aż stan osiągnie 1, i~wykonuje RMW, który
zapisuje 2. Odczyt wartości 2 z~\code{acquire} synchronizuje się
z~pierwotnym \code{release} przez sekwencję \emph{release}.
Nie można w~ogólności zastąpić \code{fetch_add} zwykłym
\code{store(2, relaxed)} i~utrzymać tej samej argumentacji.

\section{CAS: porównaj i~zamień, pętle ponawiania i~porządek niepowodzenia}

\begin{center}\small
\begin{tabularx}{\textwidth}{@{}p{0.38\textwidth}X@{}}
\toprule
Operacja atomowa & Działanie \\
\midrule
\code{compare_exchange_strong} & Jednorazowo porównuje i próbuje podmienić wartość; nie zgłasza pozornego niepowodzenia. \\
\code{compare_exchange_weak} & Może zgłosić pozorne niepowodzenie; nadaje się do pętli ponawiania. \\
\code{expected} przy CAS & Przy niepowodzeniu otrzymuje rzeczywistą wartość atomika, którą trzeba uwzględnić przy kolejnej próbie. \\
\code{memory_order_acq_rel} & Przy udanej operacji odczytująco-zapisującej łączy własności acquire i release. \\
\code{memory_order_relaxed} & Zapewnia atomowość bez dodatkowej synchronizacji z innymi danymi. \\
\bottomrule
\end{tabularx}
\end{center}
\label{sec:cpp16-cas}

\code{compare_exchange_weak(expected, desired, success, failure)}
atomowo porównuje bieżącą wartość z~\code{expected}. Gdy porównanie
się powiedzie, zapisuje \code{desired}; w~przeciwnym razie aktualizuje
\code{expected} do wartości zastanej w~atomiku. Wersja \code{weak}
może ponadto nie powieść się pozornie (\emph{spurious failure}).
W~pętli jest zwykle dobrym wyborem.

\begin{lstlisting}[caption={Atomowe maksimum z~pętlą CAS (C++11).}]
#include <atomic>
#include <cassert>

void atomic_max(std::atomic<int>& target, int candidate) {
    int old = target.load(std::memory_order_relaxed);
    while (old < candidate &&
           !target.compare_exchange_weak(
               old, candidate,
               std::memory_order_relaxed,
               std::memory_order_relaxed)) {
        // Przy niepowodzeniu old jest aktualizowane.
    }
}

int main() {
    std::atomic<int> maximum{10};
    atomic_max(maximum, 8);
    atomic_max(maximum, 17);
    assert(maximum.load() == 17);
}
\end{lstlisting}

Parametr \code{failure} opisuje porządek operacji \emph{odczytu}
w~przypadku niepowodzenia. Nie może to być \code{release} ani
\code{acq_rel}; musi też spełniać ograniczenia względem porządku
sukcesu określone standardem. Typowy wzorzec publikacji nowej
wartości to sukces z~\code{acq_rel} i~niepowodzenie z~\code{acquire}
(albo \code{relaxed}, jeśli tak wynika z~dowodu).

\subsection{\code{weak} kontra \code{strong}}

\code{compare_exchange_strong} nie zgłasza pozornego niepowodzenia
przy poprawnie porównanych wartościach. Może być użyteczne dla
pojedynczej próby; \code{weak} bywa wygodniejsze w~pętlach CAS.
Nie ma uniwersalnej gwarancji, że jedna z~tych wersji zawsze działa
szybciej na danej platformie.

\section{Bariery pamięci i~\texttt{std::atomic\_thread\_fence}}
\label{sec:cpp16-fences}

Bariera \code{std::atomic_thread_fence(order)} jest operacją
porządkowania pamięci, która sama nie odczytuje ani nie zapisuje
konkretnego atomika. Synchronizacja z~innym wątkiem wciąż wymaga
odpowiednich operacji atomowych łączących obie strony. Samo
wstawienie \code{fence} obok zwykłej zmiennej nie usuwa wyścigu danych.

\begin{lstlisting}[caption={Publikacja ze zwolnieniową i~przejmującą barierą pamięci (C++11).}]
#include <atomic>
#include <cassert>
#include <thread>

int main() {
    int payload = 0;
    std::atomic<bool> flag{false};

    std::thread producer([&] {
        payload = 42;
        std::atomic_thread_fence(std::memory_order_release);
        flag.store(true, std::memory_order_relaxed);
    });
    std::thread consumer([&] {
        while (!flag.load(std::memory_order_relaxed)) {
            std::this_thread::yield();
        }
        std::atomic_thread_fence(std::memory_order_acquire);
        assert(payload == 42);
    });

    producer.join();
    consumer.join();
}
\end{lstlisting}

To przykład synchronizacji \emph{fence--fence}: odczyt atomika
widzi zapis następujący po barierze \code{release}, a~bariera
\code{acquire} jest wykonywana po tym odczycie. W~nowym kodzie
zwykle prostsze i~czytelniejsze jest zastosowanie pary operacji
\code{release}/\code{acquire} bez jawnych barier. Bariery nie są
substytutem przemyślanego protokołu komunikacji.

\section{Gwarancje postępu: \emph{blocking}, \emph{lock-free} i~\emph{wait-free}}
\label{sec:cpp16-progress}

Wielowątkowy algorytm można analizować pod względem \emph{bezpieczeństwa}
(\emph{safety}: nic niedozwolonego się nie dzieje) i~\emph{żywotności}
(\emph{liveness}: coś pożądanego w~końcu się wydarzy). Są to różne kwestie.

\begin{description}[style=nextline]
\item[\emph{Blocking}]
  Wątek może oczekiwać na odblokowanie mutexa przez konkretny inny wątek.
  Jeśli właściciel przestanie robić postępy, pozostali również mogą utknąć.
\item[\emph{Lock-free}]
  W~standardowym ujęciu postępu algorytmowego, mimo rywalizacji wątków,
  system jako całość może kontynuować wykonywanie operacji:
  niepowodzenia jednej próby CAS zwykle oznaczają, że inny uczestnik
  wykonał postęp. Nie gwarantuje to postępu \emph{każdego} wątku.
\item[\emph{Wait-free}]
  Każda operacja danego wątku ma ograniczoną liczbę własnych kroków
  niezależnie od działań innych wątków (przy przyjętym modelu).
  To silniejszy warunek niż \emph{lock-free}.
\item[\emph{Obstruction-free}]
  Operacja postępuje, jeśli wątek może działać dostatecznie długo
  bez konfliktów ze strony innych wątków. Jest to słabsza gwarancja.
\end{description}

\emph{Lock-free} nie znaczy ,,zawsze szybciej''. Nie oznacza również
braku aktywnego oczekiwania ani braku zagłodzenia konkretnego wątku.
Odróżniajmy też własność implementacji pojedynczego atomika
(\code{is_lock_free()}) od postępu kompletnego algorytmu.

\subsection{Prosta blokada wirująca (\emph{spinlock})}

\begin{lstlisting}[caption={Poprawna blokada wirująca z flagą atomową; ilustracja (C++20).}]
#include <atomic>
#include <cassert>
#include <mutex>
#include <thread>

class SpinLock {
    std::atomic_flag flag = ATOMIC_FLAG_INIT;
public:
    void lock() noexcept {
        while (flag.test_and_set(std::memory_order_acquire)) {
            // Zamiast stale rywalizowac o cache line, czekamy.
            flag.wait(true, std::memory_order_relaxed);
        }
    }
    void unlock() noexcept {
        flag.clear(std::memory_order_release);
        flag.notify_one();
    }
};

int main() {
    SpinLock lock;
    int counter = 0;
    auto work = [&] {
        for (int i = 0; i < 10000; ++i) {
            std::lock_guard<SpinLock> guard(lock);
            ++counter;
        }
    };
    std::jthread a(work), b(work);
    a.join();
    b.join();
    assert(counter == 20000);
}
\end{lstlisting}

\code{test_and_set} ma semantykę RMW; odblokowanie jest operacją
\code{release}. Metoda \code{wait(true)} może blokować, a~zatem
powyższy \code{SpinLock} \emph{nie} jest algorytmem \emph{lock-free}.
Ponadto nie gwarantuje sprawiedliwości. W~realnym oprogramowaniu
\code{std::mutex} zwykle pozostaje lepszym wyborem, zwłaszcza gdy
sekcja krytyczna może być długa lub obciążenie jest wysokie.

\section{Problem ABA: poprawne CAS nie wystarcza}
\label{sec:cpp16-aba}

Rozważmy wierzchołek stosu reprezentowany wskaźnikiem atomowym.
Wątek \(A\) czyta wskaźnik \(p\) i~zamierza wykonać
\code{compare_exchange} z~wartości \(p\) na \code{p->next}.
W~tym czasie wątek \(B\) może:
\begin{enumerate}
\item zdjąć element \(p\) ze stosu;
\item zdjąć jeszcze jeden element;
\item ponownie umieścić \(p\) na wierzchołku stosu.
\end{enumerate}

Wartość wskaźnika uległa zmianie \(A\to B\to A\) (stąd \emph{ABA}),
ale CAS wątku \(A\) może zobaczyć ponownie tę samą wartość \(p\).
W~międzyczasie struktura stosu mogła się całkowicie zmienić.

\begin{center}
\(\boxed{\text{head}=p}\;\longrightarrow\;
  \boxed{\text{head}=q}\;\longrightarrow\;
  \boxed{\text{head}=p}\)
\end{center}

Dwa problemy trzeba odróżnić:
\begin{enumerate}
\item \emph{ABA logiczne}: porównanie tej samej wartości nie wykrywa
      zmian stanu struktury.
\item \emph{Bezpieczeństwo pamięci}: wątek \(A\) może próbować odczytać
      \code{p->next}, gdy inny wątek zniszczył już obiekt \code{*p}.
      Taki odczyt nie staje się poprawny tylko dlatego, że \code{head}
      jest atomowy.
\end{enumerate}

\subsection{Wskaźnik z~licznikiem generacji}

Jedna z~technik polega na porównywaniu pary
\((\text{wskaźnik},\text{licznik wersji})\) zamiast samego wskaźnika.
Wtedy ponowne pojawienie się tego samego adresu nie oznacza równości
całego stanu. Licznik również ma ograniczony zakres i~może się
przepełnić; trzeba uzasadnić brak niebezpiecznego \emph{wrap-around}.

\nota{Nie zakładaj, że \code{std::atomic<SomeStruct>} będzie bez blokad:
implementacja może wykorzystać wewnętrzny mutex lub bibliotekę
atomowych operacji. Możliwość zapisania stanu złożonego w~\code{atomic}
nie rozwiązuje też samodzielnie problemu czasu życia obiektów.}

\subsection{Hazard pointers, epoch-based reclamation i~RCU}

\term{Hazard pointers} pozwalają wątkom publikować informację:
,,zamierzam korzystać z~tego obiektu''. Wątek usuwający obiekt
odkłada go na listę \emph{retired} i~odracza zniszczenie do chwili,
gdy żaden aktywny hazard pointer go nie wskazuje. Poprawna procedura
zwykle musi ponownie zweryfikować odczyt wskaźnika po publikacji
hazardu; samo opublikowanie dowolnego, potencjalnie już nieważnego
adresu nie wystarcza.

\term{Epoch-based reclamation} grupuje operacje w~epoki. Usunięte
obiekty mogą zostać odzyskane dopiero wtedy, gdy wszystkie aktywne
wątki zakończyły okresy korzystania z~danych obejmujące ich usunięcie.
Dłużej zatrzymany wątek może opóźniać zwolnienie pamięci.

\term{Read-copy-update (RCU)} korzysta z~odczytów starych wersji
danych i~publikacji nowych wersji; usuwanie poprzednich odbywa się
po odpowiednim okresie ochronnym (\emph{grace period}). Szczegóły
różnią się pomiędzy implementacjami.

\begin{center}
\small
\begin{tabularx}{\textwidth}{@{}lX@{}}
\toprule
\textbf{Metoda} & \textbf{Kiedy rozważyć} \\
\midrule
Hazard pointers & Wiele odczytów wskaźników, częste usuwanie
węzłów; możliwa kosztowna administracja hazardami. \\
Epoki & Duży ruch, kontrolowane sekcje dostępu; zatrzymany czytelnik
może opóźnić odzyskanie pamięci. \\
RCU & Wzorce ,,wielu czytelników, niewielu pisarzy'', publikacja
niezmiennych wersji danych. \\
\code{shared_ptr} & Gdy łatwość dowodu poprawności jest ważniejsza
od kosztu liczenia referencji; nie obiecuje \emph{lock-free}. \\
\bottomrule
\end{tabularx}
\end{center}

W~C++20 standardowa biblioteka nie udostępnia jeszcze kompletnej,
przenośnej abstrakcji hazard pointerów. Bezpieczna implementacja
wymaga precyzyjnego protokołu pamięciowego, dlatego zamiast podawać
uproszczony, ale błędny stos ze zwalnianiem węzłów, przedstawimy
algorytm bez dynamicznego odzyskiwania pamięci: kolejkę SPSC.

\section{Kolejka SPSC: praktyczny algorytm bez mutexa}
\label{sec:cpp16-spsc}

\emph{Single Producer, Single Consumer} (SPSC) oznacza dokładnie
jednego producenta i~jednego konsumenta. Bufor cykliczny może wtedy
używać dwóch liczników:
\begin{itemize}
\item \code{head} --- numer kolejnego miejsca zapisu, modyfikowany
      wyłącznie przez producenta;
\item \code{tail} --- numer kolejnego miejsca odczytu, modyfikowany
      wyłącznie przez konsumenta.
\end{itemize}

Producent zapisuje komórkę, a~następnie publikuje nowy \code{head}
operacją \code{release}. Konsument odczytuje \code{head} przez
\code{acquire}, a~więc może bezpiecznie zobaczyć zawartość wcześniej
zapisanej komórki. Analogicznie konsument zwalnia komórkę, publikując
nowy \code{tail}, a~producent odbiera tę informację operacją
\code{acquire}. Warunek pełności bufora o~pojemności \(C\) to
\(\mathrm{head}-\mathrm{tail}=C\), a~pustości:
\(\mathrm{head}=\mathrm{tail}\).

\begin{lstlisting}[caption={Kompletna kolejka SPSC dla int, ze statyczną pamięcią (C++20).}]
#include <array>
#include <atomic>
#include <cassert>
#include <cstddef>
#include <thread>

// Tylko jeden producent i dokladnie jeden konsument.
template<std::size_t Capacity>
class SpscQueue {
    static_assert(Capacity > 0);
    static_assert((Capacity & (Capacity - 1)) == 0);
    std::array<int, Capacity> buffer_{};
    std::atomic<std::size_t> head_{0};
    std::atomic<std::size_t> tail_{0};
public:
    bool try_push(int value) {
        const auto head = head_.load(std::memory_order_relaxed);
        const auto tail = tail_.load(std::memory_order_acquire);
        if (head - tail == Capacity) return false;
        buffer_[head & (Capacity - 1)] = value;
        head_.store(head + 1, std::memory_order_release);
        return true;
    }

    bool try_pop(int& result) {
        const auto tail = tail_.load(std::memory_order_relaxed);
        const auto head = head_.load(std::memory_order_acquire);
        if (tail == head) return false;
        result = buffer_[tail & (Capacity - 1)];
        tail_.store(tail + 1, std::memory_order_release);
        return true;
    }
};

int main() {
    constexpr int Count = 100000;
    SpscQueue<1024> queue;
    long long result = 0;

    std::jthread producer([&] {
        for (int x = 1; x <= Count; ++x) {
            while (!queue.try_push(x)) std::this_thread::yield();
        }
    });
    std::jthread consumer([&] {
        int item = 0;
        for (int n = 0; n < Count; ++n) {
            while (!queue.try_pop(item)) std::this_thread::yield();
            result += item;
        }
    });
    producer.join();
    consumer.join();
    assert(result == 1LL * Count * (Count + 1) / 2);
}
\end{lstlisting}

\subsection{Dlaczego kod jest poprawny?}

Dla danych w~\code{buffer_} nie stosujemy operacji atomowych.
Nie ma jednak konfliktowego dostępu bez synchronizacji, ponieważ:
\begin{enumerate}
\item przed odczytem komórki konsument otrzymuje informację o~zapisie
      przez \code{head_.load(acquire)};
\item przed ponownym nadpisaniem tej komórki producent otrzymuje
      informację o~jej zwolnieniu przez \code{tail_.load(acquire)};
\item każdy z~liczników ma jednego piszącego, więc nie potrzeba CAS
      do aktualizowania \code{head_} ani \code{tail_}.
\end{enumerate}

Przykład zakłada, że operacji \code{try_push} nie wywołuje kilku
producentów, a~\code{try_pop} kilku konsumentów, nawet jeśli
same metody są dostępne publicznie. Złamane założenie SPSC
powoduje utratę poprawności. Ograniczenia bufora, przepełnienie
liczników i~kwestie typów niebanalnych wymagają osobnej analizy.
Użyliśmy potęgi dwójki, żeby indeksowanie przez maskę pozostało
poprawne także przy modularnym zawijaniu licznika bez znaku.

\nota{Ten przykład jest edukacyjną, \emph{nieblokującą} konstrukcją
w~sensie, że nie przejmuje mutexa, ale aktywne pętle producenta
i~konsumenta mogą czekać dowolnie długo. Metody \code{try_push}
i~\code{try_pop} są krótkie i~nie wykonują pętli CAS. Bez uzasadnienia
własności użytych atomików na danej platformie nie deklarujemy
silniejszej gwarancji wydajnościowej.}

\section{Atomowa publikacja niezmiennych wersji danych}
\label{sec:cpp16-snapshot}

Wiele rzeczywistych aplikacji (konfiguracja serwera, słowniki,
parametry modeli) posiada licznych czytelników i~rzadkich pisarzy.
Dobrą architekturą bywa utworzenie nowej, niezmiennej wersji
struktury i~jednorazowe opublikowanie wskaźnika do niej.
W~C++20 można do tego wykorzystać specjalizację
\code{std::atomic<std::shared_ptr<T>>}.

\begin{lstlisting}[caption={Bezpieczny snapshot konfiguracji (C++20).}]
#include <atomic>
#include <cassert>
#include <memory>
#include <thread>

struct Config {
    int threshold;
    double scale;
};

int main() {
    using Ptr = std::shared_ptr<const Config>;
    std::atomic<Ptr> current{
        std::make_shared<const Config>(Config{10, 1.0})
    };

    std::jthread writer([&] {
        Ptr next = std::make_shared<const Config>(Config{20, 2.0});
        current.store(std::move(next), std::memory_order_release);
    });
    std::jthread reader([&] {
        Ptr snap = current.load(std::memory_order_acquire);
        assert((snap->threshold == 10 && snap->scale == 1.0) ||
               (snap->threshold == 20 && snap->scale == 2.0));
    });
    writer.join();
    reader.join();
}
\end{lstlisting}

Czytelnik otrzymuje \emph{spójny} snapshot obu pól, a~stara wersja
jest automatycznie utrzymywana przy życiu, dopóki istnieją do niej
uchwyty \code{shared_ptr}. Nie oznacza to, że przechowywanie
\code{shared_ptr} w~atomiku jest \emph{lock-free}. Referencyjne
liczenie i~alokacje bywają istotnym kosztem przy bardzo częstych
odczytach. Dla obciążenia z~dużą przewagą czytelników należy
porównać to podejście z~\code{shared_mutex}, kopiowaniem małych
struktur i~technikami RCU.

\subsection{\code{std::atomic_ref}: synchronizacja istniejącego obiektu}

\code{std::atomic_ref<T>} (C++20) pozwala wykonywać operacje atomowe
na istniejącym, odpowiednio wyrównanym obiekcie typu \code{T}.
Przydatne bywa przy interoperacyjności ze strukturami danych,
które nie przechowują \code{std::atomic<T>} jako pola.

\begin{lstlisting}[caption={Operacje atomowe na istniejącej zmiennej (C++20).}]
#include <atomic>
#include <cassert>
#include <thread>

int main() {
    alignas(std::atomic_ref<int>::required_alignment) int counter = 0;
    {
        std::atomic_ref<int> shared_counter(counter);
        auto add = [&] {
            for (int i = 0; i < 10000; ++i)
                shared_counter.fetch_add(1, std::memory_order_relaxed);
        };
        std::jthread a(add), b(add);
        a.join();
        b.join();
    } // koniec zycia atomic_ref
    assert(counter == 20000);
}
\end{lstlisting}

Dopóki istnieje \code{atomic_ref} odnoszący się do obiektu, dostęp
do tego obiektu należy realizować przez instancje \code{atomic_ref}
--- nie tylko w~chwilach równoczesnego dostępu. Dlatego zwykłe
odczytanie \code{counter} wykonujemy \emph{po} zakończeniu czasu
życia referencji atomowej. Sam obiekt musi żyć dostatecznie długo,
a~jego wyrównanie musi spełniać wymagania \code{required_alignment}.

\section{Planowanie zadań: pule wątków, \emph{work stealing} i~\emph{backpressure}}
\label{sec:cpp16-scheduling}

Poza poprawną synchronizacją potrzebujemy rozsądnej strategii
przypisywania prac do wątków. Wątek systemowy ma koszt utworzenia,
przełączenia kontekstu i~utrzymania stosu. Z~tego powodu tysiące
krótkich zadań nie powinny automatycznie oznaczać tysiąca wątków.

\begin{description}[style=nextline]
\item[\emph{Thread pool}]
  Ustalona liczba wątków pobiera zadania z~kolejki. Ogranicza koszt
  częstego tworzenia wątków i~może ograniczać równoległe obciążenie.
\item[\emph{Work stealing}]
  Każdy wątek ma własną kolejkę, a~bezczynny pracownik próbuje
  przejąć część pracy od innych. Dobrze sprawdza się przy nierównych
  czasach trwania podzadań. Wymaga ostrożnego protokołu synchronizacji
  kolejek, szczególnie gdy korzystają z~nich różni producenci.
\item[\emph{Backpressure}]
  Mechanizm spowalniania producenta, gdy konsumenci nie nadążają.
  Bez niego kolejka bez ograniczenia pojemności może zużyć całą pamięć.
\item[\emph{Structured concurrency}]
  Sposób organizacji zadań, w~którym ich czas życia, anulowanie i~błędy
  pozostają związane z~zakresem nadrzędnej operacji. Sam \code{std::async}
  nie implementuje kompletnego frameworka takiego stylu.
\end{description}

\subsection{Pula wątków: bezpieczne zamknięcie i~propagacja wyjątków}

Następny przykład rozwija pulę z~rozdziału~\ref{chap:concurrency}. Używa mutexa
oraz zmiennej warunkowej. Łatwy dowód poprawności i~przewidywalne
zachowanie przy zamykaniu są często ważniejsze niż
niepotwierdzony \emph{lock-free}.

\begin{lstlisting}[caption={Prosta pula wątków odprowadzająca kolejkę przed zamknięciem (C++20).}]
#include <condition_variable>
#include <functional>
#include <future>
#include <iostream>
#include <memory>
#include <mutex>
#include <queue>
#include <stdexcept>
#include <thread>
#include <type_traits>
#include <utility>
#include <vector>

class ThreadPool {
    std::mutex mutex_;
    std::condition_variable cv_;
    std::queue<std::function<void()>> jobs_;
    bool closed_ = false;
    std::vector<std::jthread> workers_;

    void worker_loop() {
        for (;;) {
            std::function<void()> job;
            {
                std::unique_lock lock(mutex_);
                cv_.wait(lock, [&] {
                    return closed_ || !jobs_.empty();
                });
                if (closed_ && jobs_.empty()) return;
                job = std::move(jobs_.front());
                jobs_.pop();
            }
            job();  // packaged_task przechwyci wyjatek
        }
    }

public:
    explicit ThreadPool(unsigned count) {
        if (count == 0) throw std::invalid_argument("count=0");
        try {
            for (unsigned i = 0; i < count; ++i)
                workers_.emplace_back([this] { worker_loop(); });
        } catch (...) {
            close();  // obudz pracownikow; ich destruktory wykonaja join
            throw;
        }
    }
    ThreadPool(const ThreadPool&) = delete;
    ThreadPool& operator=(const ThreadPool&) = delete;
    ~ThreadPool() { close(); }

    void close() {
        {
            std::lock_guard lock(mutex_);
            closed_ = true;
        }
        cv_.notify_all();
    }

    template<class F>
    auto submit(F&& callable)
        -> std::future<std::invoke_result_t<std::decay_t<F>&>> {
        using R = std::invoke_result_t<std::decay_t<F>&>;
        auto task = std::make_shared<std::packaged_task<R()>>(
            std::forward<F>(callable)
        );
        auto result = task->get_future();
        {
            std::lock_guard lock(mutex_);
            if (closed_) throw std::runtime_error("pool closed");
            jobs_.emplace([task] { (*task)(); });
        }
        cv_.notify_one();
        return result;
    }
};

int main() {
    ThreadPool pool(2);
    auto a = pool.submit([] { return 40; });
    auto b = pool.submit([] { return 2; });
    std::cout << a.get() + b.get() << '\n'; // 42
    pool.close(); // od tej chwili bez nowych zadan
}
\end{lstlisting}

Ważne własności pokazanej implementacji:
\begin{itemize}
\item \code{close()} jest bezpieczne wobec równoległego \code{submit()}:
      oba działania są serializowane przez ten sam mutex;
\item po \code{close()} pracownicy \emph{kończą kolejkę}, a~dopiero
      następnie opuszczają pętlę; destruktory \code{jthread} dołączają
      pracowników, zanim zniszczone zostaną mutex i~kolejka;
\item \code{packaged_task} zapisuje wartość lub wyjątek w~stanie
      asynchronicznym, więc zwykłe wyjątki z~zadań nie wychodzą
      bezpośrednio przez \code{worker_loop};
\item ta minimalna pula nie ma ograniczenia długości kolejki,
      priorytetów, terminów ani anulowania uruchomionych prac.
\end{itemize}

\nota{Nie wolno niszczyć takiej puli z~zadania wykonywanego w~tej
samej puli: dołączenie własnego wątku prowadzi do błędu.
Podobnie zadanie czekające na inne zadanie z~tej samej puli może
zakleszczyć program, jeśli wszyscy pracownicy pozostają zajęci
oczekiwaniem. Produkcyjny executor wymaga rozwiązań dla takich sytuacji.}

\subsection{Kiedy potrzebny jest \emph{work stealing}?}

Załóżmy, że przetwarzamy drzewo lub dzielimy dużą macierz na bloki.
Część bloków kończy się po mikrosekundach, inne zajmują milisekundy.
Statyczny przydział zadań może zostawić jeden wątek przeciążony,
gdy inne nie mają już pracy. Przechwytywanie pracy pomiędzy lokalnymi
kolejkami poprawia \emph{równoważenie obciążenia}.

Typowy projekt stosuje:
\begin{enumerate}
\item właścicielski dostęp pracownika do jednego końca jego kolejki;
\item kradzież elementów przez inne wątki z~drugiego końca;
\item atomowe indeksy i~protokół obsługujący konflikt właściciela
      ze złodziejem o~ostatni element;
\item kontrolę czasu życia tablicy zadań oraz jej ewentualnej realokacji.
\end{enumerate}

Takie struktury są znacznie bardziej subtelne niż \code{std::deque}
z~mutexem. W~praktyce warto użyć sprawdzonego executora/biblioteki,
a~własną implementację traktować jako osobny, formalnie weryfikowany projekt.

\section{Anulowanie, sygnały stopu i~punkty oczekiwania}
\label{sec:cpp16-cancel}

Kooperacyjne zatrzymywanie w~C++20 jest oparte na
\code{std::stop_source}, \code{std::stop_token} i~\code{std::stop_callback}.
Żądanie zatrzymania nie zabija wątku; umożliwia mu dobrowolne
zakończenie po zweryfikowaniu stanu i~odpowiednim uporządkowaniu zasobów.

\begin{lstlisting}[caption={Wybudzanie i anulowanie czekania (C++20).}]
#include <condition_variable>
#include <iostream>
#include <mutex>
#include <thread>

int main() {
    std::mutex mutex;
    std::condition_variable_any cv;
    bool ready = false;

    std::jthread worker([&](std::stop_token stop) {
        std::unique_lock lock(mutex);
        bool completed = cv.wait(lock, stop, [&] { return ready; });
        std::cout << (completed ? "ready\n" : "cancelled\n");
    });

    worker.request_stop();
    worker.join();
}
\end{lstlisting}

W~tym przeciążeniu \code{wait} zwraca wartość predykatu po zakończeniu
czekania. Jeśli żądanie zatrzymania nastąpi przed spełnieniem
predykatu, przykład wypisze \code{cancelled}. Jeśli stan zmieniłby
się w~międzyczasie, wynik mógłby być inny. Zwykła
\code{std::condition_variable} nie posiada w~C++20 odpowiednika
tego przeciążenia z~\code{stop_token}.

\subsection{Anulowanie i~bezpieczeństwo wyjątkowe}

W~dobrym API warto określić, co dzieje się z~wynikiem po anulowaniu:
\begin{itemize}
\item czy zwracany jest częściowy wynik;
\item czy do \code{future} trafia własny wyjątek oznaczający anulowanie;
\item czy potrzebne jest odwrócenie skutków wcześniej wykonanych operacji;
\item jak zapewnić, że wszystkie zależne zadania zostaną zakończone
      przed zniszczeniem współdzielonych danych.
\end{itemize}

Żadne z~tych pytań nie ma jednej odpowiedzi dla wszystkich bibliotek.
Sama metoda \code{future.wait_for()} nie anuluje zadań, a~korutyna
z~rozdziału~\ref{chap:cpp15-coroutines} wymaga mechanizmu propagacji żądania anulowania
przez awaitery i~executor.

\section{Fałszywe współdzielenie (\emph{false sharing}) i~cache coherence}
\label{sec:cpp16-false-sharing}

Dwa niezależne atomiki mogą być zapisane blisko siebie w~pamięci,
tak że mieszczą się w~tej samej jednostce przenoszenia danych
pomiędzy pamięciami podręcznymi. Wtedy zapis przez różne rdzenie
może wymagać kosztownych transferów linii cache, choć z~punktu
widzenia C++ nie ma wyścigu danych.

\begin{lstlisting}[caption={Rozdzielenie intensywnie zapisywanych liczników (C++20; stała wyrównania od C++17).}]
#include <atomic>
#include <cstdint>
#include <iostream>
#include <new>
#include <thread>

struct alignas(std::hardware_destructive_interference_size) Counter {
    std::atomic<std::uint64_t> value{0};
};

int main() {
    Counter a, b;
    auto increment = [](Counter& c) {
        for (int i = 0; i < 100000; ++i)
            c.value.fetch_add(1, std::memory_order_relaxed);
    };
    std::jthread first([&] { increment(a); });
    std::jthread second([&] { increment(b); });
    first.join();
    second.join();
    std::cout << a.value << ' ' << b.value << '\n';
}
\end{lstlisting}

\code{std::hardware_destructive_interference_size} jest wskazówką
implementacji dotyczącą rozdzielania obiektów często modyfikowanych
przez różne wątki; w~praktyce jej obsługa zależy od kompilatora
i~biblioteki standardowej. \code{alignas} może zwiększać zużycie
pamięci, zatem nie należy bez pomiarów wyrównywać w~ten sposób
każdego pola. Rozdział~\ref{chap:cpp-memory-performance} omawia szerzej cache locality i~organizację
danych AoS/SoA.

\subsection{Co mierzyć przy ocenie skalowalności?}

Pojedynczy czas wykonania jest mało informacyjny. Przy porównywaniu
wersji jednowątkowej i~wielowątkowej warto mierzyć:
\begin{itemize}
\item przepustowość (operacje/sekundę), a~nie tylko opóźnienie
      pojedynczego zadania;
\item \emph{speedup} \(S(p)=T(1)/T(p)\), dla \(p\) wątków;
\item rozkład opóźnień, np. medianę i~wysokie percentyle;
\item liczbę alokacji, przełączeń kontekstu, braków w~cache i~konfliktów CAS;
\item oddzielnie okres rozgrzewki oraz właściwy okres pomiarowy.
\end{itemize}

Według prawa Amdahla, jeśli część \(s\) programu musi pozostać
sekwencyjna, to w~idealizacji
\[
S(p)\leq\frac{1}{s+\frac{1-s}{p}}.
\]
Wzór ten nie obejmuje kosztów synchronizacji, a~więc rzeczywiste
przyspieszenie może być jeszcze niższe. Jednocześnie program
z~dominującym kosztem dostępu do pamięci może nie skalować się
nawet przy bardzo małej części sekwencyjnej.

\section{Linearizowalność i~niezmienniki współdzielonych struktur}
\label{sec:cpp16-linearizability}

Dla współbieżnych kolekcji często potrzebujemy mocniejszej własności
niż brak wyścigów danych: \emph{linearizowalności}.
Operacja jest linearizowalna, jeśli można przypisać jej jeden
umowny moment pomiędzy wywołaniem i~zakończeniem, w~którym
wywołanie ma taki sam skutek jak operacja wykonana sekwencyjnie.

Dla prostej atomowej inkrementacji punktem linearyzacji jest
sam RMW. Dla kolejki algorytm jest bardziej subtelny:
\begin{itemize}
\item producent najpierw zapisuje komórkę, a~następnie publikuje
      nowy \code{head};
\item konsument najpierw odczytuje i~przetwarza komórkę, a~następnie
      publikuje nowy \code{tail};
\item sama komórka nie musi być atomowa, o~ile zachowana jest
      właściwa synchronizacja i~protokół własności miejsca w~buforze.
\end{itemize}

W~algorytmach złożonych, np. z~dwoma atomikami \code{size} i~\code{head},
każde z~pól może być wolne od wyścigu danych, ale ich jednoczesny
odczyt nie musi stanowić spójnego snapshotu. Osobny \emph{niezmiennik}
struktury wymaga własnego uzasadnienia.

\subsection{Dlaczego prosty \emph{seqlock} nie wystarczy dla zwykłych danych?}

Popularny pomysł polega na tym, że pisarz zwiększa licznik wersji
przed i~po aktualizacji obiektu, a~czytelnik ponawia odczyt, jeżeli
wersja zmieniła się podczas odczytu. Jeżeli pola obiektu są zwykłymi
zmiennymi C++, czytelnik nadal może czytać je jednocześnie z~zapisem,
co oznacza \emph{data race}. Ponowienie odczytu \emph{po fakcie}
nie naprawia niezdefiniowanego zachowania, które już nastąpiło.

Poprawny wariant dla przenośnego C++ musi zastosować odpowiednią
synchronizację lub atomowe pola danych, a~nie tylko atomowy licznik
wersji. Alternatywą jest pokazana wcześniej publikacja niezmiennych
snapshotów.

\section{Współbieżność a~parallel STL, Ranges i~OpenMP}
\label{sec:cpp16-interop}

Nie każda równoległa praca wymaga własnej kolejki czy atomowych
liczników. Gdy zadanie polega na niezależnym przetwarzaniu elementów,
lepszą abstrakcją może być algorytm równoległy lub dyrektywa OpenMP.

\subsection{Polityki wykonania standardowych algorytmów}

Nagłówek \code{<execution>} udostępnia od C++17 między innymi:
\begin{description}[style=nextline]
\item[\code{std::execution::seq}] Wykonanie sekwencyjne.
\item[\code{std::execution::par}] Implementacja może wykonywać pracę
      równolegle; nie daje gwarancji, że rzeczywiście użyje wielu wątków.
\item[\code{std::execution::par_unseq}] Zezwala dodatkowo na
      niesekwencyjne wykonywanie, w~tym wektoryzację; nakłada ważne
      ograniczenia na wykonywane funkcje i~synchronizację.
\end{description}

\begin{lstlisting}[caption={Równoległa transformacja bez współdzielonych zapisów (C++17).}]
#include <algorithm>
#include <execution>
#include <numeric>
#include <vector>

int main() {
    std::vector<double> a(10000), b(a.size());
    std::iota(a.begin(), a.end(), 0.0);
    std::transform(std::execution::par,
                   a.begin(), a.end(), b.begin(),
                   [](double x) { return x * x; });
}
\end{lstlisting}

Kod nie zakłada kolejności przetwarzania elementów. Wyjściowy zakres
\code{b} nie nakłada się z~\code{a}, a~każde wywołanie funkcji
zapisuje inny element. Rzeczywiste zrównoleglenie zależy od
implementacji biblioteki standardowej i~konfiguracji środowiska;
niektóre kompilacje wymagają dodatkowej biblioteki wykonawczej.

\subsection{C++ i~OpenMP: dwa poziomy abstrakcji}

OpenMP, omawiany szczegółowo w~rozdziale~\ref{chap:concurrency}, umożliwia między innymi
równoległe pętle, redukcje i~zależności pomiędzy zadaniami.
W~obliczeniach numerycznych typowe rozwiązanie wygląda tak:

\begin{lstlisting}[caption={Wektorowa redukcja z OpenMP; wymagane -fopenmp.}]
#include <iostream>
#include <vector>

int main() {
    const std::vector<double> x(100000, 1.0);
    double sum = 0.0;

    #pragma omp parallel for reduction(+:sum) schedule(static)
    for (int i = 0; i < static_cast<int>(x.size()); ++i) {
        sum += x[static_cast<std::size_t>(i)];
    }
    std::cout << sum << '\n';
}
\end{lstlisting}

Kodowanie redukcji przez \code{reduction} pozwala uniknąć współdzielenia
jednego licznika zapisywanego przez wszystkie wątki. Kolejność
sumowania liczb zmiennoprzecinkowych może być odmienna niż przy
obliczeniu sekwencyjnym, a~zatem wynik może różnić się przez
zaokrąglenia. Liczbę wątków można ustawić m.in. przez
\code{OMP_NUM_THREADS}; \code{OMP_SCHEDULE} wpływa na harmonogram
tylko przy \code{schedule(runtime)}. Powyższy przykład jawnie
wybiera \code{schedule(static)}.

\nota{Mieszanie pul wątków, OpenMP, równoległych algorytmów STL
oraz zadań asynchronicznych może prowadzić do \emph{oversubscription}:
liczba aktywnych wątków znacząco przewyższa dostępne zasoby CPU.
Przy nakładaniu kilku warstw równoległości trzeba świadomie
kontrolować budżet wątków i~mierzyć efekt końcowy.}

\section{Diagnostyka: kompilatory, sanitizatory, testy przeciążeniowe}
\label{sec:cpp16-testing}

\subsection{Kompilacja i~uruchamianie przykładów}

Dla przykładów wykorzystujących \code{std::jthread}, \code{atomic::wait}
i~\code{atomic<shared_ptr<T>>} wybieramy tryb C++20:

\begin{lstlisting}[language=bash,caption={Kompilacja wielowątkowego programu C++20.}]
g++ -std=c++20 -O2 -Wall -Wextra -Wpedantic -pthread file.cpp -o app
./app
\end{lstlisting}

Dla przykładów OpenMP w~GCC lub odpowiednio skonfigurowanym Clangu:
\begin{lstlisting}[language=bash,caption={Włączenie OpenMP.}]
g++ -std=c++20 -O2 -Wall -Wextra -fopenmp file.cpp -o app
./app
\end{lstlisting}

\subsection{ThreadSanitizer i~inne narzędzia}

\emph{ThreadSanitizer} (TSan) potrafi wykrywać wiele wyścigów danych
i~wybrane błędy synchronizacji w~uruchamianym programie:

\begin{lstlisting}[language=bash,caption={Budowanie i uruchamianie z ThreadSanitizer.}]
g++ -std=c++20 -O1 -g -pthread -fsanitize=thread file.cpp -o app_tsan
./app_tsan
\end{lstlisting}

Nie wszystkie środowiska poprawnie wspierają TSan; szczególnie
przy nietypowych bibliotekach wykonawczych, OpenMP i~własnych
mechanizmach synchronizacji należy weryfikować zakres instrumentacji.
AddressSanitizer i~UndefinedBehaviorSanitizer mogą dodatkowo wykrywać
część błędów czasu życia oraz naruszeń zasad języka, lecz nie są
zamiennikami TSan ani formalnego dowodu synchronizacji.

\subsection{Jak testować algorytm współbieżny?}

Dla poważniejszej struktury danych zaleca się następującą procedurę:
\begin{enumerate}
\item \textbf{Zdefiniuj kontrakt}: ilu producentów i~konsumentów,
      dozwolone operacje oraz zasady zamknięcia i~niszczenia struktury.
\item \textbf{Zidentyfikuj własność danych}: kto może zapisywać dane
      w~każdej chwili; kiedy przechodzi własność konkretnej komórki.
\item \textbf{Narysuj relacje}: wskaż atomik, który publikuje zapis
      i~odczyt, który tę publikację odbiera.
\item \textbf{Wskaż punkt linearyzacji}: jeśli API ma być linearizowalne,
      określ, kiedy operacja staje się widoczna logicznie.
\item \textbf{Przeanalizuj zakończenie}: co dzieje się przy anulowaniu,
      wyjątku, opóźnieniu konsumenta i~niszczeniu obiektu.
\item \textbf{Testuj wielokrotnie}: zmieniaj liczbę wątków, rozmiary
      danych i~częstotliwość przełączeń; uruchamiaj TSan tam,
      gdzie jest dostępny.
\item \textbf{Benchmarkuj osobno}: włącz optymalizacje i~mierz
      przy realistycznym obciążeniu, ale nie myl szybkiego testu
      z~dowodem poprawności.
\end{enumerate}

\section{Typowe błędy i~zalecenia projektowe}
\label{sec:cpp16-guidelines}

\begin{center}
\small
\begin{tabularx}{\textwidth}{@{}lX@{}}
\toprule
\textbf{Błąd} & \textbf{Dlaczego jest groźny / co wybrać} \\
\midrule
\code{volatile} jako synchronizacja & Nie eliminuje \emph{data race};
użyj mutexa lub atomików z~właściwym porządkiem. \\
\code{relaxed} do publikacji zwykłych danych & Nie wytwarza
relacji synchronizacji; zazwyczaj potrzebne \code{release/acquire}. \\
Bezrefleksyjne \code{seq_cst} & Zwykle poprawniejsze niż źle
dobrane słabsze porządki, ale nie zastąpi logicznego kontraktu obiektu. \\
Kasowanie węzła po udanym CAS & Inny wątek może nadal używać
wskaźnika; zaprojektuj odzyskiwanie pamięci. \\
Dwa atomiki jako ,,jedna transakcja'' & Brak automatycznego
atomowego snapshotu pary; użyj blokady lub niezmiennej wersji. \\
\emph{Spinlock} dla długiej pracy & Prowadzi do rywalizacji
o~CPU; zwykle użyj \code{std::mutex}. \\
Nieograniczona kolejka zadań & Może zużyć pamięć; dodaj limit
i~mechanizm \emph{backpressure}. \\
Nadmierna liczba wątków & Konkurencja o~rdzenie i~cache;
kontroluj równoległość i~profiluj. \\
Mieszanie OpenMP i~własnej puli bez planu & Może prowadzić
do \emph{oversubscription} i~spadku wydajności. \\
\bottomrule
\end{tabularx}
\end{center}

\subsection{Praktyczne reguły wyboru}

\begin{itemize}
\item \textbf{Mutex} wybieraj jako punkt wyjścia, jeśli współdzielisz
      kilka powiązanych pól i~potrzebujesz łatwo sprawdzalnych niezmienników.
\item \textbf{Atomik \code{relaxed}} jest właściwy dla prostego licznika
      niezależnego od publikacji innych danych.
\item \textbf{\code{release/acquire}} stosuj, gdy istnieje wyraźny
      punkt przekazania danych lub własności pomiędzy wątkami.
\item \textbf{Lock-free} rozważaj dopiero wtedy, gdy pomiary pokażą
      realny koszt blokad, a~projekt ma przejrzysty dowód poprawności
      oraz plan odzyskiwania pamięci.
\item \textbf{OpenMP / parallel STL} wybieraj do operacji na niezależnych
      elementach, redukcji i~równoległości danych.
\item \textbf{Thread pool} wybieraj do dużej liczby krótkich zadań,
      gdy potrzebujesz ograniczać liczbę wątków, priorytety lub
      przepustowość; osobno zdefiniuj zasady zamykania.
\end{itemize}

\section*{Podsumowanie rozdziału}
\addcontentsline{toc}{section}{Podsumowanie rozdziału}

Zaawansowana współbieżność nie polega na mechanicznym zastępowaniu
mutexów przez \code{std::atomic}. Najpierw trzeba określić
\emph{własność danych i~kontrakt algorytmu}, następnie uzasadnić
relacje \emph{happens-before} i~wyznaczyć punkty widoczności
poszczególnych operacji. Dopiero wtedy warto wybierać porządki
pamięciowe, typy kolejek, planistę i~sposób odzyskiwania pamięci.

W~szczególności brak wyścigu danych nie oznacza jeszcze poprawności
logicznej, a~nieblokujący dostęp do atomika nie gwarantuje
nieblokującego ani szybkiego zachowania całego programu. Bezpieczna
współbieżność to połączenie \emph{modelu formalnego}, \emph{kontroli
czasu życia} i~\emph{rzetelnego pomiaru}.

\section*{Literatura i dokumentacja}
\addcontentsline{toc}{section}{Literatura i dokumentacja}
\begin{itemize}
\item Anthony Williams, \emph{C++ Concurrency in Action}, wyd.~2,
      Manning, 2019.
\item ISO/IEC 14882:2020, rozdziały o~modelu pamięci, wątkach,
      atomikach i~synchronizacji.
\item \url{https://en.cppreference.com/w/cpp/atomic/memory_order}
      --- porządki pamięci i~relacje synchronizacji.
\item \url{https://en.cppreference.com/w/cpp/atomic/atomic}
      --- typy atomowe i~operacje RMW.
\item \url{https://en.cppreference.com/w/cpp/atomic/atomic_ref}
      --- atomowe odwołania do istniejących obiektów.
\item \url{https://en.cppreference.com/w/cpp/thread/jthread}
      --- \code{jthread}, czas życia i~kooperacyjne zatrzymywanie.
\item OpenMP Architecture Review Board,
      \emph{OpenMP Application Programming Interface};
      \url{https://www.openmp.org/specifications/}.
\end{itemize}

% Kazdy listing opisany jako kompletny program jest samodzielny.

\chapter{Zaawansowana współbieżność, model pamięci i algorytmy bez blokad}
\label{chap:cpp16-advanced-concurrency}

\section{Cel rozdziału i powiązanie z podstawami}

W rozdziale~\ref{chap:concurrency} poznaliśmy wątki, mutexy, zmienne warunkowe, operacje
atomowe, \code{std::async}, pule wątków i~OpenMP. Sama znajomość tych
interfejsów nie wystarcza jednak do samodzielnego projektowania
algorytmów współbieżnych. Trzeba odpowiedzieć na trudniejsze pytania:
\begin{itemize}
\item Kiedy zapis wykonany przez jeden wątek \emph{musi} być widoczny dla drugiego?
\item Które wyniki są dozwolone przez model pamięci, nawet gdy wydają się
      sprzeczne z intuicją opartą na sekwencyjnym wykonywaniu?
\item Czym różni się atomowość pojedynczej operacji od niepodzielności
      całego algorytmu?
\item Jak zaprojektować \emph{postęp} algorytmu i~bezpiecznie zwalniać pamięć,
      gdy inne wątki mogą jeszcze posiadać wskaźniki do obiektów?
\item Kiedy rozwiązanie bez mutexa rzeczywiście poprawia wydajność?
\end{itemize}

Skupimy się na \emph{uzasadnianiu poprawności}, a~nie tylko na składni.
W~szczególności przedstawimy proste testy modelu pamięci (\emph{litmus tests}),
algorytm kolejki SPSC, publikację niezmiennych struktur danych, problem ABA,
techniki odzyskiwania pamięci oraz praktyczne zasady pomiaru i~diagnostyki.
Nie zakładamy, że każdy \code{std::atomic<T>} jest realizowany bez blokad.

\nota{Wszystkie gwarancje sformułowane w~rozdziale dotyczą abstrakcyjnego
modelu języka C++. Rzeczywisty procesor może wykonywać instrukcje
w~innej kolejności i~stosować bufory zapisów, ale program musi być
poprawny na gruncie standardu, a~nie tylko na jednym komputerze.}

\section{Formalny porządek zdarzeń w~wielowątkowym programie}
\label{sec:cpp16-formal}

\subsection{Zdarzenia i~relacje \emph{sequenced-before}, \emph{synchronizes-with}}

Rozpatrzmy pewne operacje na pamięci jako \emph{zdarzenia}. Wśród relacji
używanych do opisu ich uporządkowania najważniejsze są:
\begin{description}[style=nextline]
\item[\emph{Sequenced-before} (\(\mathrm{sb}\))]
  Porządek określony regułami języka pomiędzy wybranymi operacjami
  w~obrębie jednego wątku. Nie należy go mylić z~dowolną kolejnością
  pojawiania się instrukcji w~kodzie źródłowym: język zawiera także
  wyrażenia z~nieokreśloną kolejnością ewaluacji.
\item[\emph{Synchronizes-with} (\(\mathrm{sw}\))]
  Synchronizacja pomiędzy odpowiednimi operacjami różnych wątków:
  na przykład zapis \code{release} i~odczyt \code{acquire}, który czyta
  wartość z~odpowiedniego zapisu lub jego sekwencji \emph{release}.
  Podobną synchronizację zapewniają właściwe operacje na mutexach.
\item[\emph{Happens-before} (\(\mathrm{hb}\))]
  Relacja porządkująca zdarzenia, budowana w~typowych przykładach
  przez \(\mathrm{sb}\), \(\mathrm{sw}\) i~ich przechodnie złożenia.
  Jeśli \(A\;\mathrm{hb}\;B\), skutki \(A\) są uporządkowane przed \(B\)
  w~sensie formalnego modelu; nie jest to stwierdzenie o~czasie zegarowym.
\end{description}

Dla przejrzystości będziemy używać następującego schematu:
\[
\underbrace{A\;\xrightarrow{\mathrm{sb}}\;R}_{\text{producent}}
\quad\xrightarrow{\mathrm{sw}}\quad
\underbrace{Q\;\xrightarrow{\mathrm{sb}}\;B}_{\text{konsument}}
\quad\Longrightarrow\quad A\;\xrightarrow{\mathrm{hb}}\;B.
\]
Operacja \(R\) jest publikacją, \(Q\) odpowiadającym jej odbiorem;
\(A\) i~\(B\) mogą dotyczyć zwykłej, nieatomowej zmiennej.

\subsection{Atomowość nie zastępuje synchronizacji}

W~przeciwieństwie do zwykłego odczytu lub zapisu operacje na
\code{std::atomic<T>} nie powodują wyścigu danych pomiędzy sobą.
Nie oznacza to jednak, że każdy odczyt atomowy widzi ostatni zapis
według zegara ściennego. Taki globalny \emph{ostatni zapis} zwykle
w~ogóle nie istnieje jako zdefiniowane pojęcie modelu.

Z~kolei jeśli dwa wątki odczytują i~modyfikują to samo miejsce pamięci
nieatomowo, a~konfliktowe operacje nie są uporządkowane relacją
\emph{happens-before}, program ma \emph{data race} i~niezdefiniowane zachowanie.
Samo oznaczenie obiektu przez \code{volatile} nie rozwiązuje tego problemu.

\subsection{Porządek modyfikacji (\emph{modification order}) i~spójność}

Dla \emph{każdego pojedynczego} obiektu atomowego istnieje całkowity
porządek jego modyfikacji (\emph{modification order}). Nie wynika z~tego
jeden wspólny porządek modyfikacji \emph{wszystkich} atomików.
Ważne konsekwencje to między innymi:
\begin{itemize}
\item zapis do \code{x} uporządkowany przed innym zapisem do \code{x}
      musi być odpowiednio umieszczony przed nim w~porządku modyfikacji;
\item jeśli odczyt \code{x} jest uporządkowany przed kolejnym odczytem
      \code{x} w~tym samym wątku, drugi odczyt nie może cofnąć się
      do wcześniejszej modyfikacji w~porządku tego atomika;
\item ograniczenia spójności dla jednego atomika nie zapewniają
      automatycznie spójnego obrazu \emph{dwóch różnych} atomików.
\end{itemize}

\section{Testy modelu pamięci: wyniki, które zaskakują}
\label{sec:cpp16-litmus}

\subsection{\emph{Store buffering}: dwa odczyty zera}

Załóżmy, że \code{x} i~\code{y} początkowo są równe zero, a~wątki
wykonują następujące operacje:
\[
\begin{array}{c|c}
\text{Wątek 1} & \text{Wątek 2}\\ \hline
x\gets 1 & y\gets 1\\
r_1\gets y & r_2\gets x
\end{array}
\]

\begin{lstlisting}[caption={Store buffering: kompletna, wolna od wyścigu danych próba (C++20).}]
#include <atomic>
#include <iostream>
#include <thread>

int main() {
    std::atomic<int> x{0}, y{0};
    int r1 = -1, r2 = -1;

    std::jthread a([&] {
        x.store(1, std::memory_order_relaxed);
        r1 = y.load(std::memory_order_relaxed);
    });
    std::jthread b([&] {
        y.store(1, std::memory_order_relaxed);
        r2 = x.load(std::memory_order_relaxed);
    });
    a.join();
    b.join();
    std::cout << r1 << ' ' << r2 << '\n';
}
\end{lstlisting}

W~tym programie \(r_1=r_2=0\) jest \textbf{dozwolonym} wynikiem
modelu C++ dla \code{relaxed}; nie ma gwarancji, że zostanie
zaobserwowany w~praktycznym uruchomieniu. Jeśli \emph{wszystkie cztery}
operacje atomowe zamienimy na \code{memory_order_seq_cst}, wynik
\((0,0)\) stanie się niemożliwy: wspólny porządek operacji
sekwencyjnie spójnych musiałby przeczyć kolejności w~każdym z~wątków.

\nota{Wynik pojedynczego eksperymentu, nawet powtórzony milion razy,
nie dowodzi niedopuszczalności innego zachowania. Test tego rodzaju
to materiał do rozumowania nad modelem, a~nie pełny dowód poprawności.}

\subsection{\emph{Message passing}: publikacja zwykłych danych}

\begin{lstlisting}[caption={Publikacja danych przy użyciu release/acquire (C++11).}]
#include <atomic>
#include <cassert>
#include <thread>

int main() {
    int payload = 0;
    std::atomic<bool> published{false};

    std::thread producer([&] {
        payload = 42;
        published.store(true, std::memory_order_release);
    });
    std::thread consumer([&] {
        while (!published.load(std::memory_order_acquire)) {
            std::this_thread::yield();
        }
        assert(payload == 42);
    });

    producer.join();
    consumer.join();
}
\end{lstlisting}

Odczyt flagi, który widzi \code{true}, synchronizuje się z~zapisem
\code{release}. Zapis \code{payload = 42} następuje \emph{happens-before}
odczytu \code{payload} przez konsumenta. Flaga \emph{publikuje}
wcześniejszy stan; nie należy po publikacji modyfikować
\code{payload} bez dalszej synchronizacji.

\section{Porządki atomowe i~ich dokładne zastosowania}
\label{sec:cpp16-orders}

\begin{center}
\small
\begin{tabularx}{\textwidth}{@{}lX@{}}
\toprule
\textbf{Porządek} & \textbf{Typowe zastosowanie i~ograniczenie} \\
\midrule
\code{relaxed} & Liczniki, statystyki, pola kontrolne, gdy
sama atomowość wystarcza, a~relacja synchronizacji wynika skądinąd. \\
\code{release} & Zapis lub operacja RMW publikująca wcześniejsze efekty
wątku. \\
\code{acquire} & Odczyt lub operacja RMW odbierająca stan od publikującego wątku. \\
\code{acq_rel} & Operacja odczyt--modyfikacja--zapis wymagająca obu kierunków. \\
\code{seq_cst} & Operacje z~dodatkową gwarancją pojedynczego wspólnego
porządku dla operacji sekwencyjnie spójnych. \\
\bottomrule
\end{tabularx}
\end{center}

\code{memory_order_consume} ma długą i~skomplikowaną historię
standaryzacji. W~praktycznym kodzie C++20 należy preferować
\code{acquire} zamiast opierania przenośnej poprawności na
\code{consume}. Uwaga: porządek \code{release} jest właściwy dla
zapisu, \code{acquire} dla odczytu, a~\code{acq_rel} dla operacji
odczyt--modyfikacja--zapis; nie wszystkie porządki są poprawne dla
każdego przeciążenia funkcji atomowej.

\subsection{Operacje \emph{read--modify--write}}

Metody \code{fetch_add}, \code{fetch_sub}, \code{exchange} i~udany
\code{compare_exchange_*} wykonują operację RMW. Są one niepodzielne
w~odniesieniu do danego atomika, ale nie tworzą automatycznie
transakcji obejmującej inne pola.

\begin{lstlisting}[caption={Przydzielanie kolejnych numerów: wystarcza relaxed (C++11).}]
#include <atomic>
#include <cassert>
#include <thread>

int main() {
    std::atomic<int> next_id{0};
    int first = -1, second = -1;
    std::thread a([&] {
        first = next_id.fetch_add(1, std::memory_order_relaxed);
    });
    std::thread b([&] {
        second = next_id.fetch_add(1, std::memory_order_relaxed);
    });
    a.join();
    b.join();
    assert(first != second);
    assert(first + second == 1);
}
\end{lstlisting}

Uzyskujemy unikatowe wartości, bo RMW mają jednoznaczny porządek
modyfikacji \code{next_id}. Nie uzyskujemy przez to żadnej dodatkowej
synchronizacji niezwiązanych danych.

\subsection{\emph{Release sequence}: łańcuch operacji RMW}

W~C++20 zapis z~porządkiem \code{release} może rozpoczynać sekwencję
\emph{release} składającą się z~następujących po nim operacji RMW
na \emph{tym samym} atomiku, wykonywanych także przez inne wątki.
Odczyt \code{acquire} wartości pochodzącej z~tej sekwencji może
synchronizować się z~początkowym zapisem \code{release}.

\begin{lstlisting}[caption={Przekazanie publikacji przez relaxed RMW (C++20).}]
#include <atomic>
#include <cassert>
#include <thread>

int main() {
    int payload = 7;
    std::atomic<int> state{0};

    std::thread producer([&] {
        payload = 42;
        state.store(1, std::memory_order_release);
    });
    std::thread intermediary([&] {
        int expected = 0;
        while ((expected = state.load(std::memory_order_acquire))
               == 0) {
            std::this_thread::yield();
        }
        state.fetch_add(1, std::memory_order_relaxed);
    });
    std::thread consumer([&] {
        while (state.load(std::memory_order_acquire) != 2) {
            std::this_thread::yield();
        }
        assert(payload == 42);
    });

    producer.join();
    intermediary.join();
    consumer.join();
}
\end{lstlisting}

Wątek pośredni czeka, aż stan osiągnie 1, i~wykonuje RMW, który
zapisuje 2. Odczyt wartości 2 z~\code{acquire} synchronizuje się
z~pierwotnym \code{release} przez sekwencję \emph{release}.
Nie można w~ogólności zastąpić \code{fetch_add} zwykłym
\code{store(2, relaxed)} i~utrzymać tej samej argumentacji.

\section{CAS: porównaj i~zamień, pętle ponawiania i~porządek niepowodzenia}

\begin{center}\small
\begin{tabularx}{\textwidth}{@{}p{0.38\textwidth}X@{}}
\toprule
Operacja atomowa & Działanie \\
\midrule
\code{compare_exchange_strong} & Jednorazowo porównuje i próbuje podmienić wartość; nie zgłasza pozornego niepowodzenia. \\
\code{compare_exchange_weak} & Może zgłosić pozorne niepowodzenie; nadaje się do pętli ponawiania. \\
\code{expected} przy CAS & Przy niepowodzeniu otrzymuje rzeczywistą wartość atomika, którą trzeba uwzględnić przy kolejnej próbie. \\
\code{memory_order_acq_rel} & Przy udanej operacji odczytująco-zapisującej łączy własności acquire i release. \\
\code{memory_order_relaxed} & Zapewnia atomowość bez dodatkowej synchronizacji z innymi danymi. \\
\bottomrule
\end{tabularx}
\end{center}
\label{sec:cpp16-cas}

\code{compare_exchange_weak(expected, desired, success, failure)}
atomowo porównuje bieżącą wartość z~\code{expected}. Gdy porównanie
się powiedzie, zapisuje \code{desired}; w~przeciwnym razie aktualizuje
\code{expected} do wartości zastanej w~atomiku. Wersja \code{weak}
może ponadto nie powieść się pozornie (\emph{spurious failure}).
W~pętli jest zwykle dobrym wyborem.

\begin{lstlisting}[caption={Atomowe maksimum z~pętlą CAS (C++11).}]
#include <atomic>
#include <cassert>

void atomic_max(std::atomic<int>& target, int candidate) {
    int old = target.load(std::memory_order_relaxed);
    while (old < candidate &&
           !target.compare_exchange_weak(
               old, candidate,
               std::memory_order_relaxed,
               std::memory_order_relaxed)) {
        // Przy niepowodzeniu old jest aktualizowane.
    }
}

int main() {
    std::atomic<int> maximum{10};
    atomic_max(maximum, 8);
    atomic_max(maximum, 17);
    assert(maximum.load() == 17);
}
\end{lstlisting}

Parametr \code{failure} opisuje porządek operacji \emph{odczytu}
w~przypadku niepowodzenia. Nie może to być \code{release} ani
\code{acq_rel}; musi też spełniać ograniczenia względem porządku
sukcesu określone standardem. Typowy wzorzec publikacji nowej
wartości to sukces z~\code{acq_rel} i~niepowodzenie z~\code{acquire}
(albo \code{relaxed}, jeśli tak wynika z~dowodu).

\subsection{\code{weak} kontra \code{strong}}

\code{compare_exchange_strong} nie zgłasza pozornego niepowodzenia
przy poprawnie porównanych wartościach. Może być użyteczne dla
pojedynczej próby; \code{weak} bywa wygodniejsze w~pętlach CAS.
Nie ma uniwersalnej gwarancji, że jedna z~tych wersji zawsze działa
szybciej na danej platformie.

\section{Bariery pamięci i~\texttt{std::atomic\_thread\_fence}}
\label{sec:cpp16-fences}

Bariera \code{std::atomic_thread_fence(order)} jest operacją
porządkowania pamięci, która sama nie odczytuje ani nie zapisuje
konkretnego atomika. Synchronizacja z~innym wątkiem wciąż wymaga
odpowiednich operacji atomowych łączących obie strony. Samo
wstawienie \code{fence} obok zwykłej zmiennej nie usuwa wyścigu danych.

\begin{lstlisting}[caption={Publikacja ze zwolnieniową i~przejmującą barierą pamięci (C++11).}]
#include <atomic>
#include <cassert>
#include <thread>

int main() {
    int payload = 0;
    std::atomic<bool> flag{false};

    std::thread producer([&] {
        payload = 42;
        std::atomic_thread_fence(std::memory_order_release);
        flag.store(true, std::memory_order_relaxed);
    });
    std::thread consumer([&] {
        while (!flag.load(std::memory_order_relaxed)) {
            std::this_thread::yield();
        }
        std::atomic_thread_fence(std::memory_order_acquire);
        assert(payload == 42);
    });

    producer.join();
    consumer.join();
}
\end{lstlisting}

To przykład synchronizacji \emph{fence--fence}: odczyt atomika
widzi zapis następujący po barierze \code{release}, a~bariera
\code{acquire} jest wykonywana po tym odczycie. W~nowym kodzie
zwykle prostsze i~czytelniejsze jest zastosowanie pary operacji
\code{release}/\code{acquire} bez jawnych barier. Bariery nie są
substytutem przemyślanego protokołu komunikacji.

\section{Gwarancje postępu: \emph{blocking}, \emph{lock-free} i~\emph{wait-free}}
\label{sec:cpp16-progress}

Wielowątkowy algorytm można analizować pod względem \emph{bezpieczeństwa}
(\emph{safety}: nic niedozwolonego się nie dzieje) i~\emph{żywotności}
(\emph{liveness}: coś pożądanego w~końcu się wydarzy). Są to różne kwestie.

\begin{description}[style=nextline]
\item[\emph{Blocking}]
  Wątek może oczekiwać na odblokowanie mutexa przez konkretny inny wątek.
  Jeśli właściciel przestanie robić postępy, pozostali również mogą utknąć.
\item[\emph{Lock-free}]
  W~standardowym ujęciu postępu algorytmowego, mimo rywalizacji wątków,
  system jako całość może kontynuować wykonywanie operacji:
  niepowodzenia jednej próby CAS zwykle oznaczają, że inny uczestnik
  wykonał postęp. Nie gwarantuje to postępu \emph{każdego} wątku.
\item[\emph{Wait-free}]
  Każda operacja danego wątku ma ograniczoną liczbę własnych kroków
  niezależnie od działań innych wątków (przy przyjętym modelu).
  To silniejszy warunek niż \emph{lock-free}.
\item[\emph{Obstruction-free}]
  Operacja postępuje, jeśli wątek może działać dostatecznie długo
  bez konfliktów ze strony innych wątków. Jest to słabsza gwarancja.
\end{description}

\emph{Lock-free} nie znaczy ,,zawsze szybciej''. Nie oznacza również
braku aktywnego oczekiwania ani braku zagłodzenia konkretnego wątku.
Odróżniajmy też własność implementacji pojedynczego atomika
(\code{is_lock_free()}) od postępu kompletnego algorytmu.

\subsection{Prosta blokada wirująca (\emph{spinlock})}

\begin{lstlisting}[caption={Poprawna blokada wirująca z flagą atomową; ilustracja (C++20).}]
#include <atomic>
#include <cassert>
#include <mutex>
#include <thread>

class SpinLock {
    std::atomic_flag flag = ATOMIC_FLAG_INIT;
public:
    void lock() noexcept {
        while (flag.test_and_set(std::memory_order_acquire)) {
            // Zamiast stale rywalizowac o cache line, czekamy.
            flag.wait(true, std::memory_order_relaxed);
        }
    }
    void unlock() noexcept {
        flag.clear(std::memory_order_release);
        flag.notify_one();
    }
};

int main() {
    SpinLock lock;
    int counter = 0;
    auto work = [&] {
        for (int i = 0; i < 10000; ++i) {
            std::lock_guard<SpinLock> guard(lock);
            ++counter;
        }
    };
    std::jthread a(work), b(work);
    a.join();
    b.join();
    assert(counter == 20000);
}
\end{lstlisting}

\code{test_and_set} ma semantykę RMW; odblokowanie jest operacją
\code{release}. Metoda \code{wait(true)} może blokować, a~zatem
powyższy \code{SpinLock} \emph{nie} jest algorytmem \emph{lock-free}.
Ponadto nie gwarantuje sprawiedliwości. W~realnym oprogramowaniu
\code{std::mutex} zwykle pozostaje lepszym wyborem, zwłaszcza gdy
sekcja krytyczna może być długa lub obciążenie jest wysokie.

\section{Problem ABA: poprawne CAS nie wystarcza}
\label{sec:cpp16-aba}

Rozważmy wierzchołek stosu reprezentowany wskaźnikiem atomowym.
Wątek \(A\) czyta wskaźnik \(p\) i~zamierza wykonać
\code{compare_exchange} z~wartości \(p\) na \code{p->next}.
W~tym czasie wątek \(B\) może:
\begin{enumerate}
\item zdjąć element \(p\) ze stosu;
\item zdjąć jeszcze jeden element;
\item ponownie umieścić \(p\) na wierzchołku stosu.
\end{enumerate}

Wartość wskaźnika uległa zmianie \(A\to B\to A\) (stąd \emph{ABA}),
ale CAS wątku \(A\) może zobaczyć ponownie tę samą wartość \(p\).
W~międzyczasie struktura stosu mogła się całkowicie zmienić.

\begin{center}
\(\boxed{\text{head}=p}\;\longrightarrow\;
  \boxed{\text{head}=q}\;\longrightarrow\;
  \boxed{\text{head}=p}\)
\end{center}

Dwa problemy trzeba odróżnić:
\begin{enumerate}
\item \emph{ABA logiczne}: porównanie tej samej wartości nie wykrywa
      zmian stanu struktury.
\item \emph{Bezpieczeństwo pamięci}: wątek \(A\) może próbować odczytać
      \code{p->next}, gdy inny wątek zniszczył już obiekt \code{*p}.
      Taki odczyt nie staje się poprawny tylko dlatego, że \code{head}
      jest atomowy.
\end{enumerate}

\subsection{Wskaźnik z~licznikiem generacji}

Jedna z~technik polega na porównywaniu pary
\((\text{wskaźnik},\text{licznik wersji})\) zamiast samego wskaźnika.
Wtedy ponowne pojawienie się tego samego adresu nie oznacza równości
całego stanu. Licznik również ma ograniczony zakres i~może się
przepełnić; trzeba uzasadnić brak niebezpiecznego \emph{wrap-around}.

\nota{Nie zakładaj, że \code{std::atomic<SomeStruct>} będzie bez blokad:
implementacja może wykorzystać wewnętrzny mutex lub bibliotekę
atomowych operacji. Możliwość zapisania stanu złożonego w~\code{atomic}
nie rozwiązuje też samodzielnie problemu czasu życia obiektów.}

\subsection{Hazard pointers, epoch-based reclamation i~RCU}

\term{Hazard pointers} pozwalają wątkom publikować informację:
,,zamierzam korzystać z~tego obiektu''. Wątek usuwający obiekt
odkłada go na listę \emph{retired} i~odracza zniszczenie do chwili,
gdy żaden aktywny hazard pointer go nie wskazuje. Poprawna procedura
zwykle musi ponownie zweryfikować odczyt wskaźnika po publikacji
hazardu; samo opublikowanie dowolnego, potencjalnie już nieważnego
adresu nie wystarcza.

\term{Epoch-based reclamation} grupuje operacje w~epoki. Usunięte
obiekty mogą zostać odzyskane dopiero wtedy, gdy wszystkie aktywne
wątki zakończyły okresy korzystania z~danych obejmujące ich usunięcie.
Dłużej zatrzymany wątek może opóźniać zwolnienie pamięci.

\term{Read-copy-update (RCU)} korzysta z~odczytów starych wersji
danych i~publikacji nowych wersji; usuwanie poprzednich odbywa się
po odpowiednim okresie ochronnym (\emph{grace period}). Szczegóły
różnią się pomiędzy implementacjami.

\begin{center}
\small
\begin{tabularx}{\textwidth}{@{}lX@{}}
\toprule
\textbf{Metoda} & \textbf{Kiedy rozważyć} \\
\midrule
Hazard pointers & Wiele odczytów wskaźników, częste usuwanie
węzłów; możliwa kosztowna administracja hazardami. \\
Epoki & Duży ruch, kontrolowane sekcje dostępu; zatrzymany czytelnik
może opóźnić odzyskanie pamięci. \\
RCU & Wzorce ,,wielu czytelników, niewielu pisarzy'', publikacja
niezmiennych wersji danych. \\
\code{shared_ptr} & Gdy łatwość dowodu poprawności jest ważniejsza
od kosztu liczenia referencji; nie obiecuje \emph{lock-free}. \\
\bottomrule
\end{tabularx}
\end{center}

W~C++20 standardowa biblioteka nie udostępnia jeszcze kompletnej,
przenośnej abstrakcji hazard pointerów. Bezpieczna implementacja
wymaga precyzyjnego protokołu pamięciowego, dlatego zamiast podawać
uproszczony, ale błędny stos ze zwalnianiem węzłów, przedstawimy
algorytm bez dynamicznego odzyskiwania pamięci: kolejkę SPSC.

\section{Kolejka SPSC: praktyczny algorytm bez mutexa}
\label{sec:cpp16-spsc}

\emph{Single Producer, Single Consumer} (SPSC) oznacza dokładnie
jednego producenta i~jednego konsumenta. Bufor cykliczny może wtedy
używać dwóch liczników:
\begin{itemize}
\item \code{head} --- numer kolejnego miejsca zapisu, modyfikowany
      wyłącznie przez producenta;
\item \code{tail} --- numer kolejnego miejsca odczytu, modyfikowany
      wyłącznie przez konsumenta.
\end{itemize}

Producent zapisuje komórkę, a~następnie publikuje nowy \code{head}
operacją \code{release}. Konsument odczytuje \code{head} przez
\code{acquire}, a~więc może bezpiecznie zobaczyć zawartość wcześniej
zapisanej komórki. Analogicznie konsument zwalnia komórkę, publikując
nowy \code{tail}, a~producent odbiera tę informację operacją
\code{acquire}. Warunek pełności bufora o~pojemności \(C\) to
\(\mathrm{head}-\mathrm{tail}=C\), a~pustości:
\(\mathrm{head}=\mathrm{tail}\).

\begin{lstlisting}[caption={Kompletna kolejka SPSC dla int, ze statyczną pamięcią (C++20).}]
#include <array>
#include <atomic>
#include <cassert>
#include <cstddef>
#include <thread>

// Tylko jeden producent i dokladnie jeden konsument.
template<std::size_t Capacity>
class SpscQueue {
    static_assert(Capacity > 0);
    static_assert((Capacity & (Capacity - 1)) == 0);
    std::array<int, Capacity> buffer_{};
    std::atomic<std::size_t> head_{0};
    std::atomic<std::size_t> tail_{0};
public:
    bool try_push(int value) {
        const auto head = head_.load(std::memory_order_relaxed);
        const auto tail = tail_.load(std::memory_order_acquire);
        if (head - tail == Capacity) return false;
        buffer_[head & (Capacity - 1)] = value;
        head_.store(head + 1, std::memory_order_release);
        return true;
    }

    bool try_pop(int& result) {
        const auto tail = tail_.load(std::memory_order_relaxed);
        const auto head = head_.load(std::memory_order_acquire);
        if (tail == head) return false;
        result = buffer_[tail & (Capacity - 1)];
        tail_.store(tail + 1, std::memory_order_release);
        return true;
    }
};

int main() {
    constexpr int Count = 100000;
    SpscQueue<1024> queue;
    long long result = 0;

    std::jthread producer([&] {
        for (int x = 1; x <= Count; ++x) {
            while (!queue.try_push(x)) std::this_thread::yield();
        }
    });
    std::jthread consumer([&] {
        int item = 0;
        for (int n = 0; n < Count; ++n) {
            while (!queue.try_pop(item)) std::this_thread::yield();
            result += item;
        }
    });
    producer.join();
    consumer.join();
    assert(result == 1LL * Count * (Count + 1) / 2);
}
\end{lstlisting}

\subsection{Dlaczego kod jest poprawny?}

Dla danych w~\code{buffer_} nie stosujemy operacji atomowych.
Nie ma jednak konfliktowego dostępu bez synchronizacji, ponieważ:
\begin{enumerate}
\item przed odczytem komórki konsument otrzymuje informację o~zapisie
      przez \code{head_.load(acquire)};
\item przed ponownym nadpisaniem tej komórki producent otrzymuje
      informację o~jej zwolnieniu przez \code{tail_.load(acquire)};
\item każdy z~liczników ma jednego piszącego, więc nie potrzeba CAS
      do aktualizowania \code{head_} ani \code{tail_}.
\end{enumerate}

Przykład zakłada, że operacji \code{try_push} nie wywołuje kilku
producentów, a~\code{try_pop} kilku konsumentów, nawet jeśli
same metody są dostępne publicznie. Złamane założenie SPSC
powoduje utratę poprawności. Ograniczenia bufora, przepełnienie
liczników i~kwestie typów niebanalnych wymagają osobnej analizy.
Użyliśmy potęgi dwójki, żeby indeksowanie przez maskę pozostało
poprawne także przy modularnym zawijaniu licznika bez znaku.

\nota{Ten przykład jest edukacyjną, \emph{nieblokującą} konstrukcją
w~sensie, że nie przejmuje mutexa, ale aktywne pętle producenta
i~konsumenta mogą czekać dowolnie długo. Metody \code{try_push}
i~\code{try_pop} są krótkie i~nie wykonują pętli CAS. Bez uzasadnienia
własności użytych atomików na danej platformie nie deklarujemy
silniejszej gwarancji wydajnościowej.}

\section{Atomowa publikacja niezmiennych wersji danych}
\label{sec:cpp16-snapshot}

Wiele rzeczywistych aplikacji (konfiguracja serwera, słowniki,
parametry modeli) posiada licznych czytelników i~rzadkich pisarzy.
Dobrą architekturą bywa utworzenie nowej, niezmiennej wersji
struktury i~jednorazowe opublikowanie wskaźnika do niej.
W~C++20 można do tego wykorzystać specjalizację
\code{std::atomic<std::shared_ptr<T>>}.

\begin{lstlisting}[caption={Bezpieczny snapshot konfiguracji (C++20).}]
#include <atomic>
#include <cassert>
#include <memory>
#include <thread>

struct Config {
    int threshold;
    double scale;
};

int main() {
    using Ptr = std::shared_ptr<const Config>;
    std::atomic<Ptr> current{
        std::make_shared<const Config>(Config{10, 1.0})
    };

    std::jthread writer([&] {
        Ptr next = std::make_shared<const Config>(Config{20, 2.0});
        current.store(std::move(next), std::memory_order_release);
    });
    std::jthread reader([&] {
        Ptr snap = current.load(std::memory_order_acquire);
        assert((snap->threshold == 10 && snap->scale == 1.0) ||
               (snap->threshold == 20 && snap->scale == 2.0));
    });
    writer.join();
    reader.join();
}
\end{lstlisting}

Czytelnik otrzymuje \emph{spójny} snapshot obu pól, a~stara wersja
jest automatycznie utrzymywana przy życiu, dopóki istnieją do niej
uchwyty \code{shared_ptr}. Nie oznacza to, że przechowywanie
\code{shared_ptr} w~atomiku jest \emph{lock-free}. Referencyjne
liczenie i~alokacje bywają istotnym kosztem przy bardzo częstych
odczytach. Dla obciążenia z~dużą przewagą czytelników należy
porównać to podejście z~\code{shared_mutex}, kopiowaniem małych
struktur i~technikami RCU.

\subsection{\code{std::atomic_ref}: synchronizacja istniejącego obiektu}

\code{std::atomic_ref<T>} (C++20) pozwala wykonywać operacje atomowe
na istniejącym, odpowiednio wyrównanym obiekcie typu \code{T}.
Przydatne bywa przy interoperacyjności ze strukturami danych,
które nie przechowują \code{std::atomic<T>} jako pola.

\begin{lstlisting}[caption={Operacje atomowe na istniejącej zmiennej (C++20).}]
#include <atomic>
#include <cassert>
#include <thread>

int main() {
    alignas(std::atomic_ref<int>::required_alignment) int counter = 0;
    {
        std::atomic_ref<int> shared_counter(counter);
        auto add = [&] {
            for (int i = 0; i < 10000; ++i)
                shared_counter.fetch_add(1, std::memory_order_relaxed);
        };
        std::jthread a(add), b(add);
        a.join();
        b.join();
    } // koniec zycia atomic_ref
    assert(counter == 20000);
}
\end{lstlisting}

Dopóki istnieje \code{atomic_ref} odnoszący się do obiektu, dostęp
do tego obiektu należy realizować przez instancje \code{atomic_ref}
--- nie tylko w~chwilach równoczesnego dostępu. Dlatego zwykłe
odczytanie \code{counter} wykonujemy \emph{po} zakończeniu czasu
życia referencji atomowej. Sam obiekt musi żyć dostatecznie długo,
a~jego wyrównanie musi spełniać wymagania \code{required_alignment}.

\section{Planowanie zadań: pule wątków, \emph{work stealing} i~\emph{backpressure}}
\label{sec:cpp16-scheduling}

Poza poprawną synchronizacją potrzebujemy rozsądnej strategii
przypisywania prac do wątków. Wątek systemowy ma koszt utworzenia,
przełączenia kontekstu i~utrzymania stosu. Z~tego powodu tysiące
krótkich zadań nie powinny automatycznie oznaczać tysiąca wątków.

\begin{description}[style=nextline]
\item[\emph{Thread pool}]
  Ustalona liczba wątków pobiera zadania z~kolejki. Ogranicza koszt
  częstego tworzenia wątków i~może ograniczać równoległe obciążenie.
\item[\emph{Work stealing}]
  Każdy wątek ma własną kolejkę, a~bezczynny pracownik próbuje
  przejąć część pracy od innych. Dobrze sprawdza się przy nierównych
  czasach trwania podzadań. Wymaga ostrożnego protokołu synchronizacji
  kolejek, szczególnie gdy korzystają z~nich różni producenci.
\item[\emph{Backpressure}]
  Mechanizm spowalniania producenta, gdy konsumenci nie nadążają.
  Bez niego kolejka bez ograniczenia pojemności może zużyć całą pamięć.
\item[\emph{Structured concurrency}]
  Sposób organizacji zadań, w~którym ich czas życia, anulowanie i~błędy
  pozostają związane z~zakresem nadrzędnej operacji. Sam \code{std::async}
  nie implementuje kompletnego frameworka takiego stylu.
\end{description}

\subsection{Pula wątków: bezpieczne zamknięcie i~propagacja wyjątków}

Następny przykład rozwija pulę z~rozdziału~\ref{chap:concurrency}. Używa mutexa
oraz zmiennej warunkowej. Łatwy dowód poprawności i~przewidywalne
zachowanie przy zamykaniu są często ważniejsze niż
niepotwierdzony \emph{lock-free}.

\begin{lstlisting}[caption={Prosta pula wątków odprowadzająca kolejkę przed zamknięciem (C++20).}]
#include <condition_variable>
#include <functional>
#include <future>
#include <iostream>
#include <memory>
#include <mutex>
#include <queue>
#include <stdexcept>
#include <thread>
#include <type_traits>
#include <utility>
#include <vector>

class ThreadPool {
    std::mutex mutex_;
    std::condition_variable cv_;
    std::queue<std::function<void()>> jobs_;
    bool closed_ = false;
    std::vector<std::jthread> workers_;

    void worker_loop() {
        for (;;) {
            std::function<void()> job;
            {
                std::unique_lock lock(mutex_);
                cv_.wait(lock, [&] {
                    return closed_ || !jobs_.empty();
                });
                if (closed_ && jobs_.empty()) return;
                job = std::move(jobs_.front());
                jobs_.pop();
            }
            job();  // packaged_task przechwyci wyjatek
        }
    }

public:
    explicit ThreadPool(unsigned count) {
        if (count == 0) throw std::invalid_argument("count=0");
        try {
            for (unsigned i = 0; i < count; ++i)
                workers_.emplace_back([this] { worker_loop(); });
        } catch (...) {
            close();  // obudz pracownikow; ich destruktory wykonaja join
            throw;
        }
    }
    ThreadPool(const ThreadPool&) = delete;
    ThreadPool& operator=(const ThreadPool&) = delete;
    ~ThreadPool() { close(); }

    void close() {
        {
            std::lock_guard lock(mutex_);
            closed_ = true;
        }
        cv_.notify_all();
    }

    template<class F>
    auto submit(F&& callable)
        -> std::future<std::invoke_result_t<std::decay_t<F>&>> {
        using R = std::invoke_result_t<std::decay_t<F>&>;
        auto task = std::make_shared<std::packaged_task<R()>>(
            std::forward<F>(callable)
        );
        auto result = task->get_future();
        {
            std::lock_guard lock(mutex_);
            if (closed_) throw std::runtime_error("pool closed");
            jobs_.emplace([task] { (*task)(); });
        }
        cv_.notify_one();
        return result;
    }
};

int main() {
    ThreadPool pool(2);
    auto a = pool.submit([] { return 40; });
    auto b = pool.submit([] { return 2; });
    std::cout << a.get() + b.get() << '\n'; // 42
    pool.close(); // od tej chwili bez nowych zadan
}
\end{lstlisting}

Ważne własności pokazanej implementacji:
\begin{itemize}
\item \code{close()} jest bezpieczne wobec równoległego \code{submit()}:
      oba działania są serializowane przez ten sam mutex;
\item po \code{close()} pracownicy \emph{kończą kolejkę}, a~dopiero
      następnie opuszczają pętlę; destruktory \code{jthread} dołączają
      pracowników, zanim zniszczone zostaną mutex i~kolejka;
\item \code{packaged_task} zapisuje wartość lub wyjątek w~stanie
      asynchronicznym, więc zwykłe wyjątki z~zadań nie wychodzą
      bezpośrednio przez \code{worker_loop};
\item ta minimalna pula nie ma ograniczenia długości kolejki,
      priorytetów, terminów ani anulowania uruchomionych prac.
\end{itemize}

\nota{Nie wolno niszczyć takiej puli z~zadania wykonywanego w~tej
samej puli: dołączenie własnego wątku prowadzi do błędu.
Podobnie zadanie czekające na inne zadanie z~tej samej puli może
zakleszczyć program, jeśli wszyscy pracownicy pozostają zajęci
oczekiwaniem. Produkcyjny executor wymaga rozwiązań dla takich sytuacji.}

\subsection{Kiedy potrzebny jest \emph{work stealing}?}

Załóżmy, że przetwarzamy drzewo lub dzielimy dużą macierz na bloki.
Część bloków kończy się po mikrosekundach, inne zajmują milisekundy.
Statyczny przydział zadań może zostawić jeden wątek przeciążony,
gdy inne nie mają już pracy. Przechwytywanie pracy pomiędzy lokalnymi
kolejkami poprawia \emph{równoważenie obciążenia}.

Typowy projekt stosuje:
\begin{enumerate}
\item właścicielski dostęp pracownika do jednego końca jego kolejki;
\item kradzież elementów przez inne wątki z~drugiego końca;
\item atomowe indeksy i~protokół obsługujący konflikt właściciela
      ze złodziejem o~ostatni element;
\item kontrolę czasu życia tablicy zadań oraz jej ewentualnej realokacji.
\end{enumerate}

Takie struktury są znacznie bardziej subtelne niż \code{std::deque}
z~mutexem. W~praktyce warto użyć sprawdzonego executora/biblioteki,
a~własną implementację traktować jako osobny, formalnie weryfikowany projekt.

\section{Anulowanie, sygnały stopu i~punkty oczekiwania}
\label{sec:cpp16-cancel}

Kooperacyjne zatrzymywanie w~C++20 jest oparte na
\code{std::stop_source}, \code{std::stop_token} i~\code{std::stop_callback}.
Żądanie zatrzymania nie zabija wątku; umożliwia mu dobrowolne
zakończenie po zweryfikowaniu stanu i~odpowiednim uporządkowaniu zasobów.

\begin{lstlisting}[caption={Wybudzanie i anulowanie czekania (C++20).}]
#include <condition_variable>
#include <iostream>
#include <mutex>
#include <thread>

int main() {
    std::mutex mutex;
    std::condition_variable_any cv;
    bool ready = false;

    std::jthread worker([&](std::stop_token stop) {
        std::unique_lock lock(mutex);
        bool completed = cv.wait(lock, stop, [&] { return ready; });
        std::cout << (completed ? "ready\n" : "cancelled\n");
    });

    worker.request_stop();
    worker.join();
}
\end{lstlisting}

W~tym przeciążeniu \code{wait} zwraca wartość predykatu po zakończeniu
czekania. Jeśli żądanie zatrzymania nastąpi przed spełnieniem
predykatu, przykład wypisze \code{cancelled}. Jeśli stan zmieniłby
się w~międzyczasie, wynik mógłby być inny. Zwykła
\code{std::condition_variable} nie posiada w~C++20 odpowiednika
tego przeciążenia z~\code{stop_token}.

\subsection{Anulowanie i~bezpieczeństwo wyjątkowe}

W~dobrym API warto określić, co dzieje się z~wynikiem po anulowaniu:
\begin{itemize}
\item czy zwracany jest częściowy wynik;
\item czy do \code{future} trafia własny wyjątek oznaczający anulowanie;
\item czy potrzebne jest odwrócenie skutków wcześniej wykonanych operacji;
\item jak zapewnić, że wszystkie zależne zadania zostaną zakończone
      przed zniszczeniem współdzielonych danych.
\end{itemize}

Żadne z~tych pytań nie ma jednej odpowiedzi dla wszystkich bibliotek.
Sama metoda \code{future.wait_for()} nie anuluje zadań, a~korutyna
z~rozdziału~\ref{chap:cpp15-coroutines} wymaga mechanizmu propagacji żądania anulowania
przez awaitery i~executor.

\section{Fałszywe współdzielenie (\emph{false sharing}) i~cache coherence}
\label{sec:cpp16-false-sharing}

Dwa niezależne atomiki mogą być zapisane blisko siebie w~pamięci,
tak że mieszczą się w~tej samej jednostce przenoszenia danych
pomiędzy pamięciami podręcznymi. Wtedy zapis przez różne rdzenie
może wymagać kosztownych transferów linii cache, choć z~punktu
widzenia C++ nie ma wyścigu danych.

\begin{lstlisting}[caption={Rozdzielenie intensywnie zapisywanych liczników (C++20; stała wyrównania od C++17).}]
#include <atomic>
#include <cstdint>
#include <iostream>
#include <new>
#include <thread>

struct alignas(std::hardware_destructive_interference_size) Counter {
    std::atomic<std::uint64_t> value{0};
};

int main() {
    Counter a, b;
    auto increment = [](Counter& c) {
        for (int i = 0; i < 100000; ++i)
            c.value.fetch_add(1, std::memory_order_relaxed);
    };
    std::jthread first([&] { increment(a); });
    std::jthread second([&] { increment(b); });
    first.join();
    second.join();
    std::cout << a.value << ' ' << b.value << '\n';
}
\end{lstlisting}

\code{std::hardware_destructive_interference_size} jest wskazówką
implementacji dotyczącą rozdzielania obiektów często modyfikowanych
przez różne wątki; w~praktyce jej obsługa zależy od kompilatora
i~biblioteki standardowej. \code{alignas} może zwiększać zużycie
pamięci, zatem nie należy bez pomiarów wyrównywać w~ten sposób
każdego pola. Rozdział~\ref{chap:cpp-memory-performance} omawia szerzej cache locality i~organizację
danych AoS/SoA.

\subsection{Co mierzyć przy ocenie skalowalności?}

Pojedynczy czas wykonania jest mało informacyjny. Przy porównywaniu
wersji jednowątkowej i~wielowątkowej warto mierzyć:
\begin{itemize}
\item przepustowość (operacje/sekundę), a~nie tylko opóźnienie
      pojedynczego zadania;
\item \emph{speedup} \(S(p)=T(1)/T(p)\), dla \(p\) wątków;
\item rozkład opóźnień, np. medianę i~wysokie percentyle;
\item liczbę alokacji, przełączeń kontekstu, braków w~cache i~konfliktów CAS;
\item oddzielnie okres rozgrzewki oraz właściwy okres pomiarowy.
\end{itemize}

Według prawa Amdahla, jeśli część \(s\) programu musi pozostać
sekwencyjna, to w~idealizacji
\[
S(p)\leq\frac{1}{s+\frac{1-s}{p}}.
\]
Wzór ten nie obejmuje kosztów synchronizacji, a~więc rzeczywiste
przyspieszenie może być jeszcze niższe. Jednocześnie program
z~dominującym kosztem dostępu do pamięci może nie skalować się
nawet przy bardzo małej części sekwencyjnej.

\section{Linearizowalność i~niezmienniki współdzielonych struktur}
\label{sec:cpp16-linearizability}

Dla współbieżnych kolekcji często potrzebujemy mocniejszej własności
niż brak wyścigów danych: \emph{linearizowalności}.
Operacja jest linearizowalna, jeśli można przypisać jej jeden
umowny moment pomiędzy wywołaniem i~zakończeniem, w~którym
wywołanie ma taki sam skutek jak operacja wykonana sekwencyjnie.

Dla prostej atomowej inkrementacji punktem linearyzacji jest
sam RMW. Dla kolejki algorytm jest bardziej subtelny:
\begin{itemize}
\item producent najpierw zapisuje komórkę, a~następnie publikuje
      nowy \code{head};
\item konsument najpierw odczytuje i~przetwarza komórkę, a~następnie
      publikuje nowy \code{tail};
\item sama komórka nie musi być atomowa, o~ile zachowana jest
      właściwa synchronizacja i~protokół własności miejsca w~buforze.
\end{itemize}

W~algorytmach złożonych, np. z~dwoma atomikami \code{size} i~\code{head},
każde z~pól może być wolne od wyścigu danych, ale ich jednoczesny
odczyt nie musi stanowić spójnego snapshotu. Osobny \emph{niezmiennik}
struktury wymaga własnego uzasadnienia.

\subsection{Dlaczego prosty \emph{seqlock} nie wystarczy dla zwykłych danych?}

Popularny pomysł polega na tym, że pisarz zwiększa licznik wersji
przed i~po aktualizacji obiektu, a~czytelnik ponawia odczyt, jeżeli
wersja zmieniła się podczas odczytu. Jeżeli pola obiektu są zwykłymi
zmiennymi C++, czytelnik nadal może czytać je jednocześnie z~zapisem,
co oznacza \emph{data race}. Ponowienie odczytu \emph{po fakcie}
nie naprawia niezdefiniowanego zachowania, które już nastąpiło.

Poprawny wariant dla przenośnego C++ musi zastosować odpowiednią
synchronizację lub atomowe pola danych, a~nie tylko atomowy licznik
wersji. Alternatywą jest pokazana wcześniej publikacja niezmiennych
snapshotów.

\section{Współbieżność a~parallel STL, Ranges i~OpenMP}
\label{sec:cpp16-interop}

Nie każda równoległa praca wymaga własnej kolejki czy atomowych
liczników. Gdy zadanie polega na niezależnym przetwarzaniu elementów,
lepszą abstrakcją może być algorytm równoległy lub dyrektywa OpenMP.

\subsection{Polityki wykonania standardowych algorytmów}

Nagłówek \code{<execution>} udostępnia od C++17 między innymi:
\begin{description}[style=nextline]
\item[\code{std::execution::seq}] Wykonanie sekwencyjne.
\item[\code{std::execution::par}] Implementacja może wykonywać pracę
      równolegle; nie daje gwarancji, że rzeczywiście użyje wielu wątków.
\item[\code{std::execution::par_unseq}] Zezwala dodatkowo na
      niesekwencyjne wykonywanie, w~tym wektoryzację; nakłada ważne
      ograniczenia na wykonywane funkcje i~synchronizację.
\end{description}

\begin{lstlisting}[caption={Równoległa transformacja bez współdzielonych zapisów (C++17).}]
#include <algorithm>
#include <execution>
#include <numeric>
#include <vector>

int main() {
    std::vector<double> a(10000), b(a.size());
    std::iota(a.begin(), a.end(), 0.0);
    std::transform(std::execution::par,
                   a.begin(), a.end(), b.begin(),
                   [](double x) { return x * x; });
}
\end{lstlisting}

Kod nie zakłada kolejności przetwarzania elementów. Wyjściowy zakres
\code{b} nie nakłada się z~\code{a}, a~każde wywołanie funkcji
zapisuje inny element. Rzeczywiste zrównoleglenie zależy od
implementacji biblioteki standardowej i~konfiguracji środowiska;
niektóre kompilacje wymagają dodatkowej biblioteki wykonawczej.

\subsection{C++ i~OpenMP: dwa poziomy abstrakcji}

OpenMP, omawiany szczegółowo w~rozdziale~\ref{chap:concurrency}, umożliwia między innymi
równoległe pętle, redukcje i~zależności pomiędzy zadaniami.
W~obliczeniach numerycznych typowe rozwiązanie wygląda tak:

\begin{lstlisting}[caption={Wektorowa redukcja z OpenMP; wymagane -fopenmp.}]
#include <iostream>
#include <vector>

int main() {
    const std::vector<double> x(100000, 1.0);
    double sum = 0.0;

    #pragma omp parallel for reduction(+:sum) schedule(static)
    for (int i = 0; i < static_cast<int>(x.size()); ++i) {
        sum += x[static_cast<std::size_t>(i)];
    }
    std::cout << sum << '\n';
}
\end{lstlisting}

Kodowanie redukcji przez \code{reduction} pozwala uniknąć współdzielenia
jednego licznika zapisywanego przez wszystkie wątki. Kolejność
sumowania liczb zmiennoprzecinkowych może być odmienna niż przy
obliczeniu sekwencyjnym, a~zatem wynik może różnić się przez
zaokrąglenia. Liczbę wątków można ustawić m.in. przez
\code{OMP_NUM_THREADS}; \code{OMP_SCHEDULE} wpływa na harmonogram
tylko przy \code{schedule(runtime)}. Powyższy przykład jawnie
wybiera \code{schedule(static)}.

\nota{Mieszanie pul wątków, OpenMP, równoległych algorytmów STL
oraz zadań asynchronicznych może prowadzić do \emph{oversubscription}:
liczba aktywnych wątków znacząco przewyższa dostępne zasoby CPU.
Przy nakładaniu kilku warstw równoległości trzeba świadomie
kontrolować budżet wątków i~mierzyć efekt końcowy.}

\section{Diagnostyka: kompilatory, sanitizatory, testy przeciążeniowe}
\label{sec:cpp16-testing}

\subsection{Kompilacja i~uruchamianie przykładów}

Dla przykładów wykorzystujących \code{std::jthread}, \code{atomic::wait}
i~\code{atomic<shared_ptr<T>>} wybieramy tryb C++20:

\begin{lstlisting}[language=bash,caption={Kompilacja wielowątkowego programu C++20.}]
g++ -std=c++20 -O2 -Wall -Wextra -Wpedantic -pthread file.cpp -o app
./app
\end{lstlisting}

Dla przykładów OpenMP w~GCC lub odpowiednio skonfigurowanym Clangu:
\begin{lstlisting}[language=bash,caption={Włączenie OpenMP.}]
g++ -std=c++20 -O2 -Wall -Wextra -fopenmp file.cpp -o app
./app
\end{lstlisting}

\subsection{ThreadSanitizer i~inne narzędzia}

\emph{ThreadSanitizer} (TSan) potrafi wykrywać wiele wyścigów danych
i~wybrane błędy synchronizacji w~uruchamianym programie:

\begin{lstlisting}[language=bash,caption={Budowanie i uruchamianie z ThreadSanitizer.}]
g++ -std=c++20 -O1 -g -pthread -fsanitize=thread file.cpp -o app_tsan
./app_tsan
\end{lstlisting}

Nie wszystkie środowiska poprawnie wspierają TSan; szczególnie
przy nietypowych bibliotekach wykonawczych, OpenMP i~własnych
mechanizmach synchronizacji należy weryfikować zakres instrumentacji.
AddressSanitizer i~UndefinedBehaviorSanitizer mogą dodatkowo wykrywać
część błędów czasu życia oraz naruszeń zasad języka, lecz nie są
zamiennikami TSan ani formalnego dowodu synchronizacji.

\subsection{Jak testować algorytm współbieżny?}

Dla poważniejszej struktury danych zaleca się następującą procedurę:
\begin{enumerate}
\item \textbf{Zdefiniuj kontrakt}: ilu producentów i~konsumentów,
      dozwolone operacje oraz zasady zamknięcia i~niszczenia struktury.
\item \textbf{Zidentyfikuj własność danych}: kto może zapisywać dane
      w~każdej chwili; kiedy przechodzi własność konkretnej komórki.
\item \textbf{Narysuj relacje}: wskaż atomik, który publikuje zapis
      i~odczyt, który tę publikację odbiera.
\item \textbf{Wskaż punkt linearyzacji}: jeśli API ma być linearizowalne,
      określ, kiedy operacja staje się widoczna logicznie.
\item \textbf{Przeanalizuj zakończenie}: co dzieje się przy anulowaniu,
      wyjątku, opóźnieniu konsumenta i~niszczeniu obiektu.
\item \textbf{Testuj wielokrotnie}: zmieniaj liczbę wątków, rozmiary
      danych i~częstotliwość przełączeń; uruchamiaj TSan tam,
      gdzie jest dostępny.
\item \textbf{Benchmarkuj osobno}: włącz optymalizacje i~mierz
      przy realistycznym obciążeniu, ale nie myl szybkiego testu
      z~dowodem poprawności.
\end{enumerate}

\section{Typowe błędy i~zalecenia projektowe}
\label{sec:cpp16-guidelines}

\begin{center}
\small
\begin{tabularx}{\textwidth}{@{}lX@{}}
\toprule
\textbf{Błąd} & \textbf{Dlaczego jest groźny / co wybrać} \\
\midrule
\code{volatile} jako synchronizacja & Nie eliminuje \emph{data race};
użyj mutexa lub atomików z~właściwym porządkiem. \\
\code{relaxed} do publikacji zwykłych danych & Nie wytwarza
relacji synchronizacji; zazwyczaj potrzebne \code{release/acquire}. \\
Bezrefleksyjne \code{seq_cst} & Zwykle poprawniejsze niż źle
dobrane słabsze porządki, ale nie zastąpi logicznego kontraktu obiektu. \\
Kasowanie węzła po udanym CAS & Inny wątek może nadal używać
wskaźnika; zaprojektuj odzyskiwanie pamięci. \\
Dwa atomiki jako ,,jedna transakcja'' & Brak automatycznego
atomowego snapshotu pary; użyj blokady lub niezmiennej wersji. \\
\emph{Spinlock} dla długiej pracy & Prowadzi do rywalizacji
o~CPU; zwykle użyj \code{std::mutex}. \\
Nieograniczona kolejka zadań & Może zużyć pamięć; dodaj limit
i~mechanizm \emph{backpressure}. \\
Nadmierna liczba wątków & Konkurencja o~rdzenie i~cache;
kontroluj równoległość i~profiluj. \\
Mieszanie OpenMP i~własnej puli bez planu & Może prowadzić
do \emph{oversubscription} i~spadku wydajności. \\
\bottomrule
\end{tabularx}
\end{center}

\subsection{Praktyczne reguły wyboru}

\begin{itemize}
\item \textbf{Mutex} wybieraj jako punkt wyjścia, jeśli współdzielisz
      kilka powiązanych pól i~potrzebujesz łatwo sprawdzalnych niezmienników.
\item \textbf{Atomik \code{relaxed}} jest właściwy dla prostego licznika
      niezależnego od publikacji innych danych.
\item \textbf{\code{release/acquire}} stosuj, gdy istnieje wyraźny
      punkt przekazania danych lub własności pomiędzy wątkami.
\item \textbf{Lock-free} rozważaj dopiero wtedy, gdy pomiary pokażą
      realny koszt blokad, a~projekt ma przejrzysty dowód poprawności
      oraz plan odzyskiwania pamięci.
\item \textbf{OpenMP / parallel STL} wybieraj do operacji na niezależnych
      elementach, redukcji i~równoległości danych.
\item \textbf{Thread pool} wybieraj do dużej liczby krótkich zadań,
      gdy potrzebujesz ograniczać liczbę wątków, priorytety lub
      przepustowość; osobno zdefiniuj zasady zamykania.
\end{itemize}

\section*{Podsumowanie rozdziału}
\addcontentsline{toc}{section}{Podsumowanie rozdziału}

Zaawansowana współbieżność nie polega na mechanicznym zastępowaniu
mutexów przez \code{std::atomic}. Najpierw trzeba określić
\emph{własność danych i~kontrakt algorytmu}, następnie uzasadnić
relacje \emph{happens-before} i~wyznaczyć punkty widoczności
poszczególnych operacji. Dopiero wtedy warto wybierać porządki
pamięciowe, typy kolejek, planistę i~sposób odzyskiwania pamięci.

W~szczególności brak wyścigu danych nie oznacza jeszcze poprawności
logicznej, a~nieblokujący dostęp do atomika nie gwarantuje
nieblokującego ani szybkiego zachowania całego programu. Bezpieczna
współbieżność to połączenie \emph{modelu formalnego}, \emph{kontroli
czasu życia} i~\emph{rzetelnego pomiaru}.

\section*{Literatura i dokumentacja}
\addcontentsline{toc}{section}{Literatura i dokumentacja}
\begin{itemize}
\item Anthony Williams, \emph{C++ Concurrency in Action}, wyd.~2,
      Manning, 2019.
\item ISO/IEC 14882:2020, rozdziały o~modelu pamięci, wątkach,
      atomikach i~synchronizacji.
\item \url{https://en.cppreference.com/w/cpp/atomic/memory_order}
      --- porządki pamięci i~relacje synchronizacji.
\item \url{https://en.cppreference.com/w/cpp/atomic/atomic}
      --- typy atomowe i~operacje RMW.
\item \url{https://en.cppreference.com/w/cpp/atomic/atomic_ref}
      --- atomowe odwołania do istniejących obiektów.
\item \url{https://en.cppreference.com/w/cpp/thread/jthread}
      --- \code{jthread}, czas życia i~kooperacyjne zatrzymywanie.
\item OpenMP Architecture Review Board,
      \emph{OpenMP Application Programming Interface};
      \url{https://www.openmp.org/specifications/}.
\end{itemize}
